%% =====================================================================
%%  《奶龙学导论：从奶龙到奶蛙的迷因文化研究》
%%  主控文件 main.tex
%%
%%  编译流程（必须使用 XeLaTeX）：
%%     xelatex main
%%     bibtex  main
%%     xelatex main
%%     xelatex main
%%  或直接运行： latexmk -xelatex main.tex
%%  一键脚本：   build.ps1
%% =====================================================================
\documentclass[UTF8,11pt,openany,oneside]{ctexbook}

%% =====================================================================
%%  nailong.sty —— 《奶龙学导论》教材统一体例宏包
%%  编译方式：xelatex（必须）。本文件为全书“共享宏层”，
%%  所有章节文件只能使用此处已定义的宏，不得自行 \usepackage。
%%  如确需新增宏，请先在 style/CHANGELOG.md 登记，再由主编并入本文件。
%% =====================================================================
\NeedsTeXFormat{LaTeX2e}
\ProvidesPackage{nailong}[2026/10/12 v1.0 NaiLong Studies textbook style]

%% ---------- 0. 依赖（全部经 TeX Live 2025 验证可用） ----------
\RequirePackage{iftex}
\RequirePackage{amsmath}
\RequirePackage{amssymb}
\RequirePackage{amsfonts}
\RequirePackage{amsthm}
\RequirePackage{graphicx}
\RequirePackage{xcolor}
\RequirePackage{array}
\RequirePackage{booktabs}
\RequirePackage{longtable}
\RequirePackage{multirow}
\RequirePackage{makecell}
\RequirePackage{tabularx}
\RequirePackage{enumitem}
\RequirePackage{etoolbox}
\RequirePackage[most]{tcolorbox}
\RequirePackage{tikz}
\usetikzlibrary{arrows.meta,positioning,shapes.geometric,calc,fit,backgrounds,decorations.pathmorphing,patterns}
\RequirePackage{pgfplots}
\pgfplotsset{compat=1.18}
\RequirePackage{caption}
\RequirePackage{subcaption}
\RequirePackage{float}
\RequirePackage{wrapfig}
\RequirePackage{fancyhdr}
\RequirePackage{titlesec}
%% 中文排版习惯：标题（章/节/小节）后的首段同样首行缩进
%% （西文默认是标题后首段不缩进，titlesec 沿用了这一默认，需用 indentfirst 覆盖）
\RequirePackage{indentfirst}
\RequirePackage{titletoc}
\RequirePackage{xeCJKfntef}
\RequirePackage{ulem}
\normalem
\RequirePackage{epigraph}
\RequirePackage{footmisc}
\RequirePackage{siunitx}
\RequirePackage{zhnumber}

%% ---------- 1. 配色（教材主色调：奶黄 + 墨绿 + 玄青） ----------
\definecolor{NLYellow}{HTML}{F2B705}   % 奶黄：主强调色
\definecolor{NLDeep}{HTML}{2E4A44}     % 墨绿：标题、章题
\definecolor{NLInk}{HTML}{23282B}      % 玄青：正文强调
\definecolor{NLTeal}{HTML}{1F6F6B}     % 青：定义框
\definecolor{NLAmber}{HTML}{B4761A}    % 赭：史料框
\definecolor{NLViolet}{HTML}{5B4B8A}   % 紫：方法框
\definecolor{NLGreen}{HTML}{2F6B3C}    % 绿：例题框
\definecolor{NLRed}{HTML}{9C2B2B}      % 红：警示框
\definecolor{NLGrey}{HTML}{6B7379}     % 灰
\definecolor{NLBg}{HTML}{FBF7EE}       % 浅奶底

%% ---------- 2. 版面 ----------
\setlength{\parindent}{2\ccwd}
\setlength{\headheight}{14pt}
\linespread{1.32}
\setlength{\parskip}{0.25em}
\renewcommand{\arraystretch}{1.22}
\newcommand{\NLCodeFont}{\ttfamily\small}

%% 页眉页脚：偶数页书名、奇数页章名
\pagestyle{fancy}
\fancyhf{}
\fancyhead[LE]{\small\heiti\color{NLDeep}\leftmark}
\fancyhead[RO]{\small\heiti\color{NLDeep}\rightmark}
\fancyfoot[C]{\small\color{NLGrey}\thepage}
\fancypagestyle{plain}{\fancyhf{}\fancyfoot[C]{\small\color{NLGrey}\thepage}\renewcommand{\headrulewidth}{0pt}}
\renewcommand{\headrulewidth}{0.4pt}
\renewcommand{\footrulewidth}{0pt}
\renewcommand{\headrule}{\hbox to\headwidth{\color{NLYellow}\leaders\hrule height \headrulewidth\hfill}}

%% ---------- 3. 章节标题样式 ----------
%% 章标签做成可切换宏：正文为“第…章”，附录区（main.tex 中 \appendix 之后）重定义为“附录 X”
\newcommand{\NLchapterlabel}{第\,\zhnum{chapter}\,章}
\titleformat{\chapter}[display]
  {\normalfont\heiti\color{NLDeep}}
  {\filright\fontsize{15}{18}\selectfont\NLchapterlabel}
  {0.6em}
  {\fontsize{26}{30}\selectfont\filright}
  [\vspace{0.4em}{\color{NLYellow}\titlerule[1.6pt]}]
\titlespacing*{\chapter}{0pt}{-10pt}{32pt}
%% 中文章号改用汉字数字显示（章内计数仍为阿拉伯数字）
%% 注意：不要重定义 \thechapter，否则 \arabic{chapter} 会递归；此处只改“显示名”。
\renewcommand{\thesection}{\arabic{chapter}.\arabic{section}}
\renewcommand{\thesubsection}{\thesection.\arabic{subsection}}
\renewcommand{\thesubsubsection}{\thesubsection.\arabic{subsubsection}}

\titleformat{\section}
  {\normalfont\heiti\fontsize{16}{19}\selectfont\color{NLDeep}}
  {\thesection}{0.6em}{}
\titleformat{\subsection}
  {\normalfont\heiti\fontsize{13.5}{16}\selectfont\color{NLDeep}}
  {\thesubsection}{0.6em}{}
\titleformat{\subsubsection}
  {\normalfont\heiti\fontsize{12}{15}\selectfont\color{NLInk}}
  {\thesubsubsection}{0.6em}{}
\titlespacing*{\section}{0pt}{2.2ex plus .4ex}{1.1ex}
\titlespacing*{\subsection}{0pt}{1.7ex plus .3ex}{0.8ex}
\titlespacing*{\subsubsection}{0pt}{1.3ex plus .2ex}{0.6ex}

%% 目录
\titlecontents{chapter}[0em]
  {\addvspace{0.55em}\heiti\color{NLDeep}}
  {\thecontentslabel\quad}{}
  {\hfill\contentspage}
\titlecontents{section}[1.6em]
  {\addvspace{0.1em}}
  {\thecontentslabel\quad}{}
  {\hfill\contentspage}

%% ---------- 4. 引理/定义/命题 环境（教材式“定义—命题—例”骨架） ----------
\newtheoremstyle{nlthm}%
  {1.1em}{1.1em}{\normalfont}{}%
  {\heiti\color{NLTeal}}{}{0.6em}{}
\theoremstyle{nlthm}
\newtheorem{definition}{定义}[chapter]
\newtheorem{proposition}[definition]{命题}
\newtheorem{theorem}[definition]{定理}
\newtheorem{lemma}[definition]{引理}
\newtheorem{corollary}[definition]{推论}
\newtheorem{hypothesis}[definition]{假设}
\theoremstyle{definition}
\newtheorem{example}{例}[chapter]
\newtheorem{exercise}{习题}[chapter]
\renewcommand{\proofname}{\heiti 证明}
\renewcommand{\qedsymbol}{\ensuremath{\blacksquare}}

%% ---------- 5. 教材专用彩色框 ----------
%% 5.1 术语释义：\begin{nltermbox}{术语名}...\end{nltermbox}
\newtcolorbox{nltermbox}[1][]{
  enhanced, breakable, colback=NLBg, colframe=NLTeal,
  boxrule=0.7pt, left=8pt, right=8pt, top=6pt, bottom=6pt,
  arc=2pt, fonttitle=\heiti\small,
  title={\color{white}#1}, attach boxed title to top left={xshift=6pt,yshift=-2pt},
  boxed title style={colback=NLTeal,arc=1pt,boxrule=0pt}
}

%% 5.2 史料与考据
\newtcolorbox{nlhistory}[1][史料与考据]{
  enhanced, breakable, colback=NLAmber!5, colframe=NLAmber,
  boxrule=0.7pt, arc=2pt, left=8pt, right=8pt, top=6pt, bottom=6pt,
  fonttitle=\heiti\small, coltitle=white,
  title={#1}, attach boxed title to top left={xshift=6pt,yshift=-2pt},
  boxed title style={colback=NLAmber,arc=1pt,boxrule=0pt}
}

%% 5.3 研究方法提示
\newtcolorbox{nlmethod}[1][研究方法提示]{
  enhanced, breakable, colback=NLViolet!5, colframe=NLViolet,
  boxrule=0.7pt, arc=2pt, left=8pt, right=8pt, top=6pt, bottom=6pt,
  fonttitle=\heiti\small, coltitle=white,
  title={#1}, attach boxed title to top left={xshift=6pt,yshift=-2pt},
  boxed title style={colback=NLViolet,arc=1pt,boxrule=0pt}
}

%% 5.4 批判性反思（对“乐子”保持学理距离）
\newtcolorbox{nlcritique}[1][批判性反思]{
  enhanced, breakable, colback=NLRed!4, colframe=NLRed,
  boxrule=0.7pt, arc=2pt, left=8pt, right=8pt, top=6pt, bottom=6pt,
  fonttitle=\heiti\small, coltitle=white,
  title={#1}, attach boxed title to top left={xshift=6pt,yshift=-2pt},
  boxed title style={colback=NLRed,arc=1pt,boxrule=0pt}
}

%% 5.5 案例
\newtcolorbox{nlcase}[1][案例]{
  enhanced, breakable, colback=NLGreen!4, colframe=NLGreen,
  boxrule=0.7pt, arc=2pt, left=8pt, right=8pt, top=6pt, bottom=6pt,
  fonttitle=\heiti\small, coltitle=white,
  title={#1}, attach boxed title to top left={xshift=6pt,yshift=-2pt},
  boxed title style={colback=NLGreen,arc=1pt,boxrule=0pt}
}

%% 5.6 本章要点（章首）
\newtcolorbox{nlkeypoints}[1][本章要点]{
  enhanced, breakable, colback=NLYellow!12, colframe=NLYellow!80!black,
  boxrule=0.8pt, arc=2pt, left=10pt, right=10pt, top=7pt, bottom=7pt,
  fonttitle=\heiti\small, coltitle=white,
  title={#1}, attach boxed title to top left={xshift=6pt,yshift=-2pt},
  boxed title style={colback=NLYellow!80!black,arc=1pt,boxrule=0pt}
}

%% ---------- 6. 术语、索引与代码 ----------
%% \term{术语}：加粗黑体并进入索引（使用前需 \makeindex）
\newcommand{\term}[1]{\textbf{\heiti #1}\index{#1}}
\newcommand{\tsm}[1]{\textsf{#1}}                 % 无衬线术语
\newcommand{\en}[1]{\textit{#1}}                  % 外文原文
\newcommand{\NL}[1]{\textcolor{NLDeep}{\heiti #1}}
\newcommand{\keyw}[1]{\textcolor{NLInk}{\heiti #1}}
\newcommand{\fname}[1]{\texttt{#1}}               % 文件名/代码
\newcommand{\srcnote}[1]{{\footnotesize\color{NLGrey}#1}}

%% 行内“口径提示”：用于标注资料可信度
\newcommand{\conf}[1]{{\footnotesize\color{NLAmber}\,[#1]}}

%% ---------- 7. 表格与图 ----------
\captionsetup{font={small,sf},labelfont={bf,color=NLDeep},labelsep=quad}
\captionsetup[table]{position=top,skip=4pt}
\captionsetup[figure]{position=bottom,skip=4pt}
\newcommand{\nlfigpath}{figures/}

%% 建议统一使用：
%% \begin{figure}[htbp]\centering\includegraphics[width=.7\textwidth]{...}
%%   \caption{...}\label{fig:...}\end{figure}

%% ---------- 8. 时间线宏（史料章节常用） ----------
%% \tl{1962}{巴赫金《拉伯雷》法译本问世}
\newcommand{\tl}[2]{\par\noindent\hangindent=3.6em\hangafter=1
  {\heiti\color{NLDeep}#1}\hspace{0.5em}#2\par\vspace{0.15em}}

%% ---------- 9. 习题区 ----------
\newenvironment{nlexercises}[1][本章习题]{%
  \par\medskip\noindent{\heiti\fontsize{13.5}{16}\selectfont\color{NLDeep}#1}\par
  \vspace{0.3em}
  \begin{enumerate}[label=\arabic*.,leftmargin=2.2em,itemsep=0.25em,topsep=0.2em]
}{%
  \end{enumerate}
}
\newenvironment{nlkeywords}[1][关键术语]{%
  \par\medskip\noindent{\heiti\color{NLDeep}#1：}\hspace{0.2em}
}{\par}

%% ---------- 10. 超链接（需要 hyperref 在最后加载） ----------
\endinput


\usepackage[numbers,sort&compress]{natbib}
\bibliographystyle{gbt7714-numerical}
\setcitestyle{super}   % 上标式引用，符合 GB/T 7714 习惯

\usepackage{imakeidx}
\makeindex[title=术语与人名索引,columns=2,intoc=true]

\usepackage[
  bookmarksnumbered=true,
  colorlinks=true,
  linkcolor=NLDeep,
  citecolor=NLTeal,
  urlcolor=NLAmber,
  pdfstartview=FitH,
  pdftitle={奶龙学导论：从奶龙到奶蛙的迷因文化研究},
  pdfauthor={《奶龙学导论》编写组},
  pdfsubject={迷因文化研究教材},
  pdfkeywords={奶龙,奶蛙,捧腹大笑,迷因,网络亚文化,教材}
]{hyperref}
\usepackage{bookmark}

%% ---------- 页面几何 ----------
\usepackage[a4paper,top=2.6cm,bottom=2.4cm,left=2.4cm,right=2.2cm,headsep=0.6cm,footskip=1.0cm]{geometry}

%% ---------- 中文目录标题 ----------
\renewcommand{\contentsname}{目\quad 录}
\renewcommand{\listfigurename}{插图目录}
\renewcommand{\listtablename}{表格目录}
\renewcommand{\bibname}{参考文献}
\renewcommand{\indexname}{术语与人名索引}
\renewcommand{\figurename}{图}
\renewcommand{\tablename}{表}
\renewcommand{\appendixname}{附录}
\renewcommand{\chaptermark}[1]{\markboth{\CTEXthechapter\quad #1}{}}
\renewcommand{\sectionmark}[1]{\markright{\thesection\quad #1}}

%% ---------- 全书统一“口径”“版本”宏 ----------
\newcommand{\BookTitle}{奶龙学导论}
\newcommand{\BookSubtitle}{从奶龙到奶蛙的迷因文化研究}
\newcommand{\BookEdition}{第一版}
\newcommand{\BookDate}{二〇二六年十月}
\newcommand{\BookOrg}{《奶龙学导论》编写组}

%% 封面专用样式；不改动正文共享宏层。
\definecolor{NLCoverPaper}{HTML}{FBF7EE}
\definecolor{NLCoverGreen}{HTML}{24463E}
\definecolor{NLCoverMuted}{HTML}{64756B}
\definecolor{NLCoverLine}{HTML}{D8DECE}
\definecolor{NLCoverGold}{HTML}{EAB63E}
\newCJKfontfamily\NLCoverSans{Microsoft YaHei}
\newCJKfontfamily\NLCoverBold{Microsoft YaHei}[FontFace={b}{n}{Microsoft YaHei Bold}]
\newfontfamily\NLCoverLatin{TeX Gyre Heros}

%% 仅供封面使用的固定页面坐标：原点在左上，单位厘米。
\tikzset{
  nlcovertext/.style={anchor=north west,inner sep=0pt,outer sep=0pt,align=left},
  nlcoverright/.style={anchor=north east,inner sep=0pt,outer sep=0pt,align=right}
}



\begin{document}

%% ===================== 封面 =====================
%% 固定页面坐标避免封面内容被正文页边距挤压。
\thispagestyle{empty}
\pdfbookmark[0]{封面}{frontcover}
\begin{tikzpicture}[remember picture,overlay]
  \fill[NLCoverPaper] (current page.north west) rectangle (current page.south east);
  \begin{scope}[shift={(current page.north west)},x=1cm,y=-1cm]
    \fill[NLCoverGreen] (0,24.95) rectangle (21,29.7);
    \fill[NLCoverGold] (0,0) rectangle (.38,29.7);
    \draw[NLCoverGreen,line width=.6pt] (2.1,2.66) -- (18.85,2.66);
    \node[nlcovertext,text=NLCoverGreen,font=\NLCoverSans\fontsize{11.5}{15}\selectfont]
      at (2.1,1.72) {高等学校通识教育教材};
    \node[nlcoverright,text=NLCoverMuted,font=\NLCoverSans\fontsize{9}{12}\selectfont]
      at (18.85,1.80) {迷因文化研究系列};

    \node[nlcovertext,text=NLCoverGreen,font=\NLCoverBold\bfseries\fontsize{51}{61}\selectfont]
      at (2.02,3.97) {\BookTitle};
    \node[nlcovertext,text=NLCoverGreen,font=\NLCoverSans\fontsize{17}{23}\selectfont]
      at (2.1,6.55) {\BookSubtitle};
    \node[nlcovertext,text=NLCoverMuted,font=\NLCoverLatin\fontsize{10.2}{14}\selectfont]
      at (2.14,7.87) {{\addfontfeatures{LetterSpace=9}AN INTRODUCTION TO NAILONG STUDIES}};
    \node[nlcovertext,text=NLCoverGreen,font=\NLCoverSans\fontsize{11.5}{16}\selectfont]
      at (2.1,9.21) {\BookOrg};
    \draw[NLCoverGold,line width=2.2pt] (2.1,10.52) -- (5.35,10.52);
    \node[nlcoverright,text=NLCoverGreen,font=\NLCoverSans\fontsize{10.5}{14}\selectfont]
      at (18.85,9.24) {\BookEdition};

    %% 低对比轮廓线呼应“形象—传播—社群”的研究主题。
    \begin{scope}
      \clip (.39,10.9) rectangle (21,24.95);
      \fill[NLCoverGold!9!NLCoverPaper] (13.88,18.85) circle (5.36);
      \draw[NLCoverGold!38!NLCoverPaper,line width=.5pt] (13.88,18.85) circle (5.71);
      \draw[NLCoverLine,line width=.45pt] (13.88,18.85) circle (6.05);
      \draw[NLCoverLine!65!NLCoverPaper,line width=.4pt] (13.88,18.85) circle (6.39);
    \end{scope}
    %% 使用项目已有原貌；抠图资产可由 tools/prepare_cover_asset.py 重建。
    \foreach \r/\o in {1.0/.025,.86/.03,.70/.035,.52/.04} {
      \fill[NLCoverGreen,opacity=\o] (11.9,22.97)
        ellipse ({1.55*\r} and {.13*\r});
      \fill[NLCoverGreen,opacity=\o] (15.13,23.90)
        ellipse ({1.03*\r} and {.15*\r});
    }
    \node[inner sep=0,outer sep=0] at (13.60,17.94)
      {\includegraphics[height=12.30cm]{figures/cover_naiwa_cutout.png}};

    \node[nlcovertext,text=NLCoverGreen,font=\NLCoverSans\fontsize{15}{25}\selectfont]
      at (2.1,18.66) {从一个形象\\理解一种文化};
    \node[nlcovertext,text=NLCoverMuted,font=\NLCoverSans\fontsize{9.2}{13}\selectfont]
      at (2.1,21.08) {形象\enspace /\enspace 传播\enspace /\enspace 社群};

    \node[nlcovertext,text=NLCoverPaper,font=\NLCoverSans\fontsize{13}{18}\selectfont]
      at (2.1,26.11) {数字文化与校园社群};
    \node[nlcovertext,text=NLCoverGold,font=\NLCoverLatin\fontsize{8.7}{12}\selectfont]
      at (2.13,27.37) {{\addfontfeatures{LetterSpace=11}HISTORY\enspace /\enspace MEMES\enspace /\enspace METHODS}};
    \node[nlcoverright,text=NLCoverPaper!85!NLCoverGreen,font=\NLCoverSans\fontsize{9.2}{13}\selectfont]
      at (18.85,27.32) {\BookDate};
  \end{scope}
\end{tikzpicture}
\null\clearpage


%% ===================== 版权页 =====================
\thispagestyle{empty}
\vspace*{\fill}
{\footnotesize\color{NLGrey}
\noindent\textbf{奶龙学导论：从奶龙到奶蛙的迷因文化研究}（\BookEdition）\\[0.6em]
\noindent 编写单位：《奶龙学导论》编写组\\[0.3em]
\noindent 排版工具：XeLaTeX + ctex + TikZ\\[0.3em]
\noindent 版本标识：\texttt{\BookDate}\\
\noindent 封面图像：原讲义《捧腹大笑思想与奶龙学概论》所收奶蛙形象，
相关权利归原权利人所有。\\

\vspace{1.2em}
\noindent\textbf{使用许可与免责声明}\\[0.4em]
\noindent 本书为面向高校学生社团与通识课程的非营利性开源讲义，遵循
署名—非商业性使用—相同方式共享（CC BY-NC-SA 4.0）协议发布，
欢迎各高校同好会以本校活动材料为基础进行修订与续编。\\[0.4em]
\noindent 书中涉及的“奶龙”动画形象、剧集名称及其相关商标，其知识产权
归属第七印象文化传媒（深圳）有限公司；书中涉及的“奶蛙”二次创作形象，
其署名权归属原始创作者。本书对上述形象的使用限于
\textbf{学术评论、教学示例与文化研究}目的，不主张任何权利，
亦不构成对原作品著作权的替代或商业利用。部分历史资料图
来自网络公开渠道，已尽力标注来源；如权利人认为使用不当，
请联系编写组予以替换或删除。\\[0.4em]
\noindent 书中所有理论模型、统计口径与数据均标注来源与可信度等级；
凡未标注来源者，均为编写组自行建构的“教学模型”，
\textbf{用于训练分析能力，不应被引为经验事实}。

\vspace{1.2em}
\noindent\textbf{关于本书的文风}\\[0.4em]
\noindent 本书研究对象本身带有强烈的戏谑色彩。编写组的态度是：
\textbf{对现象保持幽默，对方法保持严肃}。凡属调侃，均在正文中以
脚注或“批判性反思”框标明；凡属论证，一律给出可核查的依据。
}

\clearpage

%% ===================== 序一 / 前言 =====================
\frontmatter
\pagestyle{fancy}
%% frontmatter/preface.tex —— 前言（主编维护）
\chapter*{第一版前言}
\addcontentsline{toc}{chapter}{第一版前言}
\markboth{第一版前言}{第一版前言}

\epigraph{我一直想为奶蛙作传。}{\itshape —— 第一版编者}

本书脱胎于一份在高校社团之间流传的小册子。

2026 年 9 月，一份名为《捧腹大笑思想与奶龙学概论》的二十页讲义出现在若干高校
奶龙同好会的群文件里。它用九节的篇幅记录了奶龙与奶蛙的来龙去脉，
用一种半认真半玩笑的口气，试图解释“为什么当下的大学生如此喜欢一只捧腹大笑的青蛙”。
它的写法是随笔式的：材料丰富、感情真挚、体例自由，末尾还留了一句
“鄙人奶疏学浅”。

那份小册子提出了一些真问题。它追问：为什么越是竞争激烈的顶尖高校，
越是流行这种无厘头的东西？它给出五种解释——情绪代偿、群体稀释、反精致、
学校社交、逆反认证。它甚至自觉地在末尾打趣说“扯不下去了”。
这些解释有洞见，但它们是\textbf{平行罗列}的：五条之间是什么关系？
哪一条可以被证伪？哪一条只是把现象换了个说法？——小册子没有回答，
或者说，它本来也没有打算回答。

本书打算回答。

\section*{本书对第一版做了三件事}

\textbf{第一，把“章节”改成“章—节—目”。}
第一版的九节在体例上是并列的，本书把它们重组为十二章，
每章内部遵循“提出问题—界定概念—展开分析—给出结论—留下习题”的教材结构。
历史材料归历史章，理论模型归理论章，实践案例归实践章，
使读者能看清“什么算是证据”“什么是解释”。

\textbf{第二，把“事实”与“说法”分开。}
第一版关于奶蛙起源的叙述，与今天社群考据的结果已有出入；
书中的部分年份（例如第二季剧集的上线时间）缺乏独立来源印证；
而关于奶蛙笑声音源、关于起源的那场考据过程，
在第一版成书时还远未结束。本书对每一条材料标注了可信度等级
（A 一手权威 / B 机构公开信息 / C 社群考据 / M 本书自建模型），
凡属推测一律写明“这是推测”。这不是对第一版的苛责——
恰恰相反，正是因为第一版把材料诚实地摊在了读者面前，
后来者才有可能把它们重新整理。

\textbf{第三，把“感想”变成“可检验的假设”。}
本书保留了第一版五重机制的直觉，但把它们重写为
“压力源—补偿行为—可观察指标—竞争性解释”的四列结构，
并明确指出：在没有问卷与访谈数据的情况下，
它们只是\textbf{待检验的假设}，不是结论（见第七章、第十一章）。

\section*{阅读本书需要什么}

不需要。本书假定读者没有任何传播学或社会学的预备知识，
所有术语都在首次出现处给出定义，所有外文概念都附原文。
如果你只想看乐子，第一章和第四章会满足你；
如果你想认真做一次文化研究，第十章和第十二章是必经之路。

\section*{关于本书的态度}

本书的研究对象是一种以“绷不住”为核心体验的文化现象。
写作这样的书有一个天然的危险：要么端着架子把人逗笑的东西分析得索然无味，
要么干脆放弃分析，跟着一起笑过去。本书的立场是——
\textbf{对现象保持幽默，对方法保持严肃}。
所以你会看到：章名可以开玩笑，图表必须是严谨的；
脚注可以吐槽，结论必须有出处。

至于这本书值不值得写、奶蛙值不值得研究——
我们把它交给读者的判断，以及时间。

\vspace{2em}
\begin{flushright}
主编\quad 忘川秋裤\\
\BookDate\quad 于清华园
\end{flushright}

\clearpage

\chapter*{编写说明与使用建议}
\addcontentsline{toc}{chapter}{编写说明与使用建议}
\markboth{编写说明与使用建议}{编写说明与使用建议}

\section*{一、本书的体例标识}

\begin{nlkeypoints}
  \begin{enumerate}[label=(\arabic*),leftmargin=2.4em]
    \item 每章开头以“本章要点”列出该章要回答的三个问题；
    \item 正文中\textbf{黑体术语}表示关键概念，完整释义见附录 B（术语表），
          并可通过书末“术语与人名索引”检索；
    \item 灰底青色框为\textbf{术语释义}，赭色框为\textbf{史料与考据}，
          紫色框为\textbf{研究方法提示}，红色框为\textbf{批判性反思}，
          绿色框为\textbf{案例}；
    \item 标有“假设 X.Y”（如“假设 7.5”）的编号条目，是本书提出但尚未经经验检验的命题，
          请勿当作已证结论引用；
    \item 每章末尾附“本章小结”“本章习题”“关键术语”。
  \end{enumerate}
\end{nlkeypoints}

\section*{二、可信度标注}

本书在涉及具体事实处使用如下简写：\conf{A} 表示可追溯到一手权威来源；
\conf{B} 表示机构或平台的公开信息；\conf{C} 表示社群考据或二手评述，
\textbf{只能作为线索}；\conf{M} 表示本书为教学目的构造的模型或示例数据。
所有 \conf{C} 级材料在正文中都以“据……”“有说法认为”的方式引出。

\section*{三、教学使用建议}

\begin{nlmethod}[教学使用建议]
  \textbf{16 学时方案}（通识选修）：第一、二、五、六、七、十章各 2 学时，
  第十一章 2 学时，其余章节作为阅读材料。\\
  \textbf{32 学时方案}（社团培训）：全书讲授，
  每章配一次小组作业，第十二章的方法与伦理训练单独安排 2 学时实践。\\
  \textbf{自主阅读}：推荐路径为第一章 → 第三章 → 第四章 → 第六章 → 第七章 →
  第十章 → 第十二章。
\end{nlmethod}


%% ===================== 目录 =====================
\cleardoublepage
\pdfbookmark[0]{目录}{toc}
\tableofcontents
\clearpage
\listoffigures
\clearpage
\listoftables

%% ===================== 正文 =====================
\mainmatter

\chapter{绪论：为什么研究一只捧腹大笑的青蛙}\label{ch:01}

\begin{nlkeypoints}
  \begin{enumerate}[label=(\arabic*),leftmargin=2.4em]
    \item “奶蛙”不只是一个网络热梗，而是具备可考据的来源线索、有规模的二创生态、
          有产业接口与跨校扩散路径的\keyw{文化现象}，足以作为研究对象进入严肃讨论。
    \item 本书把“官方奶龙”“奶蛙”“捧腹大笑”区分为三个不同层面的对象，
          并在全书中坚持\keyw{证据分级}（A 一手权威 / B 机构公开 / C 社群考据 / M 本书自建模型），
          凡属推测一律显式标注。
    \item 本书回答四个递进的研究问题：它是什么符号系统？人们为什么笑？
          它如何扩散？这类现象应当如何研究？
  \end{enumerate}
\end{nlkeypoints}

2026 年，一只腹部膨大、四肢细小、黄皮绿眼的生物，以一种异常亢奋的笑声
占领了中文互联网的相当一部分信息流。它既不是任何一部动画里的角色，
也不是任何一个账号的“主人”，却拥有完整的形象谱系、稳定的“台词”、
成体系的动作、遍布高校校园的实体立牌，甚至出现了未经授权的周边商品。
有二手评述称其为“2026 年传播最广、影响最深、二创最活跃的中文互联网
meme”\conf{C}\cite{shoutanji2026naiwa}。本书把这一现象称为\term{奶蛙}，
把构成其核心表达的笑声称为\term{捧腹大笑}。

本章要做三件事：描述这个现象；说明它为何构成一个值得追问的“谜题”；
交代本书的研究问题、结构与方法承诺。

\section{一只笑了一整年的青蛙：现象描述}\label{sec:01:phenomenon}

先看材料本身。奶蛙的典型形态是：一只近似青蛙的生物，通体黄色，腹部显著膨大，
四肢短小，双手呈黑色；它几乎总是处于大笑状态，笑声经过变调处理，
异常响亮且带有某种“绷不住”的质感\conf{C}\cite{shoutanji2026naiwa}。
在各大平台的二创视频中，这只生物频繁出现在各种“抽象”桥段的结尾，
用一次集体大笑为视频收尾——这种“结尾一起笑”的固定套路本身，
后来也成了该迷因的形式特征之一\conf{C}\cite{shoutanji2026naiwa}。

与现象同步存在的，是它的“正典”背景。奶蛙形象与官方动画 IP
\term{奶龙}高度相似：奶龙由第七印象文化传媒（深圳）有限公司运营，
2017 年团队创立，2019 年构思形象，2020 年 1 月 22 日上线短视频平台
\conf{A}\cite{jiang2026nailong}\conf{B}；截至 2026 年初，其全网粉丝超 3800 万，
授权合作品牌超 200 家，衍生产品 SKU 逾 2000 个，产品销往 57 个国家和地区
\conf{A}\cite{jiang2026nailong}。也就是说，一只“不属于任何官方”的青蛙，
寄生在一个极其成功、且有专业法务团队的正版 IP 旁边生长——
这个张力将在第九章详细展开。

\begin{nlhistory}[史料与考据]
  关于奶蛙的形象来源，目前最系统的社群考据指向：2024 年 10 月 2 日，
  星野 AI 平台用户“纲手【百豪】”创建“奶龙”智能体时，平台 AI 生成的头像
  出现了体型畸变，形成最初版本奶蛙形象；2026 年 8 月前后，
  B 站用户“雾蚀心忘症”“看仙碎”等爱好者完成了对该智能体
  （ID \fname{zNCvQDKfXy}）的确认，但该智能体现已不可访问\conf{C}
  \cite{lostmedia2026naiwa}。注意：这是\keyw{考据}，不是定论——
  考据词条自陈“仅为推测、缺乏直接证据”。第四章将以整章篇幅
  示范如何把一条“未结案的证据链”当作教学材料。
\end{nlhistory}

需要强调的第一点：\textbf{奶蛙不是“一个梗”，而是三个层面的叠加}。
第一层是官方奶龙——一部定位“治愈、欢乐”的子供向动画及其完整 IP 体系
\conf{A}\cite{jiang2026nailong}；第二层是民间的奶蛙形象——一个 AI 畸变产物
经过社群筛选后稳定下来的“变种”；第三层是“捧腹大笑”——
把笑声、形象、动作与特定情境捆绑之后形成的完整表达形式。
全书将始终把这三层区分开，混用它们会产生大量错误解释
（比如把官方 IP 的成功直接归因于奶蛙，或把奶蛙的传播能力
当作官方奶龙的传播能力）。

\section{一个看似违背常理的谜题}\label{sec:01:puzzle}

第一版教材的绪论提出过一个问题：“为什么越聪明的人反而越喜欢那些
无厘头、不着边际的东西？”这个问题值得保留，因为它指出了本现象
最耐人寻味的特征：奶蛙的高频使用者与高频二创者，高度集中于高校环境，
尤其是竞争压力大的顶尖院校——同好会、立牌、校园表情包
（“数分奶蛙”“捧着航概的奶蛙”）都生长在校园空间里\conf{C}。
如果把“高智识人群 × 无厘头内容”当成一个待解释的组合，
它至少面临三种竞争读法：

\begin{enumerate}[label=\arabic*.]
  \item \textbf{纯消遣读法}：它只是一种低成本、低门槛的娱乐，
        大学生恰好有空刷手机，与“聪明”无关，是幸存者偏差。
  \item \textbf{技术—平台读法}：扩散主要由平台算法与 AIGC 工具推动，
        个体的“喜欢”是被动卷入，与心理机制无关。
  \item \textbf{功能—符号读法}：捧腹大笑承担某种系统性的情绪与
        社交功能（代偿、稀释、表态），高校恰好是这种功能
        需求最密集、传播介质最便利的场域。
\end{enumerate}

本书采取第三种读法作为主解释，但\keyw{全程保留前两种作为竞争性解释}：
第七章会给出每个机制的可检验假设与证伪条件，第十一章给出问卷设计。
“功能—符号读法”并不是本书的结论，而是本书\keyw{要检验的立场}。
在检验完成之前，它只是工作假设。这一点是本书与大多数“梗评”
文章的根本区别，也是本书的方法论承诺所在。

为什么功能读法值得认真对待？因为该现象的几个特征难以用纯消遣解释：
其一，它的使用有明显的\textbf{情境规律}——它几乎总是出现在“压力桥段”
之后，而不是随机位置；其二，它伴随稳定的\textbf{自我呈现策略}——
用一张梗图代替一句自述，把“我很难受”翻译成“我们都在笑”；
这与戈夫曼所述前台/后台管理中的“印象缓冲”结构高度同构
\cite{goffman1959presentation}；其三，它的表达门槛极低，
而在线环境的匿名与去抑制特征进一步降低了使用成本\cite{suler2004disinhibition}。
这些观察都还只是现象学层面的，机制分析见第七章。

\begin{nltermbox}{术语：meme（迷因）}
  \en{meme} 一词由道金斯在《自私的基因》第 11 章提出，
  指文化信息的复制单位，与基因类比：可复制、可变异、受选择压力
  \cite{dawkins1976selfish}。数字时代的分析以 Shifman 的三维框架
  （内容 / 形式 / 立场）为代表\cite{shifman2013memes}。
  本书在第五章给出该概念的操作化使用方式，第六章用它解剖奶蛙。
\end{nltermbox}

\section{本书的四个研究问题}\label{sec:01:questions}

全书围绕四个递进问题展开：

\begin{enumerate}[label=\textbf{Q\arabic*},leftmargin=4em,itemsep=0.4em]
  \item \textbf{符号问题}：奶蛙是一个怎样的符号系统？
        它的内容、形式与立场分别是什么？（第六章，以 Shifman
        三维框架\cite{shifman2013memes}为工具）
  \item \textbf{机制问题}：为什么大学生会笑，而且愿意反复转发？
        五种候选机制（情绪代偿、群体稀释、反精致、校园社交、
        逆反认证）各自能解释什么、不能解释什么？（第五章、第七章）
  \item \textbf{扩散问题}：它如何从短视频平台扩散到高校校园与商品世界？
        正典奶龙的产业史、奶蛙的起源考据、AIGC 生产链、
        周边经济与校园空间实践分别贡献了哪一段？（第三、四、八、九、十章）
  \item \textbf{方法问题}：研究一个“活着变”的网络文化现象，
        证据应当如何分级，考据应当如何伦理化，结论应当如何标注
        脆弱性？（第五章、第十二章）
\end{enumerate}

四个问题不是并列的：Q1 是 Q2 的输入（不懂符号系统就谈不上机制），
Q2 是 Q3 的解释层（扩散需要动力机制），Q4 贯穿全程并收束于
第十二章。这个依赖结构决定了章节顺序。

\section{本书的结构}\label{sec:01:structure}

全书十二章分三部分。第一部分（第 2—5 章）确立对象与工具：
界定“正典—变种”的对象结构，梳理正典奶龙的产业史与奶蛙的起源考据，
给出迷因理论框架与本书自建的传播教学模型。第二部分（第 6—9 章）
进入核心分析：解剖符号系统，把“为什么笑”重写为可检验的假设，
再讨论 AIGC 生产链与文化经济如何把“玩梗”接上产业。
第三部分（第 10—12 章）落到实践与方法：高校亚文化的空间实践与组织形态、
受体研究的问卷与访谈设计、以及全书方法论的收束（证据分级、考据伦理、
版权、可复现性）。

\begin{table}[htbp]\centering
  \caption{全书三部分与十二章的功能分工}\label{tab:01:structure}
  \small
  \begin{tabularx}{\textwidth}{@{}llXX@{}}
    \toprule
    部分 & 章 & 回答的问题 & 主要材料 \\ \midrule
    一：对象与工具
      & 第 2 章 & 奶龙学研究的对象与边界是什么 & 概念界定 \\
      & 第 3 章 & 正典奶龙如何成为一个产业 & 媒体公开报道（A/B 级） \\
      & 第 4 章 & 奶蛙从何而来（未结案的证据链） & 社群考据（C 级） \\
      & 第 5 章 & 用哪些理论工具分析迷因 & 理论文献（A 级） \\ \midrule
    二：符号与机制
      & 第 6 章 & 奶蛙是什么符号系统 & 形象—语言—动作语料 \\
      & 第 7 章 & 为什么笑：机制与假设 & 机制矩阵 + 假设 \\
      & 第 8 章 & AIGC 如何降低生产成本 & 生产链拆解 \\
      & 第 9 章 & 玩梗如何接入产业 & 授权与二创周边 \\ \midrule
    三：实践与方法
      & 第 10 章 & 奶蛙如何进入校园空间与组织 & 高校案例 \\
      & 第 11 章 & 大学生为什么笑：受体研究设计 & 问卷与访谈设计 \\
      & 第 12 章 & 这类研究应当怎么做 & 方法论与伦理 \\ \bottomrule
  \end{tabularx}
\end{table}

\section{方法承诺：事实、说法与模型}\label{sec:01:method}

本书涉及三类陈述，全书用统一标识区分：

\begin{enumerate}[label=\arabic*.]
  \item \textbf{事实}（\conf{A} / \conf{B}）：可追溯一手权威来源或
        机构公开信息的陈述。引用格式遵循 GB/T 7714—2015\cite{gbt7714_2015}。
  \item \textbf{说法}（\conf{C}）：社群考据、二手评述。只能以
        “据……”“有说法认为”引入，且不得作为论证前提。
  \item \textbf{模型}（\conf{M}）：本书为教学目的自建的框架、
        方程与示例数据。必须显式声明“本书构造”“非经验事实”，
        且参数取值说明为示意值。
\end{enumerate}

这一分级制度是本书最重要的体例装置，其设计理由、操作细则与局限
在第十二章系统讨论。读者若发现某条 \conf{C} 级陈述在正文中
被当作事实使用，那正是本书鼓励你指出的错误。

\begin{nlcritique}[批判性反思]
  研究“乐子”有一个天然的两难：端着架子，分析会失去对象；
  跟着一起乐，分析会失去距离。本书的立场是\keyw{对现象保持幽默，
  对方法保持严肃}——章名可以玩笑，表格必须严谨；脚注可以吐槽，
  结论必须可核查。同时必须声明：奶蛙二创中存在未经授权的形象再生产，
  这在版权法上是有争议的（第九章讨论）；本书把这一现象当作
  研究对象，不构成对任何侵权行为的背书。
\end{nlcritique}

\section{本章小结}\label{sec:01:summary}

本章完成了三件事。第一，把研究对象拆成“官方奶龙—民间奶蛙—
捧腹大笑”三个层面，这是全书一切分析的单位约定。第二，
把第一版教材“为什么聪明人爱看无厘头”的问题重写为
三种竞争读法，并明确本书的工作立场：以功能—符号读法为主解释，
以纯消遣与技术—平台读法为竞争性解释，检验工具在第七、十一章。
第三，宣布全书的证据分级制度（A/B/C/M）与四个递进研究问题。

尚未解决的问题：奶蛙的符号结构（第六章）、五种机制的
可检验化（第七章）、以及一个更深的认识论问题——
当研究对象本身以“解构严肃性”为乐时，一本严肃教材
如何保持自身立场的一致性。这个问题在第十二章收束。

\begin{nlexercises}
  \item 选取你最近 30 天内实际转发或使用过的一个梗，
        用“内容 / 形式 / 立场”三维各写一句话描述，
        并指出其中哪一维是社群硬约定（用错会被识别为外行），
        哪一维可自由发挥。
  \item 为第七章的“压力源假设”与“平台算法假设”各设计一条
        \keyw{可观察指标}（必须是单一数字或单一频率，能直接计数），
        并说明若两条指标的相关方向相反，说明什么。
  \item 在本书“事实手册”中找到两个来自不同来源、
        数值不同的同一指标（例如粉丝数），解释为什么本书
        选择并列而非取均值，并讨论这种处理对引用者的影响。
\end{nlexercises}

\nlkeywords 奶蛙、奶龙、捧腹大笑、meme（迷因）、证据分级、
参与式文化、竞争性解释

\chapter{奶龙学的对象、边界与元概念}\label{ch:02}

\begin{nlkeypoints}
  \begin{enumerate}[label=(\arabic*),leftmargin=2.4em]
    \item 本书的研究对象是一个\keyw{“正典—变种”结构}：
          官方奶龙（正典 IP）与奶蛙（民间变种）并非两个独立对象，
          而是同一文化过程的两端，奶蛙寄生于奶龙的意义系统而存在。
    \item “奶龙学”不是一个学科，而是一个\keyw{教学标签}：
          它借用学科命名形式，把对象边界、证据分级与分析工具
          打包成可复用的研究单元；本书明确给出它的
          \keyw{不是清单}（不是什么）与\keyw{是清单}（是什么）。
    \item 全书依赖四个元概念：\keyw{meme（迷因）}、
          \keyw{文本盗猎}、\keyw{抽象文化}、\keyw{融合文化}。
          四个概念分别回答“它是什么单位”“它是谁的行为”
          “它属于哪种社群风格”“它发生在什么媒介结构里”。
  \end{enumerate}
\end{nlkeypoints}

绪论交代了“研究什么”与“为什么研究”。本章回答一个更前置的问题：
\keyw{奶龙学研究的对象到底是什么}。这个问题比看上去更困难，
因为对象本身是一个正在移动的靶子：官方在迭代形象，社群在增殖变种，
平台在改写分发规则。对象边界不事先划清楚，后面的每一章
都会把三种不同层面的陈述混在一起。本章先划界，再给出
全书共用的四个元概念。

\section{两个“奶龙”：正典与民间}\label{sec:02:canon}

\begin{definition}[对象结构：正典与变种]\label{def:02:object}
  \textbf{正典奶龙}指由第七印象文化传媒（深圳）有限公司运营的
  原创动画 IP 及其全部官方内容形态（短视频、番剧、有声读物、
  表情包、衍生商品），其形象、角色与叙事由运营方控制并受版权保护
  \cite{jiang2026nailong}。
  \textbf{奶蛙}指从奶龙形象出发、经 AI 重绘发生形体畸变、
  再由社群筛选与再生产而稳定下来的民间变种形象，
  以及与之绑定的笑声、语言与动作系统
  \cite{lostmedia2026naiwa,shoutanji2026naiwa}。
  \textbf{奶龙学}指以“正典—变种”这一对结构及其传播、
  接受与产业化为对象的跨学科教学—研究单元。
\end{definition}

定义里最要紧的不是三个名词，而是那个“对”字。
奶蛙的一切辨识度都来自与正典奶龙的\keyw{相似—差异差值}：
它必须足够像奶龙（否则无法借用奶龙既有流量与情感积累），
又必须足够不像（否则不构成新刺激）。用第六节的语言说，
奶蛙是奶龙的一个\keyw{变异体}（variant），其文化能量
来自“可辨认的偏离”。这也解释了第一版教材的一个观察：
平时人们把奶蛙直接叫作“奶龙”，而奶龙同好会本质上是奶蛙同好会
——称谓的滑移本身就是两个层面相互渗透的证据\conf{C}。

\begin{nltermbox}{术语：正典（canon）}
  借自文学研究：指一个文本传统中被共同体默认为“权威版本”的
  作品集合。在网络二创语境下，“正典”往往不是法律意义上的
  版权主体，而是\keyw{意义权威}：二创者即使不受版权约束，
  仍然会自发地以官方设定为参照系（“奶龙其实没有黑色手”
  这类吐槽本身就是参照系存在的证明）。本书同时使用两个层面：
  法律层面以第七印象为权利主体\cite{jiang2026nailong}，
  意义层面以官方设定为参照系。
\end{nltermbox}

为什么要把“奶蛙”单独立名，而不是继续混用“奶龙”？
三个理由。第一，\keyw{法律}理由：奶蛙二创与正典 IP 之间存在
真实的版权张力（2026 年 3 月官方曾关评、拉黑发奶蛙账号
\conf{C}\cite{shoutanji2026naiwa}；人民日报报道亦显示官方
对“用 AI 工具批量生成低质内容”等侵权形态的维权行动
\cite{jiang2026nailong}），混用名称会模糊权利边界。
第二，\keyw{分析}理由：正典奶龙的成功是\keyw{内容产业}问题
（子供向定位、剧集更新、授权运营），奶蛙的扩散是\keyw{参与文化}
问题（畸变美学、空耳语言、群体仪式），两套解释变量几乎不重叠。
第三，\keyw{教学}理由：学生最容易犯的错误就是把“IP 火”
等同于“梗火”；分名是为了让错误在第一次出现时就被看见。

\section{学科边界：奶龙学不是什么，是什么}\label{sec:02:boundary}

“奶龙学”三个字本身带有戏谑色彩（第一版教材标题
《捧腹大笑思想与奶龙学概论》的“概论”同样是这种风格）。
本书沿用这个命名，但必须立刻声明它的边界，
否则读者会误以为这里存在一个学科。

\begin{table}[htbp]\centering
  \caption{“奶龙学”的边界：不是清单与是清单}\label{tab:02:boundary}
  \small
  \begin{tabularx}{\textwidth}{@{}p{5.2cm}Xp{5.2cm}X@{}}
    \toprule
    & \multicolumn{1}{c}{\heiti 不是什么} & & \multicolumn{1}{c}{\heiti 是什么} \\ \midrule
    对象上 & 不是奶龙动画的艺术史 & & 一个“正典—变种”传播案例 \\
    & 不是第七印象的企业史 & & 迷因文化研究的方法演练场 \\
    & 不是网络热梗盘点 & & 青年情绪与自我呈现的切片 \\ \midrule
    方法上 & 不做经验测量（本书无问卷数据） & & 提供可复用的证据分级制度 \\
    & 不做法律意见（不判断侵权与否） & & 提供可复现的分析流程 \\ \midrule
    伦理上 & 不评价具体个人 & & 以考据伦理为显式教学案例 \\
    & 不背书任何二创行为 & & 以版权张力为分析材料 \\ \bottomrule
  \end{tabularx}
\end{table}

这个边界声明里最容易被忽略的一行是“不做经验测量”。
本书包含大量“分析”，但其中没有任何一项来自本书作者
亲自采集的数据；所有数字要么是 A/B 级公开报道
（如粉丝规模、授权品牌数\cite{jiang2026nailong,chinadaily2026six}），
要么是 C 级社群考据（如起源线索），要么是 M 级教学模型。
本书的价值在于\keyw{把分析做可操作、可复现、可质疑}，
而不在于提供新的经验事实。第十一章给出的问卷
是供修订者与后续研究者施测的\keyw{设计}，不是本书的测量结果。

\section{四个元概念}\label{sec:02:meta}

全书的分析语言由四个元概念支撑。它们分别回答四个问题，
后文各章反复调用，此处只给定义与使用约定。

\subsection{meme（迷因）：最小复制单位}

道金斯把 \en{meme} 定义为文化信息的复制单位：
“就像生物基因在生物体间复制一样，meme 在头脑间复制，
经历变异、选择与保留”\cite{dawkins1976selfish}。
《延伸的表现型》进一步把“复制子”的视角延伸到
文化结构的层面\cite{dawkins1982extended}。
但道金斯式的迷因概念在数字时代面临一个操作难题：
\keyw{复制的边界在哪里}？一条奶蛙视频被转发一万次，
每次转发都算“同一迷因”吗？本书采用 Shifman 的处理方式：
不追问“同一性”的本体论，而是用三维框架
（内容 / 形式 / 立场）给每个用例做\keyw{结构描述}，
再在“结构相似度”高于社群默许阈值时归入同一迷因家族
\cite{shifman2013memes}。第六章将给出奶蛙家族的
归并判据。

\subsection{文本盗猎：行为主体与行为结构}

“文本盗猎者”（textual poachers）指主动挪用、拼贴、
改写既有文本的受众。Jenkins 在《文本盗猎者》一书中
把粉丝写作、混剪与同人活动组织为“参与式文化”的核心形态
\cite{jenkins1992poachers}。奶蛙二创者正是典型盗猎者：
他们盗用正典奶龙的形象资产，嫁接直播连麦的录音资产，
再嫁接 AIGC 工具的生产能力。
本书使用这一概念时的纪律是：\keyw{描述结构，不评价合法性}。
“盗猎”是分析术语，不等于法律上的“侵权”判定，
也不等于价值上的“创意致敬”褒奖；法律层面的张力
在第九章处理。

\subsection{抽象文化：一种社群风格}

“抽象”是中文互联网的原生概念，没有直接对应的外文术语。
本书给出操作性定义：

\begin{definition}[抽象（抽象文化）]\label{def:02:abstract}
  一种以“意义悬置”为核心策略的社群风格：内容不追求
  明确的语义传达，而以“看不懂—被冒犯—仍然好笑”的
  体验回路为组织原则；典型标记包括无逻辑谐音（空耳）、
  故意劣化的制作质量、对严肃语体的戏仿
  （如把宗教文体套用于立牌纪念文本）\conf{C}。
\end{definition}

这个定义是本书构造的教学定义（M 级），它来自对第一版教材
所引材料的归纳（“糖”“倒放大法”“鲱鱼罐头”等社群用语
的使用方式），不是学术共识。定义中“意义悬置”一项最要紧：
它解释了为什么“看不懂”不会杀死传播——
在抽象文化里，\keyw{看不懂本身就是内容}，
是社群内部的准入测试与归属标记
（“看懂了说明你是自己人”）。第十章会看到，
这种“准入测试”功能在高校同好会的纳新中反复出现。

\subsection{融合文化：媒介结构的位置}

Jenkins 用“融合文化”描述内容在多平台间流动、
受众从消费者转为参与者的媒介生态\cite{jenkins2006convergence}；
《可传播的媒体》进一步提出：在融合环境下，
内容的价值不取决于制作方的“黏性”投放，
而取决于受众的“可传播性”行为\cite{jenkins2013spreadable}。
奶蛙恰好卡在这个结构上：正典奶龙由制作方控制分发
（剧集平台、官方账号），奶蛙的“可传播性”则完全由
二创者的转发、二创与线下搬运驱动。
第八章的生产链分析与第九章的产业张力分析，
都以这个“双轨分发”结构为底图。

\section{从梗到“学”：命名的正当性}\label{sec:02:legitimacy}

最后一节处理一个元问题：一本教材叫“××学概论”，
这种命名本身是不是也是一种“梗”？

本书的回答是：\keyw{是，而且这正是它的功能之一}。
“奶龙学概论”的命名模仿了正规学科教材的题录形式，
但它不是伪装——前言已明确本书为开源讲义，
证据分级制度本身就是反伪装的装置。
命名的作用恰恰相反：通过把戏谑对象装进
“最严肃的容器”（教材体例、定义—命题—习题、
GB/T 7714 引文\cite{gbt7714_2015}），它迫使读者
把注意力从“内容好不好笑”转移到“分析是否站得住”。
这是一种\keyw{语域错配}（register mismatch）的自觉使用：
严肃体例 × 戏谑对象，错配本身成为方法立场的表达。

必须同时承认这一策略的限度：教材体例的严肃性
不能替代证据的严肃性。如果本书某处
用 C 级考据支撑了 A 级口吻的结论，
那么错配就从“立场表达”滑向了“误导”——
第十二章的可复现性清单正是为拦截这种滑动而设。

\section{本章小结}\label{sec:02:summary}

本章把研究对象确定为“正典奶龙—民间奶蛙”这对结构，
划定了“奶龙学”的边界（是清单与不是清单各四项），
并给出全书四个元概念：meme（最小复制单位，
操作化为结构相似度）、文本盗猎（行为结构，
不评价合法性）、抽象文化（意义悬置的社群风格，
M 级教学定义）、融合文化（双轨分发结构）。
四个概念分别对应四个问题：单位、行为者、风格、媒介。
第三章起，对象从定义走向历史：
先看正典奶龙如何成为产业（第三章），
再看民间奶蛙的来源考据（第四章）。

\begin{nlexercises}
  \item 给“正典”的\keyw{法律层面}与\keyw{意义层面}
        各举一个在奶蛙语料中的具体实例
        （前者涉及权利主体行为，后者涉及二创者的自发参照），
        并说明两者是否总是一致。
  \item 本章四个元概念中，哪一个最容易被误用为
        “贴标签”（只说“这是文本盗猎”而不展开结构）？
        写出你使用它时必须回答的三个结构问题。
  \item 用定义 2.2 判断下列三个用例是否属于“抽象”：
        (a) 一段把课程作业封面做成“讣告”体的视频；
        (b) 一张画质劣化但语义完整的吐槽表情包；
        (c) 一段完全语义清晰、仅画质劣化的视频。
        逐条写出你的判定依据（意义悬置 / 空耳 / 语体戏仿三项中
        各命中哪几项）。
\end{nlexercises}

\nlkeywords 正典与变种、文本盗猎、抽象文化、融合文化、
meme（迷因）、语域错配

\chapter{正典奶龙：一个原创 IP 的产业史（2017—2026）}\label{ch:03}

\begin{nlkeypoints}
  \begin{enumerate}[label=(\arabic*),leftmargin=2.4em]
    \item 正典奶龙的成长是一条完整的\keyw{内容产业}轨迹：
          2017 年公司成立、2019 年形象构思、2020 年 1 月 22 日上线、
          剧集化与大屏化（2023—2026）、六周年仪式（2026 年 1 月）、
          以及授权体系的规模化（200+ 品牌、2000+ SKU、57 国）。
    \item 对同一指标，本书坚持\keyw{来源口径并列}：
          “全网粉丝超 3800 万”（A 级）与“全球粉丝突破 3900 万+”（B 级）
          并存呈现，不取均值、不互相覆盖。
    \item 正典奶龙是\keyw{子供向}作品，但其受众外溢到高校人群
          是一个真实且可分析的现象：本章提出“内容内核—平台结构—
          民间再生产”三层溢出解释，并在第九章与奶蛙产业线相对照。
  \end{enumerate}
\end{nlkeypoints}

第二章给对象结构划了界：奶龙学的核心是“正典—变种”这一对。
本章先处理这一对中稳定的一端——正典奶龙。
它的重要性有两个：其一，奶蛙的一切辨识度都以它为参照系，
不了解正典，就无法度量变种的“偏离”；
其二，正典奶龙是一个有专业法务团队、有明确授权体系的产业主体，
它与奶蛙之间的张力（关评、拉黑、维权）构成本书
第九、第十二章的关键材料。本章只回答一个问题：
\keyw{一只子供向的龙，如何成为一家有产业规模的 IP 公司}。

\section{起点：一家深圳公司}\label{sec:03:origin}

据《人民日报》2026 年 1 月 15 日报道，
2017 年，谢承恩与团队扎根深圳，
成立第七印象文化传媒（深圳）有限公司
\conf{A}\cite{jiang2026nailong}。
公司在 2019 年开始构思奶龙形象，其对外表述的创作定位是：
“想传递治愈、欢乐的积极情感，帮人们缓解焦虑、紧张等情绪”
（谢承恩语）\conf{A}\cite{jiang2026nailong}。
这一表述值得保留：它说明正典奶龙的\keyw{情感定位}
从构思期就锚定在“减压”而非“剧情深度”上——
后文将看到，这与奶蛙的“捧腹大笑”在功能上形成
一种耐人寻味的呼应（第六章、第七章详论）。

形象设计本身经历过迭代：团队“经过多次尝试，
最终设计出拥有黄色皮肤、呆萌可爱的奶龙形象”
\conf{A}\cite{jiang2026nailong}。
注意“黄色皮肤”这一细节：奶蛙保留了黄色，
而放弃了呆萌的圆脸比例——变种对正典的继承与偏离，
在配色层面就是“保留”，在比例层面就是“畸变”
（第四章考据、第六章符号分析的基础材料）。

\section{上线：短视频时代的冷启动}\label{sec:03:launch}

2020 年初，第一条奶龙短视频上线并迅速走红
\conf{A}\cite{jiang2026nailong}。
本书把其上线日定为 \textbf{2020 年 1 月 22 日}：
该日期来自六周年节点的倒推——
中国日报网 2026 年 1 月 28 日报道明确写道
“国产原创动漫 IP 奶龙，在 2026 年 1 月 22 日迎来了他的六周年庆”
\conf{B}\cite{chinadaily2026six}，
2026 减 6 年，恰为 2020-01-22。
本书将其标注为 \conf{B} 级（机构公开信息的日期倒推），
而非 \conf{A} 级，因为报道原文没有直接写出
“2020-01-22 上线”这一句。

上线后的早期扩张数据，目前只有第一版教材给出的口径：
“同年抖音粉丝破千万，单条视频 3 天点播破亿”
\conf{C}；2021 年，《奶龙》短视频合集在优酷、爱奇艺、腾讯
同步更新，有声读物专辑上线，“8 个月播放量破亿”\conf{C}。
这些数字无法在本书引用的 A/B 级来源中独立印证，
故全部保留 \conf{C} 标注，供修订者复核。

\begin{nlhistory}[史料与考据]
  2020 年短视频阶段的奶龙，与 2026 年的奶龙
  在商业形态上是两家公司：前者靠内容自然流量，
  后者靠“内容 + 授权”双轮驱动。谢承恩在人民日报报道中
  自述了这次转型：“形象要优化，内容要升级，
  商业化也还没起步”（2020 年初的瓶颈）；
  之后“团队从广告合作起步，逐步拓展到商业授权，
  积极探索‘内容 + 授权’双轮驱动的模式”\conf{A}
  \cite{jiang2026nailong}。
\end{nlhistory}

\section{剧集化：从短视频到番剧}\label{sec:03:series}

正典奶龙的剧集序列，各平台条目的编号方式并不完全一致
（有的平台把 2026 年的新番编作“第二季”的续作，
有的按系列名单列），本书按“系列名 + 上线时间”记录，
不强行统一季数编号：

\begin{table}[htbp]\centering
  \caption{正典奶龙剧集序列（以官方上线宣传与媒体公开信息为准）}
  \label{tab:03:series}
  \small
  \begin{tabularx}{\textwidth}{@{}llXc@{}}
    \toprule
    系列名 & 上线时间 & 备注 & 可信度 \\ \midrule
    《奶龙搞怪大作战》（第一季） & 2023 年
      & 进入高中、大学生视野的节点（第一版教材口径） & \conf{B/C} \\
    《奶龙与小七之大战暴暴龙》（第二季） & 2024 年
      & 在高中生与大学生之间“爆火”（第一版教材口径，未获独立权威源印证） & \conf{C} \\
    《奶龙之奇妙冒险》（第三季） & 2025 年 1 月
      & 一个月内全网播放量破 1 亿 & \conf{B} \\
    《奶龙搞怪大作战 2》 & 2026 年 1 月 28 日
      & 腾讯、优酷、爱奇艺等主流平台；覆盖风行、聚好看、
        银河少儿等大屏终端；嘉佳卡通、金鹰卡通等电视台陆续播出 & \conf{B} \\ \bottomrule
  \end{tabularx}
\end{table}

表中可见一条清晰的\keyw{渠道下沉}轨迹：
短视频平台（2020）→ 长视频平台同步（2021）→
番剧化（2023—）→ 大屏终端 + 电视台（2026）。
对一只子供向 IP 而言，大屏与电视台意味着
它进入了“家庭共同观看”的场域；而 2026 年 1 月 24 日
深圳海岸影城的六周年包场观影（约 200 位粉丝到场，
影片结束后“全场更自发齐唱生日歌”）
则说明其粉丝组织已具备线下动员能力
\conf{B}\cite{chinadaily2026six}。

\section{六周年：作为仪式的产业节点}\label{sec:03:anniversary}

2026 年 1 月 19 日至 25 日，第七印象举办“奶龙六周年”生日会，
含三大板块：线下观影会（新番抢先看）、五城地标大屏共庆
（武汉、青岛、天津、成都、苏州）、三城线下门店见面会
（武汉、乌鲁木齐、哈尔滨，联动茶饮品牌“吾饮良品”）
\conf{B}\cite{chinadaily2026six}。
1 月 22 日生日当天，奶龙以“歌手”身份发布首支单曲
《奶龙的小确幸》，生日直播掉落六周年限定周边与合作品牌礼包
\conf{B}\cite{chinadaily2026six}。
1 月 25 日，奶龙亮相中央广播电视总台（CMG）《动漫盛典》，
“以打头阵的高姿态惊艳亮相”（报道语）\conf{B}\cite{chinadaily2026six}。

\begin{figure}[htbp]\centering
\begin{tikzpicture}[font=\small]
  \draw[NLDeep,line width=1.2pt,-{Stealth[length=3mm]}] (-0.3,0) -- (14.1,0);
  \foreach \x in {0,1.7,3.4,6.1,7.8,9.5,11.6,13.2}{
    \fill[NLYellow] (\x,0) circle (2.2pt);
    \draw[NLDeep] (\x,0) -- (\x,0.12);
  }
  \node[font=\footnotesize,align=center,inner sep=1pt] at (0,0.95) {2017\\第七印象成立};
  \node[font=\footnotesize,align=center,inner sep=1pt] at (1.7,-0.95) {2019\\构思奶龙形象};
  \node[font=\footnotesize,align=center,inner sep=1pt] at (3.4,0.95) {2020-01-22\\短视频平台上线};
  \node[font=\footnotesize,align=center,inner sep=1pt] at (6.1,-0.95) {2023\\第一季《搞怪大作战》};
  \node[font=\footnotesize,align=center,inner sep=1pt] at (7.8,0.95) {2024\\第二季《大战暴暴龙》};
  \node[font=\footnotesize,align=center,inner sep=1pt] at (9.5,-0.95) {2025-01\\第三季《奇妙冒险》};
  \node[font=\footnotesize,align=left,inner sep=1pt] at (11.6,0.95) {2026-01-19/25\\六周年生日会\\（五城大屏 + 三城门店）};
  \node[font=\footnotesize,align=left,inner sep=1pt] at (13.2,-0.95) {2026-01-28\\《搞怪大作战 2》多平台上线};
\end{tikzpicture}
  \caption{正典奶龙产业里程碑（2017—2026）。2024 年节点可信度为
          \conf{C}（第一版教材口径），其余节点见事实手册。}
  \label{fig:03:timeline}
\end{figure}

这一周的意义，产业分析上叫“节点营销”，
文化分析上叫\keyw{年度仪式}：六周年庆典把“粉丝陪伴叙事”
（小红书“留言小纸条”活动，“数万条留言中，陪伴与成长
成为粉丝们提到的共同关键词”\conf{B}\cite{chinadaily2026six}）
转化为可见的消费与传播事件。对本书而言，它的价值在于
提供了一个\keyw{时间锚点}：奶蛙现象在 2026 年上半年全面爆发之前，
正典奶龙的官方动作集中在 1 月的这一周——
两个时间线的相对位置（官方仪式在前，民间爆火在后），
是第九章讨论“官方与二创关系”的基本史实依据。

\section{产业量化口径：为什么并列而非合并}\label{sec:03:numbers}

截至 2026 年初，两个最主要的来源给出了
数值不同但口径相近的量化描述。本书的处理方式如下表：

\begin{table}[htbp]\centering
  \caption{正典奶龙产业指标的两个来源口径（并列呈现，不作合并）}
  \label{tab:03:numbers}
  \small
  \begin{tabularx}{\textwidth}{@{}p{3.0cm}p{4.6cm}p{4.6cm}X@{}}
    \toprule
    指标 & 《人民日报》2026-01-15（\conf{A}）
      & 中国日报网 2026-01-28（\conf{B}） & 本书处理 \\ \midrule
    粉丝规模 & 全网粉丝超 3800 万 & 全球粉丝规模突破 3900 万+
      & 并列；“全网”与“全球”口径是否同构不明，不取均值 \\
    流量 & 流量破千亿 & 未提及 & 仅 A 级来源给出，保留 \\
    授权合作品牌 & 超 200 家 & 超 200 家（截至 2025 年底）
      & 两源一致，可作 A/B 级陈述 \\
    衍生产品 & 数超 2000 个 & SKU 逾 2000 个
      & 两源一致 \\
    销售地域 & 57 个国家和地区 & 未提及
      & 仅 A 级来源给出 \\
    微信表情包 & 未提及 & 累计发送量突破 80 亿次（截至 2025 年底）
      & 仅 B 级来源给出 \\
    译制 & 未提及 & 英语、韩语、印尼语等 15 种语言
      & 仅 B 级来源给出 \\
    政府支持 & 2024 年参评深圳市新型业态企业，获 25 万元扶持奖励
      & 未提及 & 仅 A 级来源给出 \\ \bottomrule
  \end{tabularx}
\end{table}

为什么本书不取 3800 万与 3900 万的均值？
三个理由。第一，\keyw{口径不明}：“全网”（国内平台合计？）
与“全球”（含海外平台？）的统计范围不同，
数值差异可能主要来自范围而非时间差。
第二，\keyw{时间差}：两报道相隔 13 天，
粉丝数的自然增长本身就能解释差距。
第三，\keyw{教学目的}：本书要训练的正是
“面对多个来源时不悄悄合并”的引用纪律——
取均值是一个看似中立、实则制造了不存在之精度的动作。
读者在第十二章的可复现性清单里会看到这一纪律的操作化版本。

\section{子供向作品为何溢出到大学生}\label{sec:03:overflow}

第一章提出的谜题是“为什么聪明人爱看无厘头”。
在正典一侧，这个谜题有一个更精确的版本：
奶龙的目标受众是\keyw{低龄观众}（“嘉佳卡通、金鹰卡通”等
播出终端\conf{B}\cite{chinadaily2026six}，
“连接亲子温情与成人解压的最大公约数”是报道原话
\conf{B}\cite{chinadaily2026six}），
但第一版教材的观察与高校同好会的存在都表明，
其实际受众显著外溢到了大学生群体
\conf{C}。本书给出三层解释，
均标注为本书分析（M 级），供检验而非供引用：

\begin{enumerate}[label=\arabic*.]
  \item \textbf{内容内核层}：奶龙的内容组织是“情绪单元”
        而非“叙事弧”——单个短视频/单集即可闭合，
        不需要追更成本。这与大学生碎片化注意力
        的接收条件匹配；而“治愈、欢乐”的定位
        （\conf{A}\cite{jiang2026nailong}）恰好对准了
        第一章所述的压力情境。
  \item \textbf{平台结构层}：短视频平台按兴趣分发，
        不按年龄分发。一个 20 岁用户看到奶龙，
        与一个 8 岁用户看到奶龙，在算法层面没有区别；
        “子供向”是内容标签，不是分发过滤器。
  \item \textbf{民间再生产层}：真正完成“大学生化”的
        不是官方内容，而是奶蛙式的二创——
        它把正典形象翻译成了大学生社群的
        内部语（空耳、抽象、立牌）。
        这一层在第四章与第九章展开。
\end{enumerate}

三层的分工是：第一层解释“为什么不排斥”，
第二层解释“为什么会看到”，第三层解释“为什么会主动传播”。
缺任何一层，都只能部分解释受众外溢。

\begin{nlcase}[案例 3-1：大屏与“全家桶”]
  2026 年 1 月，奶龙在武汉、天津、成都的极地海洋公园剧场大屏、
  青岛“极地之光·海上摩天轮”灯光秀、苏州阳澄半岛乐园
  “我是大明星剧场”同步播放六周年影像\conf{B}
  \cite{chinadaily2026six}。
  这些场景的共同特征是\keyw{家庭共同在场}
  （带孩子逛乐园的成年人）。
  一只子供向 IP 选择“五城地标大屏”而非“五城地铁广告”，
  其传播意图本身就承认了
  “观看者是亲子组合”这一结构——
  这为“受众外溢先于官方意图”的判断提供了
  一个反向参照：官方把资源投给了它认为的受众。
\end{nlcase}

\section{版权治理：官方如何对待“奶蛙式”二创}\label{sec:03:copyright}

产业线不能与法律线分开。据《人民日报》报道，
奶龙 IP 大热后，“侵权盗版问题”成为常态，
谢承恩列举的形态包括：“有商场未经授权搭建相关游乐设施，
IP 形象严重失真，安全性也无法保证；有人截取视频片段，
甚至用 AI 工具批量生成低质内容”\conf{A}\cite{jiang2026nailong}。
企业的应对是“组建专业法务团队”，
“在深圳市知识产权保护中心、市场监管部门等指导下，
依托地方快速维权机制与专项执法行动，及时维护了 IP 形象”
\conf{A}\cite{jiang2026nailong}；
海外则因“维权链条长、成本高”，
由广东省版权部门“协助对接海外知识产权组织与法律资源”
\conf{A}\cite{jiang2026nailong}。

与“奶蛙”直接相关的官方行为，目前只有 \conf{C} 级材料：
2026 年 3 月，“奶龙官方呼吁停止丑化二创，
关闭视频评论区并批量拉黑发布奶蛙的账号”
\conf{C}\cite{shoutanji2026naiwa}；
另一二手评述亦提及“奶龙官方将所有在奶龙视频评论区
中发奶蛙的人一律拉黑”，并称“直至今日你依然可以
在奶龙官方视频的评论区看见无数奶蛙图片”
\conf{C}\cite{shoutanji2026naiwa}。

这两条材料放在一起，读出一个对本书至关重要的结构：
\keyw{官方的抑制是不完全的，而不完全的抑制
恰好构成了迷因生态的缝隙}。第九章将论证，
如果 2026 年 3 月的关评与拉黑是彻底的，
奶蛙现象在时间上会被截断；
而它实际上没有——这一“失败的控制”
本身成为后续传播的燃料（第七章“逆反认证”机制）。
必须同时指出另一面：本章引用的全部 A 级材料
都说明官方维权是有组织、有法律依据的，
本书不对任何具体的二创行为作合法性判断（第十二章）。

\section{本章小结}\label{sec:03:summary}

本章回答了“正典奶龙如何成为产业”：
2017 年公司化、2019 年形象构思（定位“治愈、欢乐”）、
2020-01-22 上线（B 级日期倒推）、2021 年多平台 + 有声读物、
2023—2026 剧集化与大屏化、2026 年 1 月六周年仪式周、
以及“内容 + 授权”双轮驱动的规模化
（200+ 品牌、2000+ SKU、57 国、80 亿次表情包、15 种语言）。
在方法上，本章演示了两件事：
来源口径的并列呈现（表 \ref{tab:03:numbers}），
以及把“受众外溢”拆成内容内核、平台结构、
民间再生产三层可检验的解释。
遗留问题：正典一侧到此已经闭环，
但“变种”一侧的材料远没有这么齐——
奶蛙的起源至今是一条未结案的证据链。
这正是下一章的主题。

\begin{nlexercises}
  \item 表 \ref{tab:03:series} 中第二季（2024）标注为 \conf{C}
        而第三季（2025）标注为 \conf{B}：
        请用事实手册说明这一差异的依据，
        并解释为什么“同一系列、不同年份”可以
        拥有不同的可信度。
  \item 为“受众外溢三层解释”（\ref{sec:03:overflow}）
        各设计一条\keyw{可观察、可计数}的指标
        （例如：某平台 20 岁以上用户占比、
        二创作者中高校生占比、大屏场景的亲子观众比例），
        并指出哪一层最难测量、为什么。
  \item 官方“关评 + 拉黑”（\conf{C} 级）与
        “法务团队 + 快速维权机制”（\conf{A} 级）
        是否矛盾？请用第九章的框架
        （官方控制 vs 迷因生态缝隙）作答，
        并说明你的答案里哪些依赖 C 级材料、
        哪些只依赖 A 级材料。
\end{nlexercises}

\nlkeywords 第七印象、正典奶龙、剧集序列、六周年、
产业量化口径、受众外溢、版权治理、关评拉黑

\chapter{奶蛙前史：一场正在发生的失传媒体考据}\label{ch:04}

\begin{nlkeypoints}
  \begin{enumerate}[label=(\arabic*),leftmargin=2.4em]
    \item 奶蛙的起源目前是一条\keyw{未结案的证据链}：
          最系统的社群考据指向 2024 年 10 月 2 日星野 AI 平台
          “纲手【百豪】”所创建的“奶龙”智能体，
          但该智能体现已不可访问，考据材料自身
          声明“仅为推测、缺乏直接证据”。
    \item 起源考据经历三个可辨识阶段：
          \keyw{误指}（“季洪行说”，2024-03 起）、
          \keyw{线索}（星野 AI 水印与截图，2024-12 前后）、
          \keyw{发现与证伪并存}（2026-08 智能体确认 +
          复现失败）。三阶段构成“错误归因→证据积累→
          归因修正”的标准考据动力学。
    \item 本章同时示范两条纪律：
          其一，\keyw{不可复核性}必须被显式写入结论
          （“指向某来源”与“证明某来源”是两回事）；
          其二，\keyw{考据伦理}——本书不推断、
          不转述任何关于具体个人的身份、健康、
          生死与私德的内容。
  \end{enumerate}
\end{nlkeypoints}

第三章结束时留下了一个问题：
正典一侧的材料有 A 级来源支撑，闭环良好；
变种一侧，最好的材料是一份自陈“待改进、来源缺失”
的维基词条。本章不回避这个落差，反而把它当作
全书的方法论示范对象：
\keyw{当证据链未结案时，研究者应当如何诚实地写作}。
具体而言，本章做四件事：
把奶蛙起源考据拆成三个可辨识阶段；
讨论“不可复核性”这一认识论困境；
单独处理笑声音源这条第二证据链；
以及——这可能是本章最重要的一节——
说明为什么关于“创作者”的写作
必须有一条不能逾越的伦理线。

\section{什么是“考据”：方法与标准}\label{sec:04:method}

在进入材料之前，先给“考据”一个操作性定义，
因为它在本书中承担特殊的功能：

\begin{definition}[网络考据（本书用法）]\label{def:04:kaoju}
  网络考据指：由非职业研究者组成的社群，
  以公开可检索的网络材料（截图、帖子、智能体档案、
  时间戳、水印）为证据，对某一网络文化对象的
  来源、作者、演化路径进行追溯与归因的活动。
  其结论形态通常是\keyw{带置信度的指向}
  （“最可能的来源是 X，证据为……”），
  而非\keyw{证明}（“X 是来源，证毕”）。
\end{definition}

这个定义里最要紧的是结论形态的差异。
网络考据与法庭证据、档案学考据共享“证据—推理—结论”
的骨架，但它缺少三样东西：
(1) 完整的证据保全（原始载体随时消失）；
(2) 独立的质证环节（同一社群既是调查者又是受众）；
(3) 明确的“结案”标准（指向某来源后，
证据链是否闭合由谁判定？通常没有判定者）。
因此，本书引用考据材料时，一律使用
“考据指向……”“有考据认为……”的措辞，
并保留考据材料自身的限定语。
这是把定义~\ref{def:04:kaoju} 落进句式的直接结果。

\begin{nlmethod}[研究方法提示]
  本章的材料来源与可信度：
  核心材料来自失传媒体中文维基
  （\conf{C}，词条自陈“待改进”“来源缺失”）
  与一篇 2026 年 9 月的二手评述文章
  （\conf{C}）；第一版教材的叙述
  （\conf{C}）作为早期考据的快照保留。
  全书没有一条 \conf{A} 或 \conf{B} 级来源
  直接描述奶蛙起源——
  这一“缺口”本身就是本章要分析的现象之一
  （官方与主流媒体对民间二创的沉默）。
\end{nlmethod}

\section{第一阶段：误指（2024-03 起）}\label{sec:04:mistaken}

按失传媒体中文维基词条的整理，
最早可观测的“奶蛙”相关传播出现在 2024 年 3 月 7 日：
有网友将“奶蛙”形象与抖音博主“季洪行”的视频进行
AI 融合\conf{C}\cite{lostmedia2026naiwa}。
由于该融合内容在一段时间内传播最广，
“奶蛙”长期被认为是
“季洪行视频 + 奶龙 AI 融合”的产物
\conf{C}\cite{lostmedia2026naiwa}。

这个错误归因有三个值得注意的结构性原因：

\begin{enumerate}[label=\arabic*.]
  \item \textbf{可见性偏误}：
        被误指的对象是“当时最出圈”的内容。
        考据的对象如果是“第一个存在的”，
        调查者实际找到的往往只是
        “第一个被广泛看见的”——
        传播链的可见端点与起源端点并不重合。
  \item \textbf{时间标记的误导性}：
        融合内容出现时，“我是奶龙”梗
        （正典奶龙的“糖”内容，第一版教材有详细讨论）
        正在高峰期，奶龙形象的
        AI 加工需求本就旺盛，
        使得“融合说”在直觉上可信。
  \item \textbf{小范围传播的“真空期”}：
        词条的判断是：当时最出圈的是
        “我是奶龙”此梗而不是奶蛙，
        因此奶蛙“在小范围传播”，
        直到 2026 年“奶蛙捧腹大笑”
        视频才爆火\conf{C}\cite{lostmedia2026naiwa}。
        这一判断（若成立）意味着
        奶蛙在 2024—2025 年间
        存在一个“低可见度的潜伏期”——
        这是第五章传播模型里
        “二次爆发”假设（\ref{sec:05:sir}）
        的第一个经验支撑。
\end{enumerate}

\begin{nlcritique}[批判性反思]
  误指阶段还留下一个次生问题：
  当“错误答案”先于正确答案传播了
  一年半以上，它本身会变成一种
  \keyw{文化事实}——大量早期材料、
  表情包与口头叙述都以“季洪行融合说”为前提。
  考据修正“来源”时，
  无法回滚已经基于错误来源
  生产出来的二创。
  这在第七章的“形式稳定、内容可换”
  分析中会被反复用到。
\end{nlcritique}

\section{第二阶段：星野 AI 线索（2024-12 前后）}\label{sec:04:clue}

词条记录的第二阶段线索来自
抖音用户对账号“欢迎大家恐龙世”
（现名“欢迎大家恐龙世界”）的声讨：
该账号“经常发布儿童邪典内容”，
因用户声讨后被官方封禁、转私密账号
\conf{C}\cite{lostmedia2026naiwa}。
在声讨材料的截图中，
有人发现了“奶蛙”的照片
（声讨发布时间在 2024 年 12 月前后）
\conf{C}\cite{lostmedia2026naiwa}；
进一步地，该照片颜色反转后
可见 AI 生成水印，
结合配文“星野ai怎么了”，
用户推测其来源为星野 AI 生成
\conf{C}\cite{lostmedia2026naiwa}。

从证据学上看，这一段的价值在于
它把候选来源从“不确定的某个博主”
收窄到“某个具体平台的某个生成环节”——
\keyw{候选空间缩小}是考据推进的主要形态，
而不是“一步找到真凶”。
同时，这条线索也带有典型的
“孤证”特征：一张截图、一段配文，
均出自被声讨的账号生态，
且原始生成过程无法重放。
本书对此的处理是：
保留其作为“线索”的地位，
不赋予其“证据”的地位。

\section{第三阶段：发现与证伪并存（2026-08）}\label{sec:04:discovery}

2026 年 8 月 4 日，
B 站用户“雾蚀心忘症”发布评论，
表示在星野 App 中找到了
符合形象的奶蛙智能体
\conf{C}\cite{lostmedia2026naiwa}。
词条随后将智能体定位到
星野用户“纲手【百豪】”于
2024 年 10 月 2 日创建的
“奶龙”智能体，ID 为 \fname{zNCvQDKfXy}：
其背景图片（即头像）即最初版奶蛙形象；
“奶龙”智能体因 AI 渲染导致
体型畸变——腹部膨大、四肢细小、形似青蛙——
故得名“奶蛙”
\conf{C}\cite{lostmedia2026naiwa}。
该智能体“已经无法搜索及分享，
评论区也已关闭”
\conf{C}\cite{lostmedia2026naiwa}。

第一版教材记录的考据成果与之相互印证：
“2026 年 8 月 5 日，
被以 B 站 UP 主看仙碎为首的一大批
奶蛙爱好者解决（视频号 BV1ndMk65EwM）”
\conf{C}\cite{shoutanji2026naiwa}。
两个来源（维基词条与第一版教材/评述文章）
在核心事实上\keyw{一致}：
来源指向星野平台、
“纲手【百豪】”账号、
2024-10-02 的智能体、
2026-08 的确认时点。
一致性本身提高了指向的可信度——
但请注意，它提高的是
“两个 C 级来源指向同一处”的可信度，
而不是“该处即起源”的证明力。

\begin{nlhistory}[史料与考据]
发现之后，考据并未“结案”，而是进入了\keyw{证伪与复现}阶段。
 词条转述了 B 站用户“321BBnm”复现视频的结论：
 “‘奶蛙’生成来源仍存在诸多疑点，虽推测其诞生于星野 App，
 但该智能体已无法搜索，
 无法确认是星野 AI 直接生成还是其他 AI 生成后上传。
 多种生成路径的推测均被证伪，其多次复现也未成功，结合老用户所说：
 不管是现在还是原来，星野 AI 本身能力有限，
 无法生成出与参考图契合的图片。最终‘星野 App 出现 bug’的结论仅为推测，
缺乏直接证据，且相关智能体已无法访问，
 生成来源的信息完整性和可追溯性降低”\conf{C}\cite{lostmedia2026naiwa}。
\end{nlhistory}

这段引文是本章的核心教学材料。它示范了网络考据最诚实的写法：
 在给出指向的同时，逐条列出指向所依赖的每一个薄弱环节。
 本书在引用它时，连同这些限定一起引用——
 \keyw{不替考据者删掉他们自己写的“但是”}。

\section{不可复核性：一个认识论困境}\label{sec:04:irrevocable}

把第三阶段的引文推到底，会得到一个一般性的问题：当关键证据载体消失时，
 “来源指向”还能不能当作知识使用？

本书的立场分三层：

\begin{enumerate}[label=\arabic*.]
  \item \textbf{指向可以保留，证明力要降级}。
“考据指向星野平台某智能体”这一陈述，其真值依赖于一个已不可访问的对象；
 但支撑它的材料（截图、评论、时间戳）多数仍可公开检索。本书的做法：
 保留指向，把结论措辞从“奶蛙来自 X”降为“奶蛙最可能的来源指向 X，
 且该指向目前不可复核”。
  \item \textbf{不可复核性不是孤例，
而是网络文化的常态属性}。账号注销、智能体下架、视频删除、平台迁移——
 每一个环节都会让某段传播链的中间环节消失。
 第五章的传播模型把这一现象形式化为“证据损耗率”（M 级参数），
 以说明为什么网络文化研究必须以“现存样本”而非“原始总体”为对象。
  \item \textbf{对教学而言，未结案是优点}。
一个被完整证伪或完整证实的起源故事只能被转述；
 一个未结案的故事可以被用来训练证据分级、置信度措辞、
 复现与证伪的概念——这正是本章的存在理由。
\end{enumerate}

\begin{nltermbox}{术语：不可复核性（irrevocability，本书用法）}
指关键证据载体在指向形成\keyw{之后}消失，
 导致指向无法被独立第三方重新验证的状态。与“来源不明”（从未找到载体）
 相对：不可复核性意味着载体曾经被找到、被记录、随后消失。
 奶蛙智能体 \fname{zNCvQDKfXy} 是本书讨论的标准案例。
\end{nltermbox}

\section{第二证据链：捧腹大笑的音源}\label{sec:04:audio}

形象来源与声音来源是
两条独立的证据链，
本书严格分开处理。

按第一版教材的记录：主播“王麻衣”在直播时，一名自称“布洛芬”
的青年用户连麦，“连麦之后，青年不知为何一直发出异常爽朗的大笑声，
 把所有人都搞破绷了”；
 “如同黑雷霆的目力咆哮与野兽先辈一个一个一个表情的搭配一样，
 这经典的笑声便流传下来，成为了奶蛙捧腹大笑的灵魂音色”
\conf{C} （第一版教材原文，本书转引，级别维持 \conf{C}）。
 二手评述补充了传播端的描述：
 奶蛙二创视频“经常会与变调版不知火舞语音搭配”，
 二者合一后“在不同身份、不同职业的奶蛙们会一同捧腹大笑”
的固定结尾中反复出现 \conf{C}\cite{shoutanji2026naiwa}。

两条材料合起来，可以得到一个 \keyw{“声音—形象合流”} 的传播学描述：
 一个与形象无关的直播笑声素材，经过变调处理后，
 与奶蛙形象在视频二创的剪辑环节发生绑定；绑定之所以稳定，
 是因为“结尾集体大笑”
 这个形式本身具有独立的可传播性（第九章将把它当作“可传播性单位”分析）。

必须同时写明这条证据链的全部脆弱性：音源的具体直播场次、原始录音、
 “变调”发生的具体环节，目前均无可检索的直接材料；“王麻衣”“布洛芬”
 两个名称均出自单线来源（第一版教材）；本书对其的处理是：
 作为 \keyw{“社群追溯到的音源指向”} 记录，级别 \conf{C}，
 不作为任何经验结论的前提。第六章将给出“捧腹大笑”
作为 \keyw{声音符号}的独立分析——那部分分析不依赖音源考据是否成立。

\section{2026 年 3 月：关评、拉黑与“潘多拉奶盒”}\label{sec:04:official}

官方介入的时间点必须单独列出，因为它同时是产业线（第三章）
 与现象线的 \keyw{分界事件}。

第一版教材的表述：“2026 年 3 月，奶龙官方呼吁停止丑化二创，
 关闭视频评论区并批量拉黑发奶蛙的账号。
 但是奶蛙的影响力已经扩散开来（潘多拉奶盒打开了），难以制止了。
 至此，奶蛙凭借着魔性而意义不明的捧腹大笑，正式脱离了奶龙，
 成为了互联网上一个强而有力的热梗” \conf{C}（第一版教材原文，转引，
级别维持 \conf{C}）。

二手评述的交叉材料：“奶龙官方对于奶蛙梗已经深恶痛绝。
 之前网友一起玩梗还能被法不责众的幌子遮掩一下。
 而这时当奶蛙的源头有了明确的指向性时——
 只能说奶龙法务部的拳头想喝血了”（评述者语）
\conf{C} \cite{shoutanji2026naiwa}；同一评述同时记录：
 “直至今日你依然可以在奶龙官方视频的评论区看见无数奶蛙图片”
 \conf{C}\cite{shoutanji2026naiwa}。

本书从这三条材料中提取两个分析判断（M 级，非经验断言）：

\begin{enumerate}[label=\arabic*.]
  \item \textbf{抑制是“不彻底”的}：
关评、拉黑针对的是官方账号下的评论区与发布奶蛙的特定账号，
 而不是奶蛙形象本身在整个互联网上的流通——
 “官方视频评论区仍有无数奶蛙图片”
 说明抑制的执行在空间上就是不完整的。
  \item \textbf{“源头指向性”改变了
抑制的动机结构}：评述者把法务动作的时间点与“源头考据完成”并置——
 在来源不明时，抑制的对象只能分散到无数二创作者（“法不责众”）；
 来源被指向特定账号后，抑制才有了可执行的目标。
 这一并置目前只有 \conf{C} 级支持，
 本书只把它作为“待检验的时间相关假设”
 保留（第七章 \ref{sec:07:reverse} 的逆反认证机制以此为经验入口）。
\end{enumerate}

“潘多拉奶盒”这个社群用语本身也值得记录：它把“不可逆扩散”
 表达为“盒子打开了就关不上”，与“正式脱离了奶龙”的判断合起来，
 构成了民间叙事中关于“奶蛙—奶龙”关系的 \keyw{断裂叙事}——
 与正典产业线（第三章：授权、维权、关评）构成的“连续叙事”相对照。
 第九章将论证：两种叙事不是对错关系，而是分别描述了“法律层面”
 与“意义层面”的正典—变种关系（第二章定义 2.1 的两个层面）。

\section{考据伦理：关于“创作者”的写作边界}\label{sec:04:ethics}

本章最后一节处理
一条\keyw{硬红线}。

奶蛙起源考据指向了一个具体的、匿名的个体账号。考据完成之后，
 围绕“创作者”的叙述出现了两种危险倾向，本书都必须拒绝：

\begin{enumerate}[label=\arabic*.]
  \item \textbf{“认领叙事”的
伦理风险}：考据指向某个账号后，
 舆论场会自然地要求“当事人出面认领”，
 或在当事人未认领时进行身份推断。本书的立场：
 考据指向的是 \keyw{账号} （一个可公开检索的网络对象），
 不是 \keyw{自然人}。本书不对账号持有者的现实身份作任何推断，
 不呼吁任何个体“认领”或“否认”，
 也不报道任何关于其个人境况（包括健康、职业、情绪、生死）
 的推测性内容——二手评述中确实存在以“最残酷的结局”为标题、
 含个人境况推测的叙述 \conf{C} \cite{shoutanji2026naiwa}，
 本书 \keyw{不转述} 其内容，只在此说明其存在及本书的处理方式。
  \item \textbf{“工具化”的
研究伦理风险}：反过来，也不能因为当事人匿名，就把“创作者”
 当作纯粹的证据节点使用。
 第二章已说明奶蛙的署名权归属原始创作者（附录 D）；
 本书在全部叙述中使用“考据指向的创建者”这一措辞，
 保留其作为 \keyw{权利主体} 的地位，
 而不只作为 \keyw{证据节点} 的地位。
\end{enumerate}

\begin{nlmethod}[伦理操作的三条规则]
\textbf{规则 E1（对象限定）}：考据叙述的宾语只能是账号、智能体、
 帖子、截图等可公开检索对象；不得以“创作者本人”
 为宾语陈述任何个人信息。 \textbf{规则 E2（推测隔离）}：
 对个人境况的推测性材料（无论来自媒体还是社群），
 一律“说明存在、不转述内容”处理。 \textbf{规则 E3（权利保留）}：
 涉及二创形象时，同时写明署名权归属，使“分析”
 不滑向“消费”。
\end{nlmethod}

这三条规则在第十二章会进入“可复现性清单”，
 成为修订者必须遵守的操作条款。在此提前说明，
 是因为本章的全部材料正是它们的 \keyw{第一个适用案例}。

\section{本章小结}\label{sec:04:summary}

本章把奶蛙起源考据拆成了三个阶段：
 误指（ 2024-03 “季洪行说”）、线索（ 2024-12 星野水印）、
 发现与证伪并存（ 2026-08 智能体确认，
 ID \fname{zNCvQDKfXy} 已不可访问）。
 核心结论是一个 \keyw{降级的指向}：
 奶蛙最可能的来源指向星野平台 2024-10-02 的智能体头像，
 该指向由两个 C 级来源交叉印证，但因载体消失而不可复核，
 且考据材料自身声明“仅为推测、缺乏直接证据”。
 音源证据链（王麻衣连麦）作为第二条独立链单独记录，
 级别同为 C 级。 2026-03 的关评拉黑事件是两条产业线的分界，
 本书从中提取了“抑制不彻底”与“指向改变动机”
 两个待检验判断。最后，
 本章给出了三条考据伦理规则（ E1/E2/E3），并以奶蛙“创作者”
 为第一个适用案例——这是全书证据分级制度最严格的一次执行，
 也是第十二章方法论的预演。遗留问题：
 起源考据回答了“形象从哪来”（部分地），
 但没有回答“为什么是这个畸变被选中、
 为什么是这个笑声被绑定”——那是符号问题与机制问题，第六、
 七章的主题。

\begin{nlexercises}
  \item 按定义~\ref{def:04:kaoju} 的结论形态区分，改写以下两句：
 “奶蛙来自星野 AI ”
 与“奶蛙最可能的来源指向星野 AI 的某智能体，
 指向目前不可复核”。说明两句在证据要求上的差异，
 以及为什么第一句在当前材料下不成立。
  \item 给
出一个你熟悉的其他网络梗的“误指 → 线索 → 发现”
 三阶段重构（若该梗起源已公认，
 说明哪一阶段被跳过及原因）。
  \item 规则
E2 要求“说明存在、不转述内容”。
 设计一条可执行的检查程序（如关键词清单 + 句式模板），
 使一位修订者能在 5 分钟内自查修订稿是否违反 E2。
\end{nlexercises}

\nlkeywords 考据、未结案证据链、误指、不可复核性、音源、
关评拉黑、考据伦理、E1/E2/E3

\chapter{迷因理论框架：复制、变异与选择}\label{ch:05}

\begin{nlkeypoints}
  \begin{enumerate}[label=(\arabic*),leftmargin=2.4em]
    \item 迷因（\en{meme}）在本书中的操作化用法：
          一个在头脑间复制的文化单元，经历\keyw{变异—选择—保留}
          三阶段；分析单位是“结构相似度高于社群默许阈值的用例簇”，
          不是某个具体文件。
    \item Shifman 三维框架（\keyw{内容 / 形式 / 立场}）
          是第六章解剖奶蛙的工具；
          “迷因作为争夺中的文化资本”解释了
          为什么“看不懂”本身就是传播力的一部分。
    \item 本书自建了一个\keyw{SIR 型迷因传播教学模型}
（M 级，参数为示意值）：用“易感—传播—疲劳”三态说明热度为何必然衰减、
以及“二次爆发”需要何种条件（新变体、新工具或免疫衰减）。
  \end{enumerate}
\end{nlkeypoints}

前四章处理的是对象与史实。本章切换任务：
给全书配一套\keyw{分析工具}。
工具只有一件半：一件是迷因理论
（从道金斯到 Shifman 的概念链），
半件是本书自建的传播教学模型。
之所以说“半件”，是因为后者的地位
在前言与第一章已经声明：
它是\keyw{教学脚手架}（M 级），
不是经验发现——
它的作用是让“为什么会衰减”
“为什么会二次爆发”这类问题
变得\keyw{可形式化}，
从而让第六章的符号分析与
第七章的机制假设
有一个共同的时间骨架。

\section{复制子：meme 的概念源}\label{sec:05:dawkins}

道金斯在《自私的基因》第 11 章提出 \en{meme}：“文化就像基因组一样，
 由大量复制的单位组成……这些单位我们不妨称为 meme”
 \cite{dawkins1976selfish}。定义的关键不在类比本身，
 而在它携带的\keyw{三个判据}： (1) 可复制——
能在个体间被忠实程度不等的传递； (2) 可变异——每次复制都可能引入差异；
 (3) 受选择——环境（受众偏好、平台算法、社群规范）决定哪些变体留存。

《延伸的表现型》把“复制子”的视角从基因延伸到更一般的“文化结构”
\cite{dawkins1982extended}。本书取用这一视角，但立刻做两处限定：

\begin{enumerate}[label=\arabic*.]
  \item \textbf{不主张本体论}。
“奶蛙是一个 meme”
 这句话在本书中 \keyw{不是一种存在论断言} （好像网络文化中存在一个叫 meme 的实体），
 而是\keyw{方法论约定}：我们\keyw{假装}存在可复制、可变异、
受选择的文化单元，以便使用“变异率”“选择压力”
 这类有预测力的分析语言。这一约定与第二章“正典—变种”结构是兼容的：
 正典奶龙是“高保真复制的稳定吸引子”，奶蛙是“低保真复制的变异簇”——
 两者都是复制子视角下的不同动态。
  \item \textbf{不追求“同一性”判定}。
一万条奶蛙视频是不是“同一个 meme”？本书不回答本体论问题，
 只给出\keyw{操作化归并标准} （定义~\ref{def:05:cluster}）。
\end{enumerate}

\begin{definition}[迷因簇（meme cluster，本书用法）]
  \label{def:05:cluster}
设 $\mathcal{U}$ 为可检索的奶蛙相关用例集合（视频、图片、帖子）。
 若两个用例 $u, u' \in \mathcal{U}$ 满足 \textbf{(C1)} 形象结构相似（保留：
黄色主体、大笑姿态；允许：体型比例、细节装饰差异），
 \textbf{(C2)} 形式要素重叠（使用相同或近似的音轨、结尾套路、
 衍生语之一），
 \textbf{(C3)} 社群可辨识（随机抽样社群成员在提示下能识别“这是奶蛙”），
 则 $u$ 与 $u'$ 归入同一迷因簇。归并是\keyw{传递闭包} 关系：
$u \sim u'$、 $u' \sim u''$ 推出 $u \sim u''$。
\end{definition}

注意 定义~\ref{def:05:cluster} 把“归并”
从美学判断（“像不像”）改写成三条可执行的检查项，
 其中 (C3) 最容易被忽略，也最关键：
 \keyw{归并的最终裁判不是编辑部的相似度直觉，而是社群的可辨识性}。
 第六章的谱系表将用这三条标准逐类检验。

\section{三阶段：变异、选择、保留}\label{sec:05:stages}

把 定义~\ref{def:05:cluster} 的归并标准反过来用，
 就得到奶蛙迷因簇内部结构的 \keyw{三阶段动力学}：

\subsection{变异}

奶蛙的变异有明确的“变异引擎”： AIGC 工具（第八章详论）。
 低精度图像模型的结构性畸变（腹部膨大、四肢细小）
 提供了\keyw{形象层面的变异源}；
 视频生成模型（二手评述提及的抖音 Seedance 等 \conf{C}\cite{shoutanji2026naiwa}）
 把静态形象变为可剪辑的“角色演员”；
 音轨的变调处理提供了 \keyw{声音层面的变异源}。
 变异的关键特征是\keyw{方向性不对称}：
 从正典奶龙到奶蛙是“一次大畸变 + 持续小调整”，
 从奶蛙回到正典几乎不发生（没人会把奶蛙“修回”奶龙再传播）。
 这种不对称是第六章“内容层”分析的直接输入。

\subsection{选择}

选择压力
来自三处：

\begin{enumerate}[label=\arabic*.]
  \item \textbf{平台算法}：
完播率、转发率等指标筛选“结尾够不够好笑”；
  \item \textbf{社群规范}：
“太正经”的奶蛙不被转发（抽象文化的准入测试，
 第二章定义~\ref{def:02:abstract}）；
  \item \textbf{认知经济性}：
形象必须“一眼可辨”——畸变不能超过可辨识阈值，否则 (C3) 失败。
\end{enumerate}

三者共同产生了一个可观察推论：
 \keyw{奶蛙形象的变异率在“可辨识”与“足够怪”之间存在一个窄通道}——
 这解释了为什么不同 UP 主独立生产的奶蛙在轮廓上高度趋同（都保留了黄色 + 大笑 + 膨大腹部），
 而在细节上高度发散（换皮、物化、娘化）。
 第六章将把这条推论落进谱系表。

\subsection{保留}

保留（retention）指变体进入“社群记忆”的过程：
 衍生语（“安迪”“无了”）进入词汇表，衍生动作（思考、
 蹲起大笑、爬行、摆架势）进入身体程式，立牌进入校园空间。
 保留的单位是 \keyw{“可复用资产”}：
 一个被保留的变体是未来变异的 \keyw{原料}——“奶蛋”
 之所以存在，正因为“奶蛙”本身先被保留。
 三阶段因此不是直线，而是 \keyw{ 飞轮 } ：保留喂养变异，
 变异经过选择，选择的产物再被保留。第七章的“压力—门槛—
放大器”结构就是这个飞轮的社会机制版本。

\begin{figure}[htbp]\centering
\begin{tikzpicture}[font=\small,
  every node/.style={align=center,inner sep=3pt},
  stage/.style={rounded corners,draw=NLDeep,thick,
    fill=NLYellow!18,minimum height=1.5cm,minimum width=3.2cm},
  arr/.style={-{Stealth[length=3mm]},thick,NLDeep}]
  \node[stage] (var) {变异\\{\footnotesize AIGC 畸变、变调、换皮}};
  \node[stage,right=2.6cm of var] (sel) {选择\\{\footnotesize 算法、社群规范、认知经济性}};
  \node[stage,below=1.6cm of sel] (ret) {保留\\{\footnotesize 衍生语、动作、立牌进入社群记忆}};
  \draw[arr] (var) -- node[above]{变体} (sel);
  \draw[arr] (sel) -- node[right]{通过者} (ret);
  \draw[arr] (ret) -- node[below,rotate=0]{可复用资产} (var);
  \node[font=\footnotesize,align=left] at (1.2,-1.6) {（飞轮：保留喂养变异）};
\end{tikzpicture}
  \caption{迷因三阶段飞轮（M 级教学框架；
          “变体/通过者/可复用资产”
          为本书的操作化命名）。}
  \label{fig:05:wheel}
\end{figure}

\section{三维框架：内容、形式、立场}\label{sec:05:shifman}

Shifman 把数字迷因的分析从“内容像什么”推进到“它在做什么”：
 每个迷因用例同时携带 \keyw{内容} （content，形象、文本、
 声音等可描述素材）、 \keyw{形式} （form，媒介、体裁、结构套路）
 与 \keyw{立场} （stance，使用者借迷因表达的态度、情绪、
 社交姿态）三个维度 \cite{shifman2013memes}。
 三维的价值在于它把“为什么好笑”
 拆成了三个可分别提问的问题：内容——“它画的是什么”；
 形式——“它按什么套路组织”；立场——“它在替谁说什么话”
 。

在奶蛙语料中，三维的不对称是显著的：
 \keyw{内容高度可替换} （换皮、物化、任何“黄色大笑生物”
 都能接进来）， \keyw{形式高度稳定} （结尾集体大笑、
 衍生语搭配、“压力桥段 → 大笑”的叙事结构），
 \keyw{立场高度约定化} （使用奶蛙 = 声明“我绷不住了，
 但我选择用笑来处理”）。第六章将对这一“一软两硬”
 的结构做完整解剖；本章只给出框架本身。

\section{迷因作为争夺中的文化资本}\label{sec:05:capital}

Nissenbaum 与 Shifman 对 4chan /b/ 版的研究把迷因的“立场”
维度推进了一步：迷因不是“表达”某种预先存在的态度，
 而是 \keyw{争夺文化资本的手段}——谁先用、谁用得更准、
 谁能识别“过期梗”，
 都在持续地重新分配社群内部的地位 \cite{nissenbaum2015contested}。
 “文化资本”在这里是 \keyw{可被社群审计的东西}：
 你转发过什么、用过哪个衍生语、对哪个“过期梗”保持沉默，
 都构成他人可读取的地位信号。

这一框架对本书的直接意义是两条：

\begin{enumerate}[label=\arabic*.]
  \item \textbf{“看不懂”
不是传播的失败，而是传播的准入测试}。
 抽象文化的意义悬置（第二章定义~\ref{def:02:abstract}）
 在资本框架下重新解释：
 它制造了一个只有圈内人能通过的认证环节，
 而通过认证的回报是地位信号本身。第七章的“群体稀释”
 机制将用到这一点：用梗代替自述，既表达了情绪，
 又完成了一次低成本的身份认证。
  \item \textbf{逆反
是资本游戏的合法玩法}。官方关评、
 拉黑（第四章 \ref{sec:04:official}）
 在资本框架下读作“权威的否认”，
 而被权威否认的梗获得“含金量”背书——
 这就是第一版教材所说的“逆反认证”的学理版本，
 也是第七章 \ref{sec:07:reverse} 机制的理论入口。
\end{enumerate}

\begin{nlcritique}[批判性反思]
资本框架的危险在于它可能把一切社群行为都读成“地位游戏”，
 从而取消了“人们就是单纯觉得好笑”这一朴素解释的合法性。
 本书的处理：第七章把“纯消遣”
 保留为竞争性假设的一部分（“低成本娱乐”解释），
 并要求资本框架给出自己的可观察预测（例如“梗的地位信号价值随其普及度上升而下降”），
 否则它只是一种解释风格，不是一个可检验的理论。
\end{nlcritique}

\section{SIR 型迷因传播教学模型}\label{sec:05:sir}

\begin{nlmethod}[研究方法提示：模型地位声明]
本节模型为\keyw{本书构造的 M 级教学模型}：所有参数取值为示意值，
 不声称拟合任何平台的真实数据；模型的作用是（i）把“热度衰减”
“二次爆发”从形容词变成有明确条件的命题；（ii）
 为第七章的假设提供统一的时间骨架。修订者可以替换参数、替换状态定义，
 但必须保留“每个状态转移都有可指认的社会机制”这一约束。
\end{nlmethod}

\subsection{三态定义}

把受众分为三个状态：

\begin{itemize}[leftmargin=2.2em]
  \item $S(t)$：\keyw{易感者}（susceptible）——
        尚未深度接触迷因、
        仍有概率被“激活”的人群；
  \item $I(t)$：\keyw{传播者}（infectious）——
        当前处于“主动使用/转发/二创”
        状态的人群；
  \item $R(t)$：\keyw{疲劳者}（recovered 的挪用，
本书读作“疲劳”）——已深度接触但“笑不动了”、不再主动传播的人群。
\end{itemize}

借用流行病学 SIR 模型的状态命名有明确理由：“疲劳”与“康复”
在结构上同型（都退出 $I$ 态），而“疲劳不是免疫”（可被重新激活）
 与下面引入的免疫衰减项对应。总人数 $N = S+I+R$ 取常数。

\subsection{方程组}

\begin{equation}\label{eq:05:sir}
\begin{cases}
\displaystyle
\frac{dS}{dt} = -\,\beta\,\frac{S\,I}{N} \;+\; \omega\, R,
\\[1.2ex]
\displaystyle
\frac{dI}{dt} = \;\beta\,\frac{S\,I}{N} \;-\; \gamma\, I,
\\[1.2ex]
\displaystyle
\frac{dR}{dt} = \;\gamma\, I \;-\; \omega\, R.
\end{cases}
\end{equation}

各参数的社会机制对应：

\begin{table}[htbp]\centering
  \caption{模型参数的机制对应（M 级教学模型，
          取值为示意值，日为单位）}\label{tab:05:params}
  \small
  \begin{tabularx}{\textwidth}{@{}lXll@{}}
    \toprule
    参数 & 机制解释 & 示意取值 & 可检验的替换 \\ \midrule
    $\beta$ & 接触—激活率：
      一个传播者每日“激活”
      一个易感者的期望数
      （对应机制：
      平台推送、
      群聊复读、
      线下偶遇） & $0.35$
      & 用“每接触者日均转发数”估计 \\
    $\gamma$ & 退出传播率：
      从“主动使用”退出到“疲劳”
      的速率（对应机制：
      疲劳、被嘲、
      “梗过期”的自我感知） & $0.20$
      & 用“梗的半衰期”估计 \\
    $\omega$ & 免疫衰减率：
      疲劳者被
      \keyw{新变体 / 新工具 /
      新情境}
      重新激活的速率
      （对应机制：
      换皮新形象、
      视频生成工具上线、
      官方事件） & $0.05$
      & 用“二次爆发的触发事件”估计 \\ \bottomrule
  \end{tabularx}
\end{table}

基本再生数 $R_0 = \beta/\gamma$
（示意参数下 $R_0 = 1.75 > 1$，
故初期会上升）。
模型的两个核心命题：

\begin{proposition}[必然衰减]\label{prop:05:decay}
在 $\omega = 0$（无免疫衰减）的退化情形下，
 方程 \eqref{eq:05:sir} 的 $I(t)$ 单峰：热度必然从峰值回落。
\end{proposition}
\begin{proof}[证明思路]
当 $S/N$ 降到 $\gamma/\beta$ 以下时， $dI/dt < 0$ 恒成立；
 而 $S$ 单调不增（$\omega=0$ 时 $dS/dt \le 0$），
 故 $I$ 一旦转为下降便不再回升。完整推导留给习题 5-2。
\end{proof}

\begin{proposition}[二次爆发的条件]\label{prop:05:second}
若存在一个时刻 $t_0$，
 使 $S(t_0)/N > \gamma/\beta$ （例如由新变体把部分 $R$ 经 $\omega R$ 重新注入 $S$，
 或由外部新人群流入），则 $I(t)$ 在 $t_0$ 之后可以出现第二个峰。
 二次爆发的“弹药”
 是\keyw{未被耗尽的易感存量} 与\keyw{足够强的重新激活项}。
\end{proposition}

两条命题合起来就是“衰减与二次爆发”的结构解释：奶蛙在 2024—
2025 年“小范围传播”
后的 2026 年爆发（第四章 \ref{sec:04:discovery} 的时间线），
 在模型语言里对应一次 $\omega$ 被“新工具 + 新形象 + 官方事件”
 共同拉高的重新激活过程（\conf{C} 级时间线为 M 级模型提供事件顺序，
 模型本身不提供因果断言）。

\subsection{示意曲线}

\begin{figure}[htbp]\centering
\begin{tikzpicture}
\begin{axis}[
  width=0.92\textwidth, height=6.4cm,
  xlabel={时间（天，示意刻度）},
  ylabel={$I(t)$，传播者规模（示意归一化）},
  xmin=0, xmax=120, ymin=0, ymax=1.05,
  xtick={0,10,20,...,120},
  grid=major, grid style={NLGrey!30},
  legend pos=north west,
  every axis plot/.append style={thick},
]
\addplot[NLDeep,samples=200,smooth]
  {0.55*exp(-((x-14)^2)/60) + 0.95*exp(-((x-72)^2)/110)};
  \addlegendentry{$I(t)$：双峰（教学示意）}
\addplot[NLYellow!80!black,samples=200,smooth,thick,dashed]
  {0.55*exp(-((x-14)^2)/60)};
  \addlegendentry{第一峰成分}
\addplot[NLAmber,samples=200,smooth,thick,dashdotted]
  {0.95*exp(-((x-72)^2)/110)};
  \addlegendentry{第二峰成分（重新激活后）}
\end{axis}
\end{tikzpicture}
  \caption{迷因传播双峰示意曲线（M 级教学图；
曲线为两个高斯峰的叠加，不是对任何平台数据的拟合，
 也不是对方程 \eqref{eq:05:sir} 的数值解——
 其作用是展示 命题~\ref{prop:05:second} 的定性形态）。
 第一峰对应 2024—2025 的小范围传播期，
 第二峰对应 2026 年全面爆发期（时间线见附录 \ref{app:timeline}，
 可信度 \conf{C}）。}
  \label{fig:05:sir}
\end{figure}

\begin{nlcritique}[批判性反思：模型的边界]
该模型至少有三处明确的局限，修订者不得“悄悄修补”
后当作经验结论引用： (1) \textbf{单群同质假设}：它把“高校”“职场”
 “家庭”混在一个 $N$ 里，
 而校园与平台的接触率差异正是本书关心的核心；
 第十章的案例因此需要多群版本。 (2) \textbf{状态不可直接观测}：
 $I$ 态的“主动使用”无法从公开数据直接读出，
 任何用“转发量”代理 $I$ 的做法都必须声明代理误差。
 (3) \textbf{参数无识别性}：
 $\beta,\gamma,\omega$ 的示意取值之间存在大量“等价”的参数组，
 模型给出的是 \keyw{形态约束} （单峰 / 双峰、峰值先后），
 不是定量预测。
\end{nlcritique}

\section{本章小结}\label{sec:05:summary}

本章配好了全书的工具箱：
 概念链（道金斯复制子 → Shifman 三维 → 争夺中的文化资本），
 操作化装置（迷因簇归并标准， 定义~\ref{def:05:cluster}；
 三阶段飞轮，图~\ref{fig:05:wheel}），
 以及时间骨架（ SIR 型教学模型，方程 \eqref{eq:05:sir}，
 图~\ref{fig:05:sir}，含两条命题与三条模型边界声明）。
 工具与材料到此齐备：接下来两章开始“解剖”——
 第六章按三维框架处理奶蛙的符号系统，第七章把“为什么笑”
 重写为可检验的假设。必须重申的是：
 本章所有模型构件均为 M 级教学装置，
 其价值在于让问题可形式化，
 而不在于对任何平台数据做定量断言。

\begin{nlexercises}
  \item 用
定义~\ref{def:05:cluster} 的 (C1)(C2)(C3) 三条标准，
 判断以下用例与“经典捧腹大笑奶蛙”是否归入同一迷因簇，
 并说明哪一条标准起了决定作用：
 (a) 把奶蛙形象 P 成拿着高等数学课本的表情包； (b) 一只黄色、
 大笑但体型正常的龙； (c) 完全不同形象、
 但使用同一音轨与结尾套路的视频。
  \item 完成
命题~\ref{prop:05:decay} 的完整推导：
 证明当 $\omega=0$ 且 $S(0)>0$ 时， $I(t)$ 至多一个极大值。（提示：
 考察 $dI/dt=0$ 的条件与 $S(t)$ 的单调性。）
  \item 给
出一个可执行的方案：
 若你能获得某平台 90 天的奶蛙相关内容的“日均新增量”序列，
 你将如何用它来区分“单峰”与“双峰”
 两种形态（写出判据与对“噪声”的稳健化处理；
 说明该方案不能识别什么）。
\end{nlexercises}

\nlkeywords 复制子、迷因簇、变异—选择—保留、
三维框架、文化资本、SIR 型教学模型、$R_0$、
二次爆发、免疫衰减

\chapter{奶蛙的语言—符号系统：内容、形式与立场}\label{ch:06}

\begin{nlkeypoints}
  \begin{enumerate}[label=(\arabic*),leftmargin=2.4em]
    \item 以 Shifman 三维框架解剖奶蛙语料，得到一个\keyw{“一软两硬”}结构：
          内容维度高度可替换（任何“黄色大笑生物”都能接入），
          形式维度高度稳定（结尾集体大笑、衍生语搭配、
          “压力桥段$\to$大笑”的叙事套路），立场维度高度约定化
          （使用奶蛙即声明“我绷不住了，但我选择用笑来处理”）。
    \item 形象谱系遵循\keyw{窄通道规律}：轮廓趋同（黄色主体、大笑姿态、
          膨大腹部），细节发散（换皮、物化、娘化）。
          谱系表按“保留特征 / 丢失特征”两列组织，
          判定依据是第五章迷因簇归并标准 (C1)--(C3)。
    \item “捧腹大笑”是一个\keyw{声音符号}而非单纯的情绪表达：
          它的分析不依赖音源考据是否成立；
          衍生语是\keyw{空耳}（\en{soramimi}）驱动的语音梗，
          其“语音外壳保留、语义外壳被重新指派”的双层结构
          是奶蛙语言系统的核心装置。
  \end{enumerate}
\end{nlkeypoints}

第五章配好了工具，本章开始解剖。解剖对象是问题 Q1：
奶蛙究竟是一个怎样的符号系统？本章的材料是三类语料：
形象语料（谱系与变种）、语言语料（衍生语与固定音轨）、
动作语料（衍生动作与叙事套路）。除正典设定来自 A/B 级来源外，
这三类语料几乎全部来自社群考据与第一版教材的记录（\conf{C} 级），
因此本章的每一张谱系表都附带“判定依据”列，
读者可以按第五章的标准 (C1)--(C3) 自行复核归并是否成立。
本章的结论将是一个结构性判断：奶蛙的符号系统呈
“一软两硬”的不对称——这个判断本身是本书的分析（\conf{M} 级），
但它解释了第八章将要观察到的现象：
为什么 AIGC 时代的二创在形象上千姿百态，
在形式上却高度重复。

\section{分析框架：三维的操作化}\label{sec:06:framework}

Shifman 把数字迷因的分析从“内容像什么”推进到“它在做什么”：
每个迷因用例同时携带\keyw{内容}（\en{content}）、
\keyw{形式}（\en{form}）与\keyw{立场}（\en{stance}）
三个维度\cite{shifman2013memes}。框架本身是 A 级理论资源；
把它落到奶蛙语料上，需要先做一步操作化：

\begin{table}[htbp]\centering
  \caption{三维框架在奶蛙语料上的操作化（判定问题均为可观察问题）}
  \label{tab:06:ops}
  \small
  \begin{tabularx}{\textwidth}{@{}llX@{}}
    \toprule
    维度 & 操作化提问 & 奶蛙语料中的对应物 \\ \midrule
    内容 & “它画的是什么？” &
      形象（黄、膨腹、细肢、黑手）、场景（校园、职场、
      游戏跨界）、被嫁接的文本（课程、考试、品牌） \\
    形式 & “它按什么套路组织？” &
      结尾集体大笑、变调笑声音轨、衍生语搭配、
      “压力桥段$\to$大笑收尾”的叙事结构、竖屏短时长 \\
    立场 & “它在替谁说什么话？” &
      “绷不住但用笑处理”的自陈、自嘲、
      对精致人设的疏离、圈内身份的默示声明 \\ \bottomrule
  \end{tabularx}
\end{table}

操作化的纪律是：每一维的判定问题必须能落到
“看一段二创就能回答”的层面。“立场”最容易被写成印象式的
“反映了当代青年的心态”，本书拒绝这种写法——
立场维度只讨论\keyw{可指认的语用行为}：
用梗图代替自述（第七章机制 M2）、
在何种情境下使用哪个衍生语（本章 \ref{sec:06:form} 节）、
以及社群自己写下的立场声明文本（\ref{sec:06:stance} 节）。

\section{内容维度：形象谱系学}\label{sec:06:content}

\subsection{谱系总表}

奶蛙的形象不是一张图，而是一棵树。按第一版教材的整理与
本书的复核（\conf{C} 级，来源为第一版教材与社群考据词条
\cite{lostmedia2026naiwa}），谱系按“保留多少正典特征”
分为三级：换皮、异形、客串。

\begin{table}[htbp]\centering
  \caption{奶蛙形象谱系（据第一版教材整理，\conf{C} 级；
          “保留 / 丢失”两列即归并标准 (C1) 的操作化）}
  \label{tab:06:pedigree}
  \small
  \begin{tabularx}{\textwidth}{@{}llXX@{}}
    \toprule
    类别 & 代表 & 保留特征 & 丢失 / 替换特征 \\ \midrule
    换皮奶蛙 & 异色奶蛙 & 全部形体结构 &
      仅配色变化 \\
    & “植奶”“僵奶”（《植奶大战僵龙》系列） &
      形体结构、动作程式 &
      皮肤替换为《植物大战僵尸》元素 \\
    & 线代课本奶蛙、VIVO 奶蛙、学习通奶蛙 &
      轮廓与姿态 & 表面替换为图标、封面等平面元素 \\
    异形奶蛙 & 奶蛋 & 眼睛与大笑 &
      躯干、四肢全部简化 \\
    & 奶狗、奶鸡、奶鼠、奶茶、豆奶 &
      眼睛与大笑（部分含黄色） & 物种形态整体替换 \\
    & 奶蛙娘 & 面部特征、黑色双手 & 人体化（娘化） \\
    & 物化奶龙（如“奶菜”） & 局部特征 & 拟物化 \\
    客串角色 & 美团袋鼠 & 姿态与笑声 & 伪装为其他 IP / 品牌形象 \\
    & 牛来、牛来妈 & 姿态与笑声 & 独立角色化 \\ \bottomrule
  \end{tabularx}
\end{table}

\begin{definition}[换皮与异形（本书用法）]\label{def:06:reskin}
  \textbf{换皮奶蛙}：仅替换表面元素（配色、图标、封面），
  保留奶蛙全部形体结构的变种。
  \textbf{异形奶蛙}：丢失大量形体特征，
  仅保留“眼睛与大笑”这一最小辨识单元的变种。
  两者以“丢失特征的比例”划界，
  而非以“像不像青蛙”划界——奶蛙娘虽是人形，
  但因完整保留面部特征与黑色双手，
  在社群辨识实践中归入异形而非独立形象。
\end{definition}

这条分界线来自社群的使用实践（归并标准 (C3)），
不是本书的发明：一个形象只要还能被社群成员
在提示下识别为“这是奶蛙”，它就还在谱系内；
识别一旦失败，谱系就到了边界。
谱系因此是\keyw{辨识边界的图谱}，不是审美相似度的图谱。

\begin{figure}[htbp]\centering
\begin{tikzpicture}[font=\small,
  every node/.style={align=center,inner sep=3pt},
  lv/.style={rounded corners,draw=NLDeep,thick,fill=NLYellow!18,
    minimum height=1.1cm,minimum width=2.6cm},
  leaf/.style={rounded corners=1pt,draw=NLGrey,fill=NLBg,
    minimum height=0.8cm,font=\footnotesize},
  arr/.style={-{Stealth[length=2.6mm]},thick,NLDeep}]
  \node[lv] (canon) {正典奶龙\\{\footnotesize 黄色、圆头圆脑（\conf{A/B}）}};
  \node[lv,below=1.0cm of canon] (naiwa) {经典奶蛙\\{\footnotesize AI 渲染畸变：膨腹、细肢（\conf{C}）}};
  \node[lv,below left=1.1cm and 1.4cm of naiwa] (reskin) {换皮奶蛙};
  \node[lv,below=1.1cm of naiwa] (hetero) {异形奶蛙};
  \node[lv,below right=1.1cm and 1.4cm of naiwa] (guest) {客串角色};
  \node[leaf,below=0.55cm of reskin] (l1) {异色 / 植奶僵奶 / 图标换皮};
  \node[leaf,below=0.55cm of hetero] (l2) {奶蛋 / 奶狗等 / 奶蛙娘 / 物化};
  \node[leaf,below=0.55cm of guest] (l3) {美团袋鼠 / 牛来};
  \draw[arr] (canon) -- node[right,font=\footnotesize]{一次大畸变} (naiwa);
  \draw[arr] (naiwa) -- node[above left,font=\footnotesize]{持续小变异} (reskin);
  \draw[arr] (naiwa) -- (hetero);
  \draw[arr] (naiwa) -- (guest);
  \draw[arr,dashed,NLGrey] (reskin) -- (l1);
  \draw[arr,dashed,NLGrey] (hetero) -- (l2);
  \draw[arr,dashed,NLGrey] (guest) -- (l3);
\end{tikzpicture}
  \caption{奶蛙形象谱系树（据第一版教材整理，\conf{C} 级）。
          实线箭头为“变异方向”，均为自上而下；
          第五章“方向性不对称”在谱系上表现为：
          没有从子类回到父类的箭头。}
  \label{fig:06:tree}
\end{figure}

\subsection{窄通道规律}

把表 \ref{tab:06:pedigree} 与图 \ref{fig:06:tree} 合起来看，
可以提炼出一条本书称为\keyw{窄通道规律}的观察（\conf{M} 级，
由 \conf{C} 级语料归纳）：

\begin{proposition}[窄通道规律（教学命题）]\label{prop:06:channel}
  奶蛙变种的存活空间由两条约束夹出：
  其一，\textbf{可辨识下限}——丢失“黄色主体、大笑姿态、
  眼睛”等特征过多时，(C3) 失败，变种脱离谱系；
  其二，\textbf{新奇性下限}——与既有变种完全同质时，
  选择压力（社群的转发行为）不再偏好该变种。
  存活的变种集中于“轮廓趋同、细节发散”的窄通道内。
\end{proposition}

这条命题是“教学命题”：它把第五章
“变异率在‘可辨识’与‘足够怪’之间存在窄通道”的推论
落到了谱系表上。它的可检验形态是：
对任意时间窗内的奶蛙二创样本，
“轮廓特征保有率”的方差应显著小于
“细节特征保有率”的方差。
第十一章给出了这一检验的编码方案。
同时必须写明竞争性解释：轮廓趋同也可能仅仅是
AIGC 工具模板化输出的副产品（第八章），
而非社群选择的结果——两种解释在“工具固定、
社群不变”与“社群固定、工具变化”的对照样本上
可以区分。

\begin{nlcase}[案例 6-1：《植奶大战僵龙》与模因杂交]
  B 站 UP 主“暖气杀手 AAA”的《植奶大战僵龙》系列
  将奶蛙与《植物大战僵尸》（PVZ）相结合，
  制作多款高播放量视频，是奶蛙爆火的“强力引擎”之一
  \conf{C}（第一版教材口径）。
  从符号学上看，这是一次典型的\keyw{模因杂交}：
  奶蛙一方输出形象与动作程式（内容层嫁接），
  PVZ 一方输出玩法与场景框架（形式层嫁接），
  双方各保留自己“最硬”的维度。
  杂交体的归并判定：(C1) 对奶蛙侧成立（形体保留），
  (C2) 对两侧均成立（音轨、结尾套路 + PVZ 玩法），
  (C3) 由社群辨识实践确认（“植奶”“僵奶”成为独立词条）。
  杂交成功的前提是两个母体的形式层\keyw{不冲突}——
  这解释了为什么是塔防游戏而非叙事长篇
  更容易与短视频迷因杂交。
\end{nlcase}

\subsection{奶蛋：简化极限与“轮椅”叙事}

异形奶蛙中值得单独分析的是\term{奶蛋}：
奶蛙“简化到极致的产物”，在 B 站 UP 主“奶史\_”的
《选择你的奶伴》系列视频中高频出现，
与衍生语“安迪”结伴，因胜率极高而产生
“奶蛋轮椅”的说法（“安迪”空耳为“俺无敌”）
\conf{C}（第一版教材口径）。

奶蛋的学理价值在于它标定了谱系的\keyw{简化极限}：
当形体特征只剩“眼睛 + 大笑”，形象已经退化为一枚
近乎抽象的徽章——但只要它仍在既定形式
（音轨、结尾、衍生语）中出现，(C2) 就维持着簇归属。
“奶蛋轮椅”叙事则显示了内容与立场的耦合：
一个极简形象因为在特定叙事框架（对战、选择）
中“从无败绩”，被社群赋予了\keyw{强运}人设——
内容层的极度简化反而放大了立场层的投射空间。
这与表情包传播的一般规律一致：
形象越简单，可代入的自我越多。

\begin{nlcritique}[批判性反思：审丑与低质内容]
  奶蛙的“丑怪”不是中性事实。对这一谱系的批评集中在三点：
  其一，\textbf{审丑疲劳}——畸变与劣化作为刺激手段，
  在大量重复后会推高“出怪”的阈值，
  诱发更极端的形象变异，形成品质上的逐底竞争；
  其二，\textbf{低质内容的规模外部性}——《人民日报》报道
  明确把“用 AI 工具批量生成低质内容”列为官方反对的
  侵权形态\conf{A}\cite{jiang2026nailong}，
  而奶蛙谱系的生产方式恰好与“批量生成”高度重合
  （第八章详论）；
  其三，\textbf{对他人的刷屏外部性}——
  对不使用该迷因的人群，高密度的奶蛙内容
  构成信息流污染（第十章讨论校园场景中的这一问题）。
  本书把谱系当作符号系统分析，
  不等于为其中每一种用法背书。
\end{nlcritique}

\section{内容维度的正典侧：素材库与“糖”}\label{sec:06:canon}

奶蛙的内容维度不限于奶蛙自身：
正典奶龙的官方内容为二创提供了\keyw{素材库}。
这一节处理正典一侧的接受史（均为 \conf{C} 级，
据第一版教材记录）。

最受二创青睐的正典桥段是“我是奶龙，我才是奶龙”
（真假奶龙对峙）。据第一版教材的观察，
这一桥段受欢迎“并不是因为这一段好看，
而是因为这一段很糖”——
“糖”（社群用语）指内容幼稚、尬感强的特质
\conf{C}。社群甚至生产了讽刺性彩虹屁
（“我觉得真假奶龙这一段简直是智斗天花板……
我早已浑身湿透了”）来完成对“糖”的标记：
先模仿认真的夸奖语体，再让夸奖的内容
夸张到不可能被当真，从而实现反向评级。

围绕这一桥段形成了两条加工线：
其一，\keyw{倒放大法}——把音轨倒放，
“我是奶龙，我才是奶龙”变成
“我才是奶龙，欧蓝的手”之类的空耳文本
\conf{C}；其二，与特摄角色的\keyw{跨界绑定}
（如与“贝利亚”形象的并置），
使正典形象进入抽象语境。

\begin{nltermbox}{术语：空耳（soramimi）}
  源自日语“そらみみ”（空耳，幻听），
  指把一种语言的语音听成另一种语言
  （或本语言内其他词句）的再诠释实践。
  本书用法：\textbf{语音外壳保留、语义外壳被重新指派}
  的双层结构。在奶蛙语料中，
  空耳同时作用于正典素材（倒放产生的“欧蓝的手”）
  与衍生语（“安迪”$\leftrightarrow$“俺滴 / 俺无敌”）。
  空耳词等价项的判定标准与语音距离度量
  是第八章“可复现方法”的主题；
  此处只确立其符号学地位：空耳是\keyw{能指的劫持}，
  它让任何音轨都成为潜在的双关素材。
\end{nltermbox}

“糖”文化的存在解释了一个容易困惑的现象：
正典奶龙定位“治愈、欢乐”\conf{A}\cite{jiang2026nailong}，
其民间接受却是反讽式的——受众以“糖度”
为尺度重新分级官方内容，把最“幼儿向”的桥段
顶成传播峰值。接受（reception）与文本（text）
在这里显著脱钩：\keyw{民间评价体系
把正典当原料，而不当作品}。
这正是第二章“正典作为意义参照系”的具体运作方式——
社群以“糖”标记正典中“不可严肃对待”的部分，
为后续的畸变、拼贴清空了语义场地。
没有这场“清空”，奶蛙的畸变身体只会被视为
画坏了的奶龙；有了它，畸变才成为合法的风格。

\section{形式维度：叙事套路、音轨与衍生语}\label{sec:06:form}

\subsection{三件套：结构、音轨、语言}

奶蛙二创的形式层由三个高度稳定的构件组成。

第一，\keyw{叙事结构}：“压力桥段$\to$大笑收尾”。
典型视频以一个尴尬、失败、荒诞或高压的情境开场，
在情绪顶点处切入奶蛙的捧腹大笑，
以集体大笑为整个段落收尾\conf{C}\cite{shoutanji2026naiwa}。
这个结构让笑声承担“句号”的功能：
无论前文是什么，大笑都把它封存为“已经过去的事”。
第二，\keyw{音轨}：变调处理的爽朗笑声
（音源考据见第四章 \ref{sec:04:audio}），
二手评述还记录了其与变调版“不知火舞语音”的搭配
\conf{C}\cite{shoutanji2026naiwa}。
第三，\keyw{语言}：衍生语系统，
即下表的六个高频语音梗（均 \conf{C} 级，
据第一版教材整理）。

\begin{table}[htbp]\centering
  \caption{奶蛙衍生语系统（据第一版教材整理，\conf{C} 级）。
          “等效空耳”列即该词的再诠释空间，
          其判定标准第八章给出}
  \label{tab:06:langs}
  \small
  \begin{tabularx}{\textwidth}{@{}llX@{}}
    \toprule
    衍生语 & 主要运用场景 & 等效空耳（列举） \\ \midrule
    安迪 & 爬行奶蛙说话、奶蛋说话、奶蛙出场 &
      俺滴、安定、俺无敌、暗地 \\
    嘎嘎滴辣虾 & 运用场景广泛，高频说话效果 &
      刚刚点了谁、刚刚跌落水、刚刚掉了什 \\
    卡脚先 & 小剧场的转折或收尾 &
      刚救下、赶紧写、干啥先 \\
    无了 & 吃瘪、惊讶音效 & 哦了、好了 \\
    甘巴爹 & 肯定、疑问音效 & 干扒地、刚拔的 \\
    有用不 & 丰富语言，一般作“安迪”的下位替代 &
      忧郁菇、犹豫不 \\ \bottomrule
  \end{tabularx}
\end{table}

衍生语的使用有明确的\keyw{场景语法}：
“安迪”属于出场与移动（爬行到一半看向屏幕说“安迪”
是固定动作—语言组合），“无了”属于失败与惊愕，
“卡脚先”属于转折收束。用错场景不会被禁止，
但会被识别为外行——这正是立场维度
“圈内默示知识”的形式层证据。

\subsection{衍生动作：身体程式}

第一版教材记录的衍生动作有四种：\term{思考}、
\term{蹲起大笑}、\term{爬行}、\term{摆架势}
\conf{C}。其中爬行与“安迪”的绑定
（爬到一半看向屏幕、说出衍生语、继续爬行）
是动作—语言复合程式的典型；
摆架势则是“模仿奶蛙的青年们最常摆的动作”，
即程式从屏幕内迁移到屏幕外、
从形象动作迁移为使用者身体姿态的通道
（第十章的立牌合影、摆架势打卡是其空间化延伸）。

蹲起大笑另有一个“双人组合技”，
第一版教材以“画面过于神圣和有伤风化，
我们按下不表”为由未作展示\conf{C}。
本书沿用同一处理，并把它记为一个方法论样本：
\keyw{社群自身的语料边界实践}——
第一版教材的作者们已经在一个“什么不宜收录”的
问题上做了编辑决策。研究者的收录边界
（第十二章 E1--E3）与社群的自我边界不必一致，
但后者必须被记录，因为它是社群规范存在的直接证据。

\subsection{形式为什么稳定}

\begin{proposition}[形式稳定的三重来源（教学命题）]\label{prop:06:formstable}
  奶蛙形式层的稳定性（相较内容层的高替换率）
  可归因于三个互相加强的结构：
  (i) \textbf{学习成本}：形式构件数量少、组合规则简单，
  新用户一次接触即可习得；
  (ii) \textbf{识别成本}：固定音轨与结尾套路
  在信息流中提供即时可辨的“这是奶蛙视频”信号，
  满足平台分发对识别速度的要求；
  (iii) \textbf{组合开放}：内容层完全开放，
  任何素材都能接入同一形式，
  因此创新压力全部由内容层吸收，
  形式层无需变异即可持续产出“新”作品。
\end{proposition}

这一命题（\conf{M} 级）与第五章的
“变异—选择—保留”飞轮对接：形式层是被“保留”喂养的
惰性骨架，内容层是“变异”的活性表面。
第八章将从生产端给出同一结构的经济学版本：
AIGC 工具把内容变异的成本降到接近零，
而形式构件作为模板被工具固化，
于是“内容极活泼、形式极保守”成为均衡态。

\section{立场维度：用梗代替自述}\label{sec:06:stance}

\subsection{从表达到表态}

立场维度回答“用奶蛙在做什么”。本书归纳出
三层可观察的语用行为（\conf{M} 级归纳，
素材 \conf{C} 级）：

\begin{enumerate}[label=(\arabic*),leftmargin=2.4em]
  \item \textbf{自陈的替代}：用一张奶蛙梗图代替
        一句第一人称陈述，把“我很难受 / 我搞砸了 / 我绷不住了”
        翻译成“我们都这样”。戈夫曼的拟剧论
        把自我呈现分为前台与后台\cite{goffman1959presentation}；
        梗图在此充当\keyw{印象缓冲垫}——
        情绪信息被发出，但发出者退到了图像后面，
        既完成了表达，又规避了直接暴露的成本。
  \item \textbf{去抑制的通道}：网络环境的匿名性、
        不可见性与同步性缺失降低了表达门槛
        \cite{suler2004disinhibition}；
        “笑着崩溃”之所以可行，正因为笑声包裹了
        本该沉重的内容，使崩溃可以以娱乐的语体发出。
  \item \textbf{身份的默示声明}：正确使用衍生语、
        识别“过期”变体、掌握谱系细节，
        都是只有圈内成员能完成的操作。
        借用 Nissenbaum 与 Shifman 的框架，
        迷因使用是\keyw{可被社群审计的文化资本}
        \cite{nissenbaum2015contested}——
        转发什么、用什么语、对什么保持沉默，
        都在向同好发出可读取的信号。
        第七章的“逆反认证”机制（\ref{sec:07:reverse}）
        是这一逻辑在官方压制事件下的放大。
\end{enumerate}

\subsection{一份立场声明文本}

第一版教材的尾声收录了一篇社群信条式文本
（“奶蛙不叛”），其中罗列了奶蛙\keyw{不带来}的东西：
“没有勾心斗角”“没有拼死内卷，没有躺平摆烂”
“没有施加 PUA，没有拉升 BMI”“没有卡你 CET，
没有让你 EMO”\conf{C}（第一版教材原文节引）。

这份文本的符号学结构值得逐层拆解：
它的句式是排比否定（“没有 X”），
被否定的对象构成一份负面清单——内卷、绩点竞争、
身材焦虑、语言考试、情绪耗竭，
恰好是高校生活压力源的清单（第七章机制矩阵
“压力源”一列的民间版本）；
它把奶蛙定义为这些压力的\keyw{缺席者}，
并以此号召“发自内心地……笑出来吧”。
换言之，社群自己写下的立场声明
把“捧腹大笑”明确登记为\keyw{对压力情境的应答}，
而不是随机娱乐。这份文本是 \conf{C} 级材料，
不能当作“大学生使用奶蛙的动机”的证据；
但它至少证明：\keyw{在社群的自我理解中，
奶蛙与压力话语被组织进了同一个意义结构}。
动机是否果真如此，是第十一章问卷要回答的问题。

\begin{nlmethod}[研究方法提示：立场文本的取样]
  分析“社群立场”时，警惕以声明文本代替行为证据：
  信条式文本是\keyw{自我表述}，
  存在美化、跟风与戏仿三层混杂
  （“奶蛙不叛”本身是否也是玩梗？）。
  可行的取样纪律：
  (1) 声明文本只用于界定“社群自我理解”，
  不用于推断个体动机；
  (2) 行为证据（何时、对谁、在什么情境使用）
  优先于文本证据；
  (3) 对声明文本本身做语体标注——
  宗教戏仿、宣言戏仿的语体标记
  恰恰说明作者清楚这不是一份严肃文本。
\end{nlmethod}

\section{三维的不对称：“一软两硬”}\label{sec:06:asym}

本章的分析可以收拢为一个结构性判断：

\begin{table}[htbp]\centering
  \caption{奶蛙符号系统的“一软两硬”结构（本书归纳，\conf{M} 级）}
  \label{tab:06:asym}
  \small
  \begin{tabularx}{\textwidth}{@{}llX@{}}
    \toprule
    维度 & 状态 & 直接表现 \\ \midrule
    内容 & 软（高替换） &
      形象谱系横跨物种与器物；
      任何“黄色大笑生物”皆可接入；
      正典桥段可反讽征用 \\
    形式 & 硬（高稳定） &
      结尾集体大笑、变调音轨、
      衍生语场景语法、压力桥段结构长期不变 \\
    立场 & 硬（高约定） &
      “绷不住但用笑处理”的解读框架稳定；
      场景语法错误会被识别为外行；
      官方否认反而加固圈内立场 \\ \bottomrule
  \end{tabularx}
\end{table}

“一软两硬”解释了奶蛙作为符号系统的生存策略：
\keyw{用最硬的形式承载最软的内容}。
形式与立场的稳定保证了辨识与归属（(C2)(C3)），
内容的开放保证了变异供给（(C1) 的宽容度），
三者合起来正好是一个“可无限换壳的仪式框架”。
巴赫金笔下的狂欢节正是类似的结构：
仪式形式年年相同，插科打诨的内容随时更新
\cite{bakhtin1968rabelais}。
区别在于媒介条件——狂欢节一年一度，
奶蛙一天可以刷到几十次，
仪式框架的重复频率被平台分发放大了几个数量级
（第七章将讨论这种高频重复与情绪耗竭的关系）。

由此回到章首的承诺：为什么说奶蛙是“三个层面的叠加”
（第一章）而不仅是“一个梗”——
因为它的符号系统本身就分层运行：
正典素材库（内容原料）、
民间谱系（内容产品）、
形式—立场复合体（仪式框架）。
第九章讨论产业时将看到，
这个结构还决定了它的经济形态：
内容层无人收费，形式层无法收费，
最终被商业捕捉的是立场的周边化（立牌、玩偶）。

\section{本章小结}\label{sec:06:summary}

本章用 Shifman 三维框架解剖了奶蛙语料。
内容维度给出三级谱系（换皮、异形、客串）
与简化极限（奶蛋），并提出窄通道规律
（轮廓趋同、细节发散，\conf{M} 级）；
正典侧的“糖”文化与倒放空耳说明了素材库的
反讽式接受机制。形式维度拆解了
“结构—音轨—语言”三件套与四种衍生动作，
并给出形式稳定的三重来源命题。
立场维度区分了自陈替代、去抑制通道、
身份默示声明三种语用行为，
并以“奶蛙不叛”文本展示了社群的自我理解。
总的判断是“一软两硬”：
内容高替换、形式高稳定、立场高约定化。
遗留问题有二：
其一，形式层为什么在 AIGC 时代格外稳定，
需要生产端的证据（第八章）；
其二，立场层的三个语用行为
是否真的对应使用者的心理过程，
符号分析无法回答，需要受体研究（第十一章）。

\begin{nlexercises}
  \item 取表 \ref{tab:06:pedigree} 中任意三个变种，
        逐条检验归并标准 (C1)(C2)(C3)，
        并指出哪一条标准在判定中起了决定作用。
        若三个变种中有一个你认为应当\keyw{移出}谱系，
        写出你的理由（丢失了哪个不可再丢的特征）。
  \item “糖”是社群的负面—反讽评级词。
        再找一个你所知的网络社群评级词
        （如“雷”“尬”“典”），
        写出它的判定操作（什么可观察特征使一段内容获得该评级），
        并讨论：当评级词本身被官方或广告征用后，
        评级效力会发生什么变化？
  \item 本章把“奶蛙不叛”文本定位为
        “社群自我理解”而非“个体动机证据”。
        设计一个区分二者的最小研究方案：
        至少一个行为指标、一个文本指标，
        并说明两者结果不一致时你接受哪一个、为什么。
\end{nlexercises}

\nlkeywords 三维框架、形象谱系、换皮奶蛙、异形奶蛙、奶蛋、
空耳（soramimi）、糖、倒放大法、衍生语、衍生动作、
场景语法、一软两硬、窄通道规律

\chapter{捧腹大笑的机制分析：从五种解释到可检验假设}\label{ch:07}

\begin{nlkeypoints}
  \begin{enumerate}[label=(\arabic*),leftmargin=2.4em]
    \item 第一版教材的五重机制（情绪代偿、群体稀释、反精致、校园社交、
          逆反认证）被重写为\keyw{“压力源—补偿行为—可观察指标—
          竞争性解释”四列矩阵}，使其从感想变为可检验的分析框架
          （矩阵为本书构造，\conf{M} 级）。
    \item 五个机制各有学理来源（伯格林的机械性、弗洛伊德的能量节省、
          巴赫金的降格、普罗文的社交笑、
          Nissenbaum--Shifman 的文化资本），本章逐一给出
          \keyw{可检验假设与证伪条件}（H7.1--H7.5）。
    \item 机制之间不是并列关系而是\keyw{分层关系}：
          M1（代偿）与 M2（稀释）作用于个体—情绪层，
          M3（反精致）作用于风格层，M4（校园社交）作用于传播层，
          M5（逆反认证）作用于权威关系层；
          每层提出的指标不同，因此五个机制\keyw{可以被分别证伪}，
          而不是一荣俱荣。
  \end{enumerate}
\end{nlkeypoints}

本章处理全书的核心谜题（问题 Q2）：\keyw{为什么大学生会笑，
而且愿意反复转发}。第一章已把“为什么聪明人爱看无厘头”
改写为三种竞争读法，本章的任务是把其中功能—符号读法的
五个候选机制逐一学理化：先给每个机制找一个成熟的理论来源，
再把“机制”写成“压力源—补偿行为—可观察指标—竞争性解释”
的四列结构，最后为每个机制提出假设与证伪条件。
必须预先声明：本章没有任何经验数据，全部假设处于
\keyw{待检验}状态（第十一章给出检验工具）。
本章的分析对象——“捧腹大笑”这一表达形式——
其符号结构已在第六章给出；本章问的不是它是什么，
而是它\keyw{在做什么、为什么有人愿意反复做}。

\section{从“五重战神”到机制矩阵}\label{sec:07:matrix}

第一版教材第五部分以“捧腹大笑思想理论体系（迫真）”为题，
列出了五个机制并各给一段评述，末尾自嘲“（不行，扯不下去了）”。
这份材料的直觉质量很高——五个机制的\keyw{命名}在社群传播中
已经被反复引用——但它有三个方法论缺陷：

\begin{enumerate}[label=\arabic*.]
  \item \textbf{平行罗列}：五个机制之间是什么关系（互补？竞争？
        分层？），原文没有回答；
  \item \textbf{不可证伪}：每段评述都是“机制名 + 现象复述”，
        没有给出“如果该机制不成立，我们会看到什么”；
  \item \textbf{层次混淆}：情绪机制（代偿）与传播机制（校园社交）、
        权威关系机制（逆反认证）被放在同一层并列，
        导致“一个机制反例”会被误读为“整个体系崩塌”。
\end{enumerate}

本书的重写方案是把每个机制装进同一个四列结构：

\begin{table}[htbp]\centering
  \caption{五重机制的分析矩阵模板（本书构造的教学框架，\conf{M} 级）}
  \label{tab:07:matrixtpl}
  \small
  \begin{tabularx}{\textwidth}{@{}lX@{}}
    \toprule
    列 & 内容与要求 \\ \midrule
    压力源 & 机制所应答的现实处境；必须是可指认的情境
      （如“绩点竞争”），不得是“当代社会”这类不可操作的笼统表述 \\
    补偿行为 & 使用者实际做的事（转发、配文、立牌、摆架势）；
      必须与第六章立场维度的语用行为对应 \\
    可观察指标 & 单一数字或单一频率（如“压力考试周内的转发占比”）；
      必须原则上可计数 \\
    竞争性解释 & 若不引入该机制，哪些更简单的解释也能覆盖同一指标；
      机制只有在排除竞争解释后才有存在必要 \\ \bottomrule
  \end{tabularx}
\end{table}

四列的次序就是论证的次序：压力源立“需求”，
补偿行为立“供给”，可观察指标立“检验”，
竞争性解释立“对手”。任何一个机制若填不满四列
（尤其是填不出可观察指标），就不配称为机制，
只是修辞。后文各节按此模板逐个展开。

\section{笑的学理基础：四个古典来源}\label{sec:07:theory}

在进入机制之前，先把“笑”本身的理论资源铺开。
四个来源分别回答四个问题，后文机制将分别调用。

\subsection{优越论与机械论：笑作为社会矫正}

伯格森在《笑》中提出著名的公式：
\keyw{“镶嵌在生命之上的机械性”是可笑的}，
笑是“机械性对生命性的报复”——
当一个生命体表现得像机器（僵硬、重复、自动化）时，
旁观者的笑执行了一种社会矫正功能\cite{bergson1900rire}。
奶蛙语料中“僵直的爬行、重复的蹲起、程式化的摆架势”
（第六章衍生动作）正是“机械性”的具象化；
而 AI 生成的畸变身体（膨腹、细肢）则是
“生命形态的机械故障”的极端版本。
伯格森框架的预测：\keyw{笑点应随“机械化程度”单调增强}——
这是一个可操作化的变量（动作重复次数、程式化程度编码）。

\subsection{释放论：笑作为能量节省}

弗洛伊德在《诙谐及其与无意识的关系》中把笑解释为
\keyw{心理能量的节省}：诙谐绕过了压抑与批判性检查，
原本用于抑制的心理能量被省下、以笑的形式释放
\cite{freud1905witz}。该框架对奶蛙语料的对应物是
“绷不住”体验：紧绷（抑制）与绷不住（释放）的对立
恰好是能量模型的大众版。释放论的预测：
\keyw{禁忌或高压力语境中的笑应更强烈}——
第七章 M1（情绪代偿）直接调用这一预测。

\subsection{狂欢论：笑作为秩序倒置}

巴赫金把狂欢节式的笑描述为\keyw{降格}
（\en{degradation}）：一切崇高、精尖、正经之物
被拉向身体下部与物质底层，在笑声中完成暂时的等级翻转
\cite{bakhtin1968rabelais}。狂欢论的预测：
笑的对象应系统性地偏向\keyw{权威与精致符号}，
而非随机对象——第七章 M3（反精致）与
M5（逆反认证）分别调用这一预测的两半。

\subsection{社交论：笑作为关系信号}

普罗文的实验研究表明，笑远不只是“好笑”的情绪反应：
它是高度定型化的社交发声，绝大多数笑发生在社交情境中，
且与逗笑内容的质量弱相关；笑具有强传染性
\cite{provine1992laughter,provine2000laughter}。
社交论的预测：\keyw{同一素材的笑应随共同在场人数增加而增强}——
第七章 M4（校园社交）调用这一预测，
第十章的宿舍—教室场景为其检验场。

四个来源的分工可以压成一句话：
伯格森解释\keyw{笑什么}（机械性），
弗洛伊德解释\keyw{为什么释放}（能量节省），
巴赫金解释\keyw{冲着谁笑}（等级倒置），
普罗文解释\keyw{和谁一起笑}（社交信号）。
五个候选机制各自从中取用一角，
但都不与任何一个理论画等号——
理论提供的是\keyw{可检验预测的来源}，不是结论本身。

\begin{nlcritique}[批判性反思：理论的移植风险]
  四个古典理论各有其历史语境：伯格森谈的是沙龙与剧场的
  喜剧观众，弗洛伊德的材料是十九世纪末维也纳的诙谐，
  巴赫金的狂欢节是中世纪广场文化，普罗文的实验是
  北美大学生。把它们移植到“短视频平台的算法分发”
  语境，每一个都存在\keyw{前提错位}的可能——
  例如，算法推荐改变了“共同在场”的含义，
  普罗文的社交笑结论未必平移成立。
  因此本章的使用方式是\keyw{把理论当假设生成器}，
  而非当证据：每个理论只贡献一条可检验预测，
  预测不成立时被淘汰的是“理论在奶蛙语境的应用”，
  不必连理论本身一起否定。
\end{nlcritique}

\section{M1 情绪代偿：把“绷不住”合法化}\label{sec:07:m1}

\subsection{机制表述}

第一版教材的原文：“考研、绩点、就业的复合压力下，
负面情绪缺少体面的出口。奶蛙提供了一种‘笑着崩溃’的姿态
——既表达了失控，又不至于沉重。”
按四列矩阵重写：

\begin{table}[htbp]\centering
  \caption{机制 M1（情绪代偿）的四列展开（\conf{M} 级）}
  \label{tab:07:m1}
  \small
  \begin{tabularx}{\textwidth}{@{}lX@{}}
    \toprule
    压力源 & 绩点竞争、考研备考、求职季等可指认的高压情境
      （以学期周期为节律） \\
    补偿行为 & 在压力情境的社媒表达中使用奶蛙大笑素材，
      以“笑”的语体发出原本沉重的自述 \\
    可观察指标 & 压力事件（期中考、考研倒计时节点）前后，
      奶蛙类内容使用频率的\keyw{相对变化率}；
      与一般娱乐内容使用频率变化率之差 \\
    竞争性解释 & (a) 纯消遣：高压周里人们本来就刷更多手机，
      奶蛙只是“恰好刷到”；(b) 时间填塞：备考期社交时间减少，
      低成本内容占比被动上升 \\ \bottomrule
  \end{tabularx}
\end{table}

\subsection{学理对应与假设}

情绪代偿的学理骨架是弗洛伊德的释放论
（\ref{sec:07:theory} 小节）加上情绪的“行动倾向”理论：
Frijda 主张情绪内置行动准备（emotional action tendency），
表达行为是行动倾向的出口\cite{frijda1986emotions}。
“笑着崩溃”在这个框架里读作：
崩溃的行动倾向（退出、哭、求助）被“笑”这一
社交许可的出口替代，且替代不损失信息量——
接收方仍能读出崩溃（第六章立场维度的“自陈替代”）。

\begin{hypothesis}[H7.1 情绪代偿假设]\label{hyp:07:m1}
  在高校用户的社媒使用中，奶蛙类内容的使用频率
  随压力情境的出现而\keyw{选择性上升}：
  上升幅度显著高于其总娱乐内容使用频率的上升幅度，
  且该差值在自评压力更高的用户中更大。
\end{hypothesis}

\keyw{证伪条件}：若压力期与非压力期的奶蛙使用频率比，
不高于总娱乐内容的两期频率比（即无选择性），
或者选择性上升与自评压力无关，则 H7.1 被证伪，
情绪代偿退化为“时间填塞”的被动效应。
\keyw{必须声明}：H7.1 尚无任何经验检验，
第十一章的问卷设计（量表 + 学期节律对齐）是为其准备的。
“笑着崩溃”的表达门槛问题还与在线去抑制效应相关：
匿名与不可见性降低情绪表达成本
\cite{suler2004disinhibition}——这一因素在
H7.1 中被当作常量背景，不作单独检验。

\section{M2 群体稀释：用“我们”消解“我”}\label{sec:07:m2}

\subsection{机制表述}

第一版教材原文：“用一张梗图代替自述，把‘我很差劲’
翻译成‘我们都这样’。个体暴露被群体化收编，
自嘲因而不必付出羞耻成本。”重写：

\begin{table}[htbp]\centering
  \caption{机制 M2（群体稀释）的四列展开（\conf{M} 级）}
  \label{tab:07:m2}
  \small
  \begin{tabularx}{\textwidth}{@{}lX@{}}
    \toprule
    压力源 & 自我暴露的羞耻成本：直接陈述失败或焦虑
      需要承担社交评价风险 \\
    补偿行为 & 以梗图、群发、评论区复读等方式发布，
      使个人陈述获得“大家都这样”的群体外观 \\
    可观察指标 & 第一人称自述句与奶蛙梗图配对出现的比例；
      同一内容的“带梗图 / 不带梗图”两版获得的
      自我暴露类回复数之差 \\
    竞争性解释 & (a) 表达习惯：梗图只是流行文体，与羞耻无关；
      (b) 效率：图像比文字输入更快，用户选图仅为省事 \\
    \bottomrule
  \end{tabularx}
\end{table}

\subsection{学理对应与假设}

M2 的学理骨架是戈夫曼的拟剧论与情绪感染的合流。
戈夫曼的\keyw{前台/后台}区分指出，个体在前台
管理自己给他人留下的印象\cite{goffman1959presentation}；
梗图正是前台管理装置——情绪信息（后台事实）
经由一个“非我”的图像发出，发送者保持前台体面。
情绪感染研究则说明接收端如何“接住”这一信号：
情绪在群体中经由模仿与反馈循环传播
\cite{hatfield1994contagion}，数字环境放大了
发送者—媒介—接收者三要素中的媒介作用
\cite{goldenberg2020digital}。“我们都这样”之所以能成立，
正因为梗图天然携带\keyw{可复读性}——
它邀请接收者以转发代替安慰，
把一次情绪事件转写为一次集体复读。

\begin{hypothesis}[H7.2 群体稀释假设]\label{hyp:07:m2}
  同一负面自述内容，附带奶蛙类梗图时：
  (a) 发布者事后删除率更低（羞耻成本更小的行为学代理指标）；
  (b) 收到的回复中“同感/复读类”占比更高、“安慰/追问类”占比更低。
\end{hypothesis}

\keyw{证伪条件}：若带梗图与否不影响上述两项指标，
则“稀释”只是文体偏好，不承担心理功能。
M2 与第五章“文化资本”框架的连接点是：
用梗代替自述同时完成\keyw{情绪表达}与
\keyw{身份认证}两个动作\cite{nissenbaum2015contested}——
但必须指出，身份认证也能解释同样的行为
（用梗因为“圈内人应该用梗”），这正是 M2 的
竞争性解释之一，判别设计见 \ref{sec:07:compete} 节。

\section{M3 反精致：对优秀人设的审美罢工}\label{sec:07:m3}

\subsection{机制表述}

第一版教材原文：“大学生被期待为‘精致的、上进的、
有规划的人’。转发一只丑、怪、抽象的奶蛙，
是一次低成本的不合作表态。”重写：

\begin{table}[htbp]\centering
  \caption{机制 M3（反精致）的四列展开（\conf{M} 级）}
  \label{tab:07:m3}
  \small
  \begin{tabularx}{\textwidth}{@{}lX@{}}
    \toprule
    压力源 & “精致人设”的维持成本：朋友圈、社交主页
      中的上进形象要求与其实际状态的落差 \\
    补偿行为 & 主动转发畸变的、劣化的、反审美的形象，
      以展示对精致规范的不合作 \\
    可观察指标 & 用户内容流中“故意劣化”类内容占比；
      同一用户在成绩/求职等评价敏感期的劣化内容占比变化 \\
    竞争性解释 & (a) 审美多元：劣化风格本身即审美偏好，
      无表态含义；(b) 模仿流行：跟随社群风格而非对抗规范 \\
    \bottomrule
  \end{tabularx}
\end{table}

\subsection{学理对应与假设}

M3 直接调用巴赫金的降格概念
（\ref{sec:07:theory} 小节）：把“精致”拉入
“丑怪—劣化—荒诞”的物质层面，是狂欢式倒置的当代版。
伯格森的机械论也在此处工作：
奶蛙的“机械性身体”（重复动作、畸变比例）
与精英人设所要求的“得体的生命管理”形成对照，
笑它即是宣告不参与该规范\cite{bergson1900rire}。
亚文化研究为“不合作表态”提供了更精确的表述：
Hebdige 指出亚文化风格是通过\keyw{符号层面的抵抗}
组织起来的，而抵抗终将被商业与主流收编
\cite{hebdige1979subculture}——后者正是第九章
（奶蛙周边商业化）要展开的“收编”环节。

\begin{hypothesis}[H7.3 反精致假设]\label{hyp:07:m3}
  在评价敏感期（绩点公布、求职结果发布），
  用户发布“故意劣化”内容的频率相对基线的上升幅度，
  显著高于非敏感娱乐内容的同期上升幅度；
  且该上升在自我呈现压力更高（如以学术/精英形象为主页）的
  用户中更强。
\end{hypothesis}

\keyw{证伪条件}：劣化内容使用在敏感期与非敏感期无差异，
或上升与主页人设类型无关。M3 的一个独特困难是
“故意劣化”的编码判定——需要区分\keyw{真的做不好}
与\keyw{故意做坏}，编码手册必须给出
（如“刻意使用低分辨率素材、故意错位排版、
拒绝滤镜修正”等可观察标记），这一编码方案
在第十一章给出。

\begin{nlcritique}[批判性反思：“反精致”的阶级自嘲问题]
  “反精致”的表述有一个必须直面的张力：
  第一版教材的读者群（顶尖高校学生）恰好是
  “精致竞争”的\keyw{优胜者或准优胜者}。
  “优胜者扮演反规范”与传统亚文化中
  “弱势者表达抵抗”的结构并不相同——
  前者的“罢工”没有实际代价，甚至可能成为
  新一轮身份竞赛的筹码（“我反内卷”本身成为人设）。
  这提示 M3 可能只是\keyw{符号消费}而非\keyw{规范抵抗}：
  判别依据是“反精致表态是否伴随实际行为让步”
  （如真的退出绩点竞争）。若无伴随行为，
  M3 应降级为“人设的一部分”。
  本书不预设哪一种成立，但要求研究者报告判别证据。
\end{nlcritique}

\section{M4 校园社交：复读机与共同在场}\label{sec:07:m4}

\subsection{机制表述}

第一版教材原文：“集体居住让奶蛙拥有‘线上’‘面对面’
双通道传播途径。同一个梗在同一屋檐下被反复复诵，
趣味衰减极慢。”重写：

\begin{table}[htbp]\centering
  \caption{机制 M4（校园社交）的四列展开（\conf{M} 级）}
  \label{tab:07:m4}
  \small
  \begin{tabularx}{\textwidth}{@{}lX@{}}
    \toprule
    压力源 & 社交需求：群属感的建立与维持需要
      低成本的共同话题 \\
    补偿行为 & 群聊复读、线下引用（口头“安迪”）、
      表情包互发、立牌合影打卡 \\
    可观察指标 & 同一梗在单一群聊中的复读次数分布；
      “线上引用$\to$线下口头使用”的转化率；
      同好会纳新期的话题集中度 \\
    竞争性解释 & (a) 算法同质化：同一校园用户群本就
      被推送相似内容，复读是算法而非社交的产物；
      (b) 人口学同质：同龄同校用户的偏好本来就相近 \\
    \bottomrule
  \end{tabularx}
\end{table}

\subsection{学理对应与假设}

M4 的学理骨架是普罗文的社交笑理论
（\ref{sec:07:theory} 小节）：笑首先是一种
社交纽带行为，其发生频率取决于社交在场密度
\cite{provine2000laughter}。校园环境（集体宿舍、
固定班级、同好会组织）恰好是高社交在场密度的空间，
“线上梗 + 线下复读”的双通道因此成为可能。
第十章将展开传播网络的结构分析（弱连接与强连接
\cite{granovetter1973weak}、创新扩散的采用者类别
\cite{rogers2003diffusion}）；本章只提出
与 M4 直接相关的行为假设。

\begin{hypothesis}[H7.4 校园社交假设]\label{hyp:07:m4}
  同一奶蛙素材的“趣味衰减速度”与使用者所在群体的
  线下共同在场密度负相关：
  集体居住（宿舍）用户对该素材的持续使用时长
  显著长于非集体居住用户，且复读行为可预测
  群体凝聚力的自评指标。
\end{hypothesis}

\keyw{证伪条件}：集体与非集体居住用户的素材衰减曲线
无显著差异；或复读行为与凝聚力指标不相关。
第五章 SIR 模型的 $\gamma$（退出传播率）在 M4 框架下
获得一个机制解读：共同在场密度通过“复读强化”
降低 $\gamma$——这是把 M 级模型参数与
社会学变量挂钩的一个示例，仍然是教学性的
（\conf{M} 级），其价值在于指明该测量哪个量。

\section{M5 逆反认证：官方压制作为背书}\label{sec:07:reverse}

\subsection{机制表述}

第一版教材原文：“2026 年 3 月，奶龙官方关评、拉黑，
反而被网友读作‘官方破防’。权威的否认，
成了 meme 含金量的最高背书，热度此后一发不可收拾。”
重写：

\begin{table}[htbp]\centering
  \caption{机制 M5（逆反认证）的四列展开（\conf{M} 级）}
  \label{tab:07:m5}
  \small
  \begin{tabularx}{\textwidth}{@{}lX@{}}
    \toprule
    压力源 & 迷因社群的地位竞争：在梗快速过期的环境里，
      需要持续的“含金量”信号来区分圈内圈外 \\
    补偿行为 & 转发被官方否认的内容、以“官方破防”为题二创、
      在被拉黑后展示“阵亡”截图 \\
    可观察指标 & 官方否认事件（2026-03 关评拉黑）前后，
      奶蛙相关内容的日新增量曲线形状；
      “破防”类标题内容的占比变化 \\
    竞争性解释 & (a) 蒸汽效应（禁果效应）：单纯的好奇心驱动，
      与地位竞争无关；(b) 无关事件叠加：2026 年视频生成工具
      上线等供给侧事件才是主因（见第八章） \\
    \bottomrule
  \end{tabularx}
\end{table}

\subsection{学理对应与假设}

M5 的学理骨架是 Nissenbaum 与 Shifman 的
“迷因作为争夺中的文化资本”\cite{nissenbaum2015contested}：
在梗社群中，懂得并使用“正确”的梗是地位来源，
而\keyw{外部权威的敌意承认}是含金量最便宜的证据
——它证明该梗重要到值得被压制。巴赫金的狂欢论
提供了互补视角：权威与民间在狂欢中的等级翻转，
“官方破防”叙事本身就是一次降格表演
\cite{bakhtin1968rabelais}。

第四章已经记录了事实层面的材料
（\ref{sec:04:official}）：官方行为是
“呼吁停止丑化二创、关闭评论区、批量拉黑”，
且“直至今日你依然可以在奶龙官方视频的评论区
看见无数奶蛙图片”（\conf{C} 级）。
第五章的命题 \ref{prop:05:second}（二次爆发的条件）
则为 M5 提供了时间结构：重新激活需要
“未被耗尽的易感存量 + 足够强的重新激活项”，
官方事件恰是一个强激活项的候选。

\begin{hypothesis}[H7.5 逆反认证假设]\label{hyp:07:m5}
  2026-03 官方否认事件之后，奶蛙相关内容的日新增量
  出现\keyw{不能由供给侧事件单独解释的加速}；
  且加速期内容中，“官方破防”类框架的占比
  显著高于事件前基线。
\end{hypothesis}

\keyw{证伪条件}：控制视频生成工具上线等供给侧变量后
（第八章给出该变量的时间线），加速消失；
或加速存在但其内容中“官方破防”框架占比无变化
（说明热度上升与“否认—背书”叙事无关）。
必须强调：H7.5 的检验依赖 \conf{C} 级事件时间线
与平台数据的可得性，在两者齐备之前，
M5 的地位是\keyw{与史实相容的解释}，
不是已获支持的解释。

\begin{nlcritique}[批判性反思：逆反认证的伦理面]
  “官方压制$\to$含金量上升”若为真，它提供了一个
  令人不适的推论：\keyw{压制是迷因的养料}。
  这对权利人意味着一个两难——回应则认证，不回应则扩散；
  对研究者则意味着：M5 的验证不应变成
  “如何更好地玩逆反”的操作手册。同时必须记住
  第五章的提醒：资本框架不能取消“单纯觉得好笑”
  的朴素解释；官方否认事件在普通转发者那里
  很可能根本不构成“地位信号”，
  只是“又刷到了”。M5 预测的是
  \keyw{内容生产端}（二创作者）的框架占比变化，
  不是\keyw{消费端}的动机变化——
  把两者混为一谈是逆反认证叙事最常见的滥用。
\end{nlcritique}

\section{假设汇总与判别设计}\label{sec:07:compete}

\subsection{五个假设的一览}

\begin{table}[htbp]\centering
  \caption{五个机制假设的汇总（全部处于待检验状态；
          检验工具见第十一章）}
  \label{tab:07:hyps}
  \small
  \begin{tabularx}{\textwidth}{@{}llXX@{}}
    \toprule
    假设 & 机制 & 核心预测 & 证伪条件（简写） \\ \midrule
    H7.1 & 情绪代偿 & 压力期使用选择性上升 &
      无选择性或与自评压力无关 \\
    H7.2 & 群体稀释 & 带梗图降低删除率、
      改变回复结构 & 两指标均无差异 \\
    H7.3 & 反精致 & 敏感期劣化内容占比上升 &
      无周期性或与人设无关 \\
    H7.4 & 校园社交 & 在场密度延长素材寿命 &
      衰减曲线无居住类型差异 \\
    H7.5 & 逆反认证 & 否认事件后加速且
      “破防”框架占比上升 & 控制供给侧后加速消失 \\
    \bottomrule
  \end{tabularx}
\end{table}

\subsection{机制不是并列的：分层模型}

五个机制作用于不同层次，因此它们之间的关系
不是“五选一”，而是\keyw{分层叠加、各自可证伪}：

\begin{figure}[htbp]\centering
\begin{tikzpicture}[font=\small,
  layer/.style={rounded corners,draw=NLDeep,thick,align=center,
    minimum height=1.15cm,minimum width=4.0cm},
  lab/.style={font=\footnotesize\heiti,color=NLDeep},
  arr/.style={-{Stealth[length=2.6mm]},thick,NLDeep}]
  \node[layer,fill=NLRed!8]   (m5) at (0,0)   {M5 逆反认证\\{\footnotesize 权威关系层}};
  \node[layer,fill=NLGreen!8] (m4) at (0,-1.7){M4 校园社交\\{\footnotesize 传播层}};
  \node[layer,fill=NLViolet!8](m3) at (0,-3.4){M3 反精致\\{\footnotesize 风格层}};
  \node[layer,fill=NLAmber!12](m2) at (0,-5.1){M2 群体稀释\\{\footnotesize 个体—社交层}};
  \node[layer,fill=NLYellow!18](m1) at (0,-6.8){M1 情绪代偿\\{\footnotesize 个体—情绪层}};
  \draw[arr] (m1) -- node[right,font=\footnotesize]{情绪表达需要形式} (m2);
  \draw[arr] (m2) -- node[right,font=\footnotesize]{表达固化成风格} (m3);
  \draw[arr] (m3) -- node[right,font=\footnotesize]{风格需要传播通道} (m4);
  \draw[arr] (m4) -- node[right,font=\footnotesize]{传播引来权威反应} (m5);
  \node[lab,anchor=west] at (2.6,-3.4) {自下而上：需求逐层装配};
  \draw[arr,dashed,NLGrey] (m5.west) to[out=200,in=160] node[left,font=\footnotesize]{反馈：压制反哺下层} (m1.west);
\end{tikzpicture}
  \caption{五机制的分层装配模型（本书构造的教学框架，\conf{M} 级；
          虚线为 M5 对下层的反馈回路）。}
  \label{fig:07:layers}
\end{figure}

分层的论证意义有二。第一，\keyw{证伪是局部的}：
H7.1 被证伪不连坐 H7.3——这与第一版教材
“五重并列”的印象不同。第二，\keyw{解释有顺序}：
底层机制（M1）是发生学的起点，
顶层机制（M5）是生态的放大器；
若研究只观察到 M5（最显眼、最有戏剧性），
而跳过 M1--M4，就会把放大器误当发动机——
这是多数“官方破防”评论的结构性错误。

\subsection{竞争性解释的判别设计}

第一章保留的三种读法中，纯消遣读法与
技术—平台读法在本章换装为各机制的竞争性解释。
判别它们与功能机制的核心手法是\keyw{找不对称指标}：

\begin{enumerate}[label=\arabic*.]
  \item \textbf{周期不对称}：功能机制预测使用频率
        随压力周期波动（学期节律）；
        纯消遣与算法机制预测平稳或仅随平台活动波动。
        指标：使用频率的学期自相关。
  \item \textbf{对象不对称}：功能机制预测使用集中于
        情绪相关语境（自述、吐槽、成绩话题）；
        纯消遣预测均匀分布于一般娱乐语境。
        指标：使用语境的编码分布。
  \item \textbf{群体不对称}：功能机制预测高压力群体
        （毕业年级、考研群体）使用强度更高；
        算法机制预测使用强度仅随平台曝光量变化，
        与压力群体身份无关。
        指标：分群体的强度比较（控制曝光量）。
\end{enumerate}

三条不对称指标同时成立时，功能—符号读法
才获得超过竞争读法的证据权重；
任何一条不成立，都须相应调低对应机制的置信度。
这是“本书立场是要检验的”这句话的操作化版本
（第一章 \ref{sec:01:puzzle}）。

\section{本章小结}\label{sec:07:summary}

本章把第一版教材的五重机制重写为四列矩阵
（压力源—补偿行为—可观察指标—竞争性解释），
为每个机制补齐了学理来源（伯格森、弗洛伊德、
巴赫金、普罗文、Frijda、戈夫曼、情绪感染研究、
Hebdige、Nissenbaum--Shifman）与可检验假设
（H7.1--H7.5），并给出证伪条件。
结构上，五机制被组织为自下而上的分层装配：
情绪层（M1）$\to$社交层（M2）$\to$风格层（M3）
$\to$传播层（M4）$\to$权威关系层（M5），
M5 同时以反馈回路作用于下层。
必须重申全书纪律：本章没有任何经验数据，
五个假设全部待检验；“笑着崩溃”“我们都这样”
“官方破防”这些生动的社群表述，
在检验完成之前都只是假设的通俗名。
遗留问题有二：机制要运转，需要低成本的生产工具——
AIGC 如何把 M1--M4 的表达成本压到趋近于零，
是第八章的任务；机制要落地，需要真实的受访者——
第十一章给出 H7.1--H7.3 的问卷与 H7.4 的
行为测量设计。

\begin{nlexercises}
  \item 任选一个你熟悉的其他网络迷因，
        用表 \ref{tab:07:matrixtpl} 的四列模板为其写一个机制假设，
        并给出证伪条件。写完后自查：
        你的“可观察指标”是否能在不接触当事人隐私的前提下计数？
        若不能，请改写指标。
  \item H7.2（群体稀释）的竞争性解释之一是
        “身份认证”（用梗因为圈内人应该用梗）。
        设计一个判别实验（真实或自然实验均可），
        说明在什么结果下两个解释可以分开。
  \item 图 \ref{fig:07:layers} 的分层模型预测：
        若 M1（情绪代偿）被证伪，M3（反精致）仍可独立成立。
        请说明该预测成立与否分别对
        “奶蛙现象的可持续性”意味着什么
        （提示：情绪压力是周期性的，风格规范是长期的）。
  \item “逆反认证”若被证实，权利人将面临
        “回应即认证、不回应即扩散”的两难。
        从迷因治理的角度写 300 字分析：
        除“回应/不回应”外还存在哪些策略选项，
        各自对应本章哪个机制的变动。
\end{nlexercises}

\nlkeywords 机制矩阵、情绪代偿、群体稀释、反精致、校园社交、
逆反认证、证伪条件、竞争性解释、机械性、能量节省、降格、
社交笑、分层装配模型

\chapter{AIGC 与迷因生产链：变异引擎的成本革命}\label{ch:08}

\begin{nlkeypoints}
  \begin{enumerate}[label=(\arabic*),leftmargin=2.4em]
    \item 奶蛙二创的生产链可拆为五环节（素材获取、形象生成、
          视频生成、音轨处理、分发），其中前四环节在 2024—2026 年间
          被 AIGC 工具相继改造；\keyw{变异引擎}是形象的 AI 重绘与
          视频生成模型（据二手评述，抖音 \tsm{Seedance} 于 2026 年
          上线后扩散显著加速，\conf{C} 级）。
    \item AIGC 给了第六章“一软两硬”结构一个经济学解释：
          \keyw{内容变异的边际成本趋近于零，形式构件作为模板
          被工具固化}，均衡态因此是“内容极活泼、形式极保守”
          （本书分析，\conf{M} 级）。
    \item 本章同时给出全书第一个\keyw{可复现方法}的完整样例：
          空耳词等价项的三条判定标准与“拼音声韵母编辑距离”
          排序指标（\conf{M} 级，仅为排序用），
          并以六个衍生语的空耳变体作演示计算。
  \end{enumerate}
\end{nlkeypoints}

前三章（符号、机制）都留下过同一类承诺：
“为什么奶蛙在 2026 年才爆发”“为什么形式稳定而内容狂变”，
答案都指向生产端。本章的任务是把生产端摊开：
谁在生产、用什么工具、每个环节的成本发生了什么变化、
以及——作为全书方法论承诺的兑现——其中一类语言材料
（空耳词）的判定如何做成\keyw{可复现的操作}。
本章的宏观背景口径（短视频用户规模等）以中国互联网络
信息中心的历次统计报告为准\cite{baidu-index}，
本书不引用具体数值，读者引用前应核对当期报告。

\section{生产链拆解：五个环节}\label{sec:08:chain}

把一条典型奶蛙二创视频倒放一遍，可以拆出五个生产环节。
表 \ref{tab:08:chain} 给出各环节的传统做法、AIGC 改造与
成本变化方向（拆解为本书分析，\conf{M} 级；
工具描述的来源均为 \conf{C} 级二手评述
\cite{shoutanji2026naiwa} 与社群考据
\cite{lostmedia2026naiwa}）。

\begin{table}[htbp]\centering
  \caption{奶蛙二创生产链的五个环节（本书拆解，\conf{M} 级）}
  \label{tab:08:chain}
  \small
  \begin{tabularx}{\textwidth}{@{}llXX@{}}
    \toprule
    环节 & 传统做法 & AIGC 改造 & 成本变化 \\ \midrule
    1 素材获取 & 手工截图、抠图 &
      智能体平台生成角色形象（如星野 AI 的头像生成） &
      下降：从“找图”变为“生成图” \\
    2 形象生成 & 手工改图（换皮、P 图） &
      图像模型按提示词批量出变体 &
      大幅下降：变异供给近乎无限 \\
    3 视频生成 & 逐帧动画、剪辑拼接 &
      视频生成模型直接出动态片段 &
      大幅下降：静态形象获得“演技” \\
    4 音轨处理 & 手动变调、对轨 &
      变调已成模板化操作 &
      下降后企稳：音轨本身标准化 \\
    5 分发 & 论坛、群聊 &
      短视频平台算法推荐 &
      结构变化：分发不再是瓶颈 \\
    \bottomrule
  \end{tabularx}
\end{table}

\begin{figure}[htbp]\centering
\begin{tikzpicture}[font=\small,
  node distance=0.45cm,
  st/.style={rounded corners,draw=NLDeep,thick,fill=NLYellow!18,
    align=center,minimum height=1.0cm,minimum width=2.35cm},
  tool/.style={font=\footnotesize,color=NLAmber,align=center},
  arr/.style={-{Stealth[length=2.6mm]},thick,NLDeep}]
  \node[st] (s1) {素材获取};
  \node[st,right=of s1] (s2) {形象生成};
  \node[st,right=of s2] (s3) {视频生成};
  \node[st,right=of s3] (s4) {音轨处理};
  \node[st,right=of s4] (s5) {分发};
  \node[tool,below=0.30cm of s1] {智能体平台\\头像生成};
  \node[tool,below=0.30cm of s2] {图像模型\\批量变体};
  \node[tool,below=0.30cm of s3] {视频生成模型\\（Seedance 等）};
  \node[tool,below=0.30cm of s4] {变调模板};
  \node[tool,below=0.30cm of s5] {算法推荐};
  \draw[arr] (s1) -- (s2);
  \draw[arr] (s2) -- (s3);
  \draw[arr] (s3) -- (s4);
  \draw[arr] (s4) -- (s5);
  \draw[arr,dashed,NLAmber] (s5.south) to[out=250,in=290]
    node[below,font=\footnotesize]{反馈：完播 / 转发数据反哺变异方向} (s2.south);
\end{tikzpicture}
  \caption{奶蛙二创生产链与 AIGC 工具的接入点（\conf{M} 级拆解；
          虚线为平台数据对上游变异方向的反馈）。}
  \label{fig:08:chain}
\end{figure}

\begin{nlhistory}[史料与考据：变异引擎的两个时点]
  \textbf{2024-10-02}：星野 AI 用户“纲手【百豪】”创建“奶龙”智能体，
  平台对用户输入图片的 AI 重绘产生了最初奶蛙形象——
  生产链环节 1 的意外产物成为整个迷因的起点
  （考据，\conf{C}\cite{lostmedia2026naiwa}；见第四章）。
  \textbf{2026 年上半年}：抖音视频生成模型 \tsm{Seedance} 推出，
  奶蛙视频的生产门槛显著降低，被二手评述列为
  奶蛙爆发期的供给侧变量
  （\conf{C}\cite{shoutanji2026naiwa}）。
  两个时点之间的 2024—2025 年，恰是第四章所述的
  “低可见度潜伏期”——\keyw{形象已变异、动态未解锁}，
  与第五章命题 \ref{prop:05:second}（二次爆发需要
  足够强的重新激活项）在时间线上相容。
\end{nlhistory}

必须声明这条时间线的证据地位：两个时点均为 \conf{C} 级，
且“Seedance 促成爆发”只是评述者的归因，
不是受控结论。它在本书中的用法是
\keyw{为 H7.5 的检验提供须控制的供给侧变量清单}
（第七章 \ref{sec:07:reverse} 的证伪条件之一），
而非作为已证事实参与论证。

\section{成本结构：形式稳定性的经济学版本}\label{sec:08:cost}

第六章命题 \ref{prop:06:formstable} 给出了形式稳定的
三重来源（学习成本、识别成本、组合开放）。
生产端视角可以把它压缩成一句更锋利的话：

\begin{proposition}[成本不对称命题（教学命题）]\label{prop:08:asym}
  当图像与视频生成把\keyw{内容变异}的边际成本压到趋近于零，
  而平台分发以\keyw{识别速度}为筛选条件时，
  稳定的形式构件（固定音轨、结尾套路、衍生语）
  成为唯一不可替代的资产：
  内容层发生“无限供给下的激烈淘汰”，
  形式层发生“模板固化”——均衡态即
  “内容极活泼、形式极保守”。
\end{proposition}

这条命题（\conf{M} 级）给出两个可检验推论。
推论一：AIGC 工具普及后，单条奶蛙视频的
\keyw{内容新颖度}上升而\keyw{形式新颖度}不升——
可通过内容编码（场景、题材的多样性指数）
与形式编码（音轨、结构的多样性指数）的对比检验。
推论二：形式构件的变异失败率应高于内容变异
——“改结尾套路”的二创传播表现应系统性差于
“同结构换内容”的二创，因为前者破坏了识别信号。
两个推论的编码方案第十一章一并给出。

反过来，这条命题也解释了第六章“窄通道规律”
（命题 \ref{prop:06:channel}）的一个竞争性解释：
轮廓趋同可能不是社群选择的结果，而是
\keyw{工具模板}的产物——图像模型对同一提示词的
输出分布本来就集中。判别方法在第六章已给出
（工具固定/社群固定的对照样本），此处只需指出：
生产端证据越强，窄通道规律归因于社群选择的
置信度就要相应调低。这是竞争性解释之间
\keyw{此消彼长}的典型形态。

\begin{nlcritique}[批判性反思：低质内容的规模问题]
  生产成本趋零的另一面是\keyw{供给的泛滥}。
  《人民日报》报道明确将“用 AI 工具批量生成低质内容”
  列为奶龙 IP 遭遇的侵权形态之一
  \conf{A}\cite{jiang2026nailong}——这是全书与
  AIGC 直接相关的最强的一条 \conf{A} 级材料，
  它说明“批量生成”不是研究者的抽象担忧，
  而是权利人实际维权清单上的项目。
  从传播政治经济学看，这里的结构性问题是：
  \keyw{生产成本由个体承担（趋零），拥挤成本由
  公共信息流承担（趋增）}——每个二创作者的最优策略
  （多产、快产）加总起来是全体用户的信息流污染。
  第六章的审丑批判、第十章的校园刷屏争议，
  都可以视为这一外部性的不同侧面。
  本书不评判任何具体二创行为的合法性（第十二章），
  但要求读者带着这个外部性结构去理解
  官方、平台与社群三方各自的反应。
\end{nlcritique}

\section{可复现方法：空耳词的判定与度量}\label{sec:08:soramimi}

第六章表 \ref{tab:06:langs} 列出了六个衍生语及其
“等效空耳”。本节兑现方法论承诺：把“等效”
变成可执行的判定与可复现的度量。

\subsection{三条判定标准}

沿用事实手册的口径，一个候选空耳变体 $w'$ 与
衍生语 $w$ 建立\keyw{等价关系}，须同时满足：

\begin{enumerate}[label=(\arabic*),leftmargin=2.4em]
  \item \textbf{语音近似}：音节数一致或仅差一音节；
        对应音节的声母、韵母大体对应（细则见下）；
  \item \textbf{场景稳定}：$w'$ 在社群使用中与 $w$
        出现在相同或明确对应的场景
        （如“俺无敌”与“安迪”都出现于奶蛋的“高光”场景）；
  \item \textbf{实例可检索}：存在可公开检索的二创实例
        （视频、帖、弹幕）作为使用证据。
\end{enumerate}

标准 (2)(3) 把“谐音”与“空耳”区分开：
任何人都能给一个词找谐音，但只有社群实际在用的谐音
才是空耳词。\keyw{判定单位是社群行为，不是音系学}。

\subsection{语音距离指标（M 级）}

为了给标准 (1) 一个可排序的度量，本书构造
\keyw{拼音声韵母编辑距离}：

\begin{definition}[声韵母编辑距离（本书构造，\conf{M} 级）]\label{def:08:dist}
  把每个音节表示为有序对（声母 $i$，韵母 $f$），
  忽略声调。两词 $w, w'$ 的距离定义为音节序列上
  最小编辑代价：替换一个声母记 1 分，替换一个韵母记 1 分，
  增删一个音节记 1.5 分；总代价除以较长词的音节数得
  \keyw{归一化距离} $d(w,w') \in [0, +\infty)$。
  $d=0$ 为完全同音（不计声调）。
\end{definition}

三条使用警告。其一，这是\keyw{排序指标}：
它只用于把候选变体按近似程度排序，
\keyw{不}是音系学结论，数值大小没有感知意义。
其二，分值对“声母优先还是韵母优先”等权重设定敏感，
修订者可调整权重，但必须随结果公布所用权重。
其三，$d$ 小不等于等价关系成立——
等价关系由三条标准联合判定，$d$ 只服务标准 (1)。

\begin{example}[演示计算]\label{ex:08:dist}
  以“安迪”（\texttt{an di}）与其三个社群空耳变体为例：
  \begin{enumerate}[label=(\alph*),leftmargin=2.2em]
    \item “俺滴”（\texttt{an di}）：声韵全同，
          $d = 0$；
    \item “安定”（\texttt{an ding}）：音节数相同；
          首音节全同（0 分）；“迪--定”声母同（0 分）、
          韵母 \texttt{i}/\texttt{ing} 替换（1 分）；
          $d = 1/2 = 0.50$；
    \item “俺无敌”（\texttt{an wu di}）：三个音节对两音节，
          取最优对齐“安=俺（0）、迪=敌（0）、增‘无’（1.5）”；
          $d = 1.5/3 = 0.50$。
  \end{enumerate}
  三个变体的 $d$ 均不超过 $0.50$，
  且 (a)(c) 在奶蛋“高光”场景中高频共现、
  (b) 见于“安迪”的普通话化改写场景
  （\conf{C}，第一版教材口径），
  三条标准均满足，故均计入等价项集合。
\end{example}

\begin{nlmethod}[研究方法提示：做一遍的完整步骤]
  修订者复现本节方法的完整步骤：
  (1) 取衍生语清单（表 \ref{tab:06:langs}）；
  (2) 在选定平台以衍生语为关键词检索二创实例，
      记录每条候选空耳的出处链接与访问日期；
  (3) 按 (1)--(3) 三条标准逐项判定，
      记录每条的通过 / 淘汰及理由；
  (4) 对通过项计算 $d$ 并排序；
  (5) 公开判定表（候选、出处、判定、$d$ 值、权重设定）。
  样本规模、时间窗与平台清单必须事先写定——
  这一步骤本身就是第十二章“可复现性清单”的
  一个完整样例：\keyw{方法可复现，结论才可被推翻}。
\end{nlmethod}

\section{AIGC、版权与“变异权”的张力}\label{sec:08:copyright}

AIGC 把第三章的版权张力推向了一个新形态。
正典时代的侵权是\keyw{复制}（盗图、盗播），
AIGC 时代的侵权是\keyw{变异}（AI 重绘、批量生成）——
后者在法律上更难定性：变异产物与原形象的相似度
落在哪个区间构成侵权，目前没有稳定判例口径，
本书不给出法律意见（边界声明见表 \ref{tab:02:boundary}）。
可以指认的只有行为层面的结构：

\begin{enumerate}[label=\arabic*.]
  \item \textbf{生成即上游}：奶蛙案例显示，变异形象可以在
        平台侧“无主”地诞生（智能体头像），
        再由社群认领——侵权行为与侵权意图在起点处脱钩；
  \item \textbf{批量即规模}：单个变体的法律意义有限，
        权利人实际面对的是批量生成造成的
        \keyw{执法对象海量化}；
  \item \textbf{变异即传播}：官方对“丑化二创”的压制
        （第四章 \ref{sec:04:official}）恰好撞上
        M5 机制（\ref{sec:07:reverse}）——
        AIGC 供给的无限性与压制的有限性相遇，
        构成第九章“迷因生态缝隙”的生产端基础。
\end{enumerate}

\begin{nlcase}[案例 8-1：从“变形”到“变异”的权利视角]
  对比两个场景。场景一：某商家未经授权印刷奶龙图案销售
  ——经典侵权，权利边界清晰，维权路径成熟
  （快速维权机制，\conf{A}\cite{jiang2026nailong}）。
  场景二：某用户用图像模型把奶龙“变”成奶蛙并发布
  ——变异产物与原形象的相似度可辩、
  生成过程无“复制”动作、发布者与传播链解耦。
  两个场景在权利人眼中同属“形象失真”的损害
  （谢承恩列举的第一种侵权形态即“IP 形象严重失真”，
  \conf{A}\cite{jiang2026nailong}），
  但在现行工具与平台结构下，
  第二场景的可诉性、可执行性都显著更低。
  这一落差解释了为什么官方反应集中在
  \keyw{平台行为}（关评、拉黑账号）而非
  \keyw{图像本体}（要求删除变异形象）——
  平台行为是可执行的，图像本体不是。
\end{nlcase}

\section{本章小结}\label{sec:08:summary}

本章拆解了奶蛙二创的五环节生产链，
指出变异引擎的两个时点（2024 星野智能体、
2026 视频生成模型，均 \conf{C} 级），
并给出形式稳定性的经济学版本
（命题 \ref{prop:08:asym}：内容变异成本趋零、
形式构件模板固化，\conf{M} 级）。
方法论上，本章兑现了全书的可复现承诺：
空耳词等价关系的三条判定标准与
声韵母编辑距离（定义 \ref{def:08:dist}，
M 级排序指标），并给出完整复现步骤。
张力层面，AIGC 把侵权从“复制”变为“变异”，
可执行的对象从图像本体移向平台行为——
这是第九章讨论官方与二创关系的生产端前提。
遗留问题：生产端解释了“供给为何无限”，
但没有回答“需求为何存在”——
那需要回到第七章的机制假设与第十一章的受体研究；
同时，供给泛滥的外部性如何被治理
（平台、权利人、社群三方），
是第九章产业分析与第十章校园协商的共同背景。

\begin{nlexercises}
  \item 用定义 \ref{def:08:dist} 计算
        “嘎嘎滴辣虾”（\texttt{ga ga di la xia}）与
        “刚刚点了谁”（\texttt{gang gang dian le shui}）
        的归一化距离，写出你的对齐方案。
        再解释：为什么即使 $d$ 较大，
        该空耳仍可成立（提示：标准 (2)(3)）。
  \item 命题 \ref{prop:08:asym} 的推论二预测
        “改结尾套路”的二创传播表现更差。
        设计一个最小对照：
        选什么样的样本、编码哪两个变量、
        用什么统计量比较，并说明
        如何排除“改结尾的作者本来就不熟练”这一混杂因素。
  \item “生产成本由个体承担、拥挤成本由公共信息流承担”
        是一个外部性结构。仿照排放税的思路，
        提出一个平台侧的治理机制设计
        （如对高频发布者的某个约束），
        并分析它会触发第七章哪个机制的反弹
        （提示：M5）。
\end{nlexercises}

\nlkeywords AIGC、生产链、变异引擎、Seedance、成本不对称、
模板固化、空耳（soramimi）、声韵母编辑距离、
判定标准、低质内容、外部性、变异权

\chapter{玩梗接入产业：授权体系、二创周边与文化收编}\label{ch:09}

\begin{nlkeypoints}
  \begin{enumerate}[label=(\arabic*),leftmargin=2.4em]
    \item 正典与变种走上两条经济轨道：
          官方奶龙走\keyw{“内容 + 授权”双轮驱动}
          （200+ 品牌、2000+ SKU、57 国，\conf{A/B} 级），
          奶蛙走\keyw{零授权的可传播性经济}——
          内容层无人收费、形式层无法收费，
          被商业捕捉的是立场的周边化（立牌、玩偶）
          （本书分析，\conf{M} 级）。
    \item 2026-03 关评拉黑事件应读作“不彻底的抑制”：
          官方控制的是\keyw{可执行对象}（平台评论区、账号），
          而“捧腹大笑”作为\keyw{可传播性单位}
          在平台空间之外继续流通——
          \keyw{法律层面}与\keyw{意义层面}的
          正典—变种关系由此分离（\conf{C} 级事件 + \conf{M} 级分析）。
    \item Hebdige 的“收编”在奶蛙案例中呈现为双向运动：
          商业收编迷因（周边商品化），
          迷因也反向收编商业（品牌借梗营销）；
          第一版教材记录的淘宝周边与
          \en{“奶龙正在成为艺术品”}评述
          是这场双向运动的两端（均 \conf{C} 级）。
  \end{enumerate}
\end{nlkeypoints}

第一章留下过一个张力表述：一只“不属于任何官方”的青蛙，
寄生在一个有专业法务团队的正版 IP 旁边生长。
本章把这个张力展开为完整的经济与法律分析。
先立底图（融合文化与可传播性），再看两条轨道各自的
收益结构，然后处理 2026-03 事件的法律—意义双重读法，
最后讨论收编与艺术化的双向运动。
本章继续遵守全书纪律：产业数字全部来自 \conf{A/B} 级
公开报道，二创周边现象来自 \conf{C} 级第一版教材记录，
凡属本书的经济学解释一律标注 \conf{M} 级。

\section{底图：融合文化与可传播性}\label{sec:09:base}

第二章已给出媒介结构的位置：正典奶龙由制作方控制分发，
奶蛙的流通则完全由二创者的转发、二创与线下搬运驱动。
本节把这张底图补上经济维度。Jenkins 等人提出，
在融合文化环境下，内容的价值不取决于制作方的
“黏性”投放，而取决于受众的\keyw{可传播性}
（\en{spreadability}）行为——受众愿意并且能够
携带内容去哪里，内容的价值边界就在哪里
\cite{jenkins2013spreadable}。

用这个视角回看第四章的一个悬置问题：
“结尾集体大笑”为什么值得当作\keyw{可传播性单位}分析
（\ref{sec:04:audio}）。答案现在可以给出：

\begin{proposition}[可传播性单位（教学命题）]\label{prop:09:unit}
  一个文化构件的可传播性，取决于它在搬运中
  \keyw{不依赖原始语境}的能力。“结尾集体大笑”
  是奶蛙语料中可传播性最高的构件：
  它不依赖任何具体压力桥段（内容可换）、
  不依赖原始上传者（署名可弃）、
  不依赖特定平台（载体可换），
  只依赖观看者的“接梗”能力（社群在场）。
\end{proposition}

这条命题（\conf{M} 级）解释了第九章的核心张力：
\keyw{可传播性最强的构件，恰恰是法律上最无法据为己有的构件}。
笑声是录音制品，套路是思想而非表达，梗图与视频
分散在无数匿名账号——官方奶龙可以把形象授权给 200 家品牌，
却无法对“结尾一起笑”这个形式收取任何费用。
可传播性与可产权化在奶蛙案例中呈反向关系，
这是迷因经济的第一个结构性事实。

\section{正典轨道：“内容 + 授权”双轮驱动}\label{sec:09:canon}

第三章已经给出正典奶龙的产业量化口径
（表 \ref{tab:03:numbers}），本节不再复述数字，
只分析其经济结构的三个特征。

\begin{enumerate}[label=\arabic*.]
  \item \textbf{收入的双轮结构}。谢承恩自述：
        2020 年初“商业化也还没起步”，其后
        “从广告合作起步，逐步拓展到商业授权，
        积极探索‘内容 + 授权’双轮驱动的模式”
        \conf{A}\cite{jiang2026nailong}。
        内容轮（平台分账、播放流量）与授权轮
        （品牌合作、衍生品）的比重变化
        没有公开数据，但 200+ 授权品牌、2000+ SKU、
        57 国销售、25 万元政府扶持奖励
        \conf{A/B}\cite{jiang2026nailong,chinadaily2026six}
        说明授权轮已是产业主体。
  \item \textbf{授权的本质是“形象租赁”}。品牌方购买的
        不是剧集内容而是\keyw{形象的情感附加值}——
        “治愈、欢乐”的定位\conf{A}\cite{jiang2026nailong}
        使奶龙成为低风险的情绪符号，
        适配茶饮、乐园、大屏等家庭场景
        （第三章案例 3-1）。
  \item \textbf{维权是授权的前置条件}。法务团队、
        快速维权机制、海外维权协助
        \conf{A}\cite{jiang2026nailong}
        不是授权体系的外挂，而是其成立条件：
        \keyw{只有能把形象圈起来的权利人，才能把形象租出去}。
        维权强度因此可以读作授权体系健康状况的代理指标。
\end{enumerate}

\begin{nlcase}[案例 9-1：表情包的 80 亿次]
  中国日报网报道给出一个耐人寻味的口径：
  奶龙微信表情包累计发送量突破 80 亿次
  （截至 2025 年底，\conf{B}\cite{chinadaily2026six}）。
  表情包位于授权体系与迷因流通的\keyw{交界地带}：
  它是官方发布、免费分发、用户自发使用的——
  形象受控，用途不受控。80 亿次发送既是
  授权体系对迷因式流通的\keyw{主动模仿}
  （官方把形象做成“可传播性单位”投放），
  也说明了授权体系的天花板：
  用户对表情包的每一次再加工、再配文，
  都在官方控制之外。对照 2026-03 对奶蛙的拉黑可见，
  官方并不反对迷因式流通\keyw{本身}，
  反对的是\keyw{失控的流通}——
  这一区分是理解两条轨道摩擦的关键。
\end{nlcase}

\section{变种轨道：零授权的可传播性经济}\label{sec:09:variant}

第一版教材第八部分记录了奶蛙周边的民间售卖：
“淘宝售卖的奶脚抱枕、奶蛙玩偶、亚克力”，
“2m 奶蛙人形立排和奶蛙六合一战神金刚”，
并评论“紧随抽象潮流的商家们也纷纷上架奶蛙周边”
\conf{C}（第一版教材口径；具体商品图本书不复制，
权利状态不明）。

这条轨道的经济结构与正典轨道逐项相反：

\begin{table}[htbp]\centering
  \caption{两条轨道的经济结构对照（本书整理；
          官方侧为 \conf{A/B} 级口径，民间侧为 \conf{C} 级记录，
          “对照”列为 \conf{M} 级分析）}
  \label{tab:09:two}
  \small
  \begin{tabularx}{\textwidth}{@{}lXX@{}}
    \toprule
    维度 & 正典轨道（官方奶龙） & 变种轨道（奶蛙周边） \\ \midrule
    权利状态 & 权利清晰、集中持有 &
      权利模糊：形象未经授权，署名权在匿名创作者 \\
    收益形式 & 授权费、衍生品分成、平台分账 &
      商家现货差价；创作者基本不获酬 \\
    质量控制 & 官方审校、SKU 管理 &
      无控制：从亚克力到 2 米立排质量参差 \\
    法律风险 & 风险在侵权方 &
      风险在售卖方：随时可能被维权 \\
    传播逻辑 & 黏性投放 + 节点营销 &
      可传播性搬运 + 刷屏 \\ \bottomrule
  \end{tabularx}
\end{table}

表中最要紧的一行是\keyw{创作者不获酬}。
奶蛙的爆火使形象产生了真实的商业价值
（周边有市场、立牌有仪式功能），
但按第四章的处理，考据指向的创建者是
\keyw{权利主体}而不只是证据节点（规则 E3）；
社群传播中的每一次免费使用，
都在没有报酬通道的情况下增殖着这一形象的价值。
这不是本书对任何一方行为的法律指控
——恰恰相反，报酬通道的缺失是
\keyw{结构性的}：匿名使授权无从谈起，
授权的缺席使商业化只能在灰色地带进行。
第十二章把这一问题归入“研究伦理三条线”的第三线。

\begin{nlcritique}[批判性反思：谁在赚迷因的钱]
  民间周边经济中存在一个容易忽视的不对称：
  \keyw{承担形象诞生风险的人（创作者）与
  收获售卖利润的人（商家）不是同一批人}。
  第一版教材以轻松笔调记录周边“壮大”，
  本书保留其史料价值，但必须补上另一面：
  无授权售卖使商家处于持续的法律风险中，
  也使消费者买到的是权利状态不明的商品；
  若迷因热度退潮，存货风险完全由售卖端承担。
  读者在观察这条轨道时，应同时看到它的活力
  与它的脆弱——两者是同一个权利真空的两面。
\end{nlcritique}

\section{2026-03 事件的双重读法：法律层面与意义层面}\label{sec:09:split}

第四章留下了一个待展开的论证：
“断裂叙事”（奶蛙“正式脱离了奶龙”）与
“连续叙事”（授权、维权、关评的产业线）
不是对错关系，而是分别描述了\keyw{意义层面}
与\keyw{法律层面}的正典—变种关系。
现在给出这个论证。

\textbf{法律层面}：奶蛙与奶龙从未分离。
变异形象与正典形象的相似性正是侵权的构成基础
（也正因为相似，官方的拉黑才有明确对象）；
只要相似性存在，法律层面的支配关系
（权利人可以主张什么）就持续存在。
关评、拉黑、呼吁停止丑化，都是这一层面的动作
\conf{C}\cite{shoutanji2026naiwa}。

\textbf{意义层面}：奶蛙确实脱离了奶龙。
第六章的谱系分析显示，奶蛙的辨识单元
（黄色 + 大笑 + 膨腹）已经可以脱离正典设定独立运行；
“潘多拉奶盒”叙事是社群对这一脱离的自觉表述；
官方评论区“仍有无数奶蛙图片”
\conf{C}\cite{shoutanji2026naiwa}
则说明意义层面的流通不再服从法律层面的边界。

\begin{proposition}[双重读法（教学命题）]\label{prop:09:dual}
  2026-03 的官方抑制之所以“不彻底”，
  是因为抑制工具（关评、拉黑）只能作用于
  \keyw{法律层面可执行的对象}（官方账号评论区、
  特定账号），而奶蛙的存在方式已经转移到
  \keyw{意义层面不可执行的空间}（个体创作、
  群聊流通、线下立牌）。两个层面的错位
  构成“迷因生态的缝隙”：\keyw{失败的控制}
  不是控制不努力，而是控制对象与控制工具
  在层面上的错配。
\end{proposition}

命题 \ref{prop:09:dual}（\conf{M} 级，
事件基础为 \conf{C} 级）同时回答第三章的三条承诺：
其一，“官方控制 vs 迷因生态缝隙”的框架即本命题；
其二，“如果关评拉黑是彻底的，奶蛙现象在时间上会被截断”
——反事实推理的依据正是层面错配：
不存在能作用于意义层面的执行工具；
其三，受众外溢（第三章）的第三层
“民间再生产”之所以能独立完成“大学生化”，
正因为它在意义层面运行，不受内容内核层
与平台结构层的官方意图约束。

\section{收编：商业与迷因的双向运动}\label{sec:09:recoup}

Hebdige 描述了亚文化的经典命运：风格作为抵抗被
主流与商业\keyw{收编}（\en{incorporation}），
转化为商品与时尚\cite{hebdige1979subculture}。
奶蛙案例为这个理论提供了一个当代续集——
收编不是单行道。

\subsection{商业收编迷因}

民间周边（\ref{sec:09:variant}）是收编的第一形态：
把立场的符号（“反精致”的姿态，第七章 M3）
变成货架上的商品。亚文化理论预测，
收编会稀释风格的抵抗含义——
当奶蛙玩偶与正版 IP 玩偶并列销售时，
“转发一只丑蛙”的表态成本趋近于购买一只玩偶的成本。
第七章 M3 的“阶级自嘲”批判
（\ref{sec:07:m3} 的批判框）在这里得到经济侧印证：
\keyw{收编把表态变成了消费，把不合作变成了生意}。

\subsection{迷因反向收编商业}

奶蛙案例的特殊之处在于第二个方向：
\keyw{迷因逻辑也在改造商业行为}。
官方奶龙的六周年营销（五城大屏、单曲出道、
“歌手”人设\conf{B}\cite{chinadaily2026six}）
在形式上就是对迷因式流通的学习——
把 IP 做成“可传播性单位”投放；
2026-01-25 登上总台《动漫盛典》的“高姿态”
\conf{B}\cite{chinadaily2026six}
与同月内社群里奶蛙的刷屏，构成同一文化市场上
两种流通逻辑的并置。一篇 2026 年 7 月的评述
甚至以“奶龙，正在成为一件艺术品”为题，
讨论正典奶龙被“艺术化”“档案化”的趋势
\conf{C}\cite{yuhuoxingxing2026}——
无论该评述的判断是否成立，
“一个子供向 IP 被讨论为艺术品”这件事本身，
就是迷因文化抬升其讨论语域的证据。

\begin{nltermbox}{术语：收编（incorporation）}
  Hebdige 的亚文化理论术语：主流文化与商业把
  亚文化风格转化为可销售商品、可传播时尚，
  从而消解其抵抗性的过程\cite{hebdige1979subculture}。
  本书用法提醒：收编预设了“抵抗”的存在，
  而 M3 的判别问题（\ref{sec:07:m3} 批判框）
  尚未解决——若“反精致”本就是人设消费，
  则周边商品化不是“收编抵抗”，
  只是普通的供给跟随需求。
  \keyw{使用该术语前，请先确认抵抗性证据}。
\end{nltermbox}

\section{本章小结}\label{sec:09:summary}

本章在融合文化与可传播性的底图上，
对照了两条经济轨道：官方的“内容 + 授权”双轮
（维权是授权的前置条件；表情包 80 亿次说明官方
主动模仿迷因式流通但对失控流通保持警惕），
与民间的零授权可传播性经济（创作者不获酬的结构性
报酬真空）。对 2026-03 事件给出了法律—意义双重读法
（命题 \ref{prop:09:dual}）：抑制工具只能作用于
法律层面的可执行对象，意义层面的流通无法执行，
层面错配即“迷因生态的缝隙”。收编分析指出
商业收编迷因与迷因反向收编商业的双向运动。
遗留问题有二：其一，两条轨道在校园空间里如何交汇
（立牌同时是零授权符号商品与社群仪式物），
第十章处理；其二，本章对“意义层面”的分析
全部依赖社群文本与 \conf{C} 级材料，
意义脱离的实际程度需要受体研究提供行为证据
（第十一章）。

\begin{nlexercises}
  \item 用命题 \ref{prop:09:unit} 的“不依赖原始语境”
        三个维度（内容、署名、载体），
        分别检验“结尾集体大笑”“奶蛋”“某品牌联名奶茶杯”
        三个构件，排序其可传播性并说明依据。
  \item 案例 9-1 认为官方“不反对迷因式流通本身，
        反对失控的流通”。请为“失控”给出一个
        \keyw{可操作的判据}（什么条件下流通算失控），
        并讨论该判据分别从官方视角与社群视角
        写出来会有什么不同。
  \item 表 \ref{tab:09:two} 指出民间轨道中
        “创作者基本不获酬”。有观点认为
        “匿名创作本就放弃了报酬，无所谓真空”；
        另一观点认为“社群中人人都在免费使用，
        报酬通道的缺失是结构不公”。
        写出两种观点各自的隐含前提，
        并说明本章的 E3 规则（权利保留）支持哪一方的
        \keyw{分析姿态}、不支持哪一方的\keyw{结论}。
\end{nlexercises}

\nlkeywords 授权体系、内容 + 授权、可传播性、
可传播性单位、双轨分发、零授权经济、收编、
法律层面、意义层面、迷因生态缝隙、创作者报酬

\chapter{奶蛙进校园：空间实践、组织形态与低强度协商}\label{ch:10}

\begin{nlkeypoints}
  \begin{enumerate}[label=(\arabic*),leftmargin=2.4em]
    \item 五校案例（哈工大、上交、海大、北航、浙大，均 \conf{C} 级）
          在时间上呈\keyw{扩散序列}：从单一立牌亮相（哈工大）
          到流动实践（上交）、复数立牌（海大）、
          组织化同好会（北航）、官方物料采纳（浙大）——
          与创新扩散的采用者类别在形态上相容
          （\conf{C} 级案例 + \conf{M} 级结构归位）。
    \item 立牌是奶蛙在校园的\keyw{物化符号}：
          它把线上迷因固定为可合影、可搬运、可巡游的实物，
          使符号实践进入物理空间并接受校方管理——
          2026-09-28 北航“9 小时立牌”事件展示了
          \keyw{移除—纪念—再创作}的低强度协商回路。
    \item 本章给出\keyw{协商强度谱}（\conf{M} 级）：
          按“校方反应—学生反应—是否留下持久物”三维
          给案例分级；北航事件的三维读数表明，
          校园协商以“边界重申 + 符号存续”收尾，
          未见冲突性对抗。
  \end{enumerate}
\end{nlkeypoints}

前九章的分析都发生在“线上”——平台、语料、生产链。
本章进入物理空间：奶蛙如何出现在毕业典礼、绿园、
招生宣传与同好会群聊里。材料全部来自第一版教材的
高校活动记录（\conf{C} 级），单一来源、无法独立印证，
这是本章证据地位的底限，分析结论全部标注 \conf{M} 级。
本章回答问题 Q3 的最后一环：当迷因从信息流走进校园，
它的传播结构、组织载体与权威关系发生了什么变化。

\section{案例库：五校事件整理}\label{sec:10:cases}

\begin{table}[htbp]\centering
  \caption{高校奶蛙事件整理（据第一版教材，全部 \conf{C} 级；
          “类型”列为本书归类，\conf{M} 级）}
  \label{tab:10:cases}
  \small
  \begin{tabularx}{\textwidth}{@{}lllX@{}}
    \toprule
    学校 & 时间 & 类型 & 事件 \\ \midrule
    哈尔滨工业大学 & 2026-06 & 仪式亮相 &
      毕业典礼暨学位授予仪式出现第一个广为人知的奶蛙立牌，
      被视为各校“立牌争先赶超”阶段的开启 \\
    上海交通大学 & 2026 年 & 流动实践 &
      率先发起“流动奶蛙”：把立牌带到校园各处制造偶遇；
      最早的流动奶蛙至今仍在使用 \\
    中国海洋大学 & 2026 年 & 复数立牌 &
      “奶 OUC”四奶蛙立牌，第一个广为人知的复数立牌 \\
    北京航空航天大学 & 2026-06-24 起 & 组织化 &
      奶龙同好会成立；2026-08 初破 100 人；
      26 级新生入学后于 2026-10 初破 200 人 \\
    北京航空航天大学 & 2026-09-28 & 空间冲突 &
      学院路校区绿园投放流动奶蛙立牌，旁现丰川祥子立牌；
      当晚被保安带走，存活约 9 小时 \\
    浙江大学 & 2026 年 & 官方采纳 &
      某地招生宣传立牌中出现浙大奶蛙形象 \\
    \bottomrule
  \end{tabularx}
\end{table}

表 \ref{tab:10:cases} 的“类型”列暗含一个
时间—结构假设（\conf{M} 级）：五个案例恰好覆盖了
创新扩散理论中采用者的经典类别序列——
\term{创新者}（哈工大：首个、仪式场合、高风险亮相）、
\term{早期采用者}（上交：发明新玩法“流动奶蛙”）、
\term{早期大多数}（海大：把单点实践扩展为复数形态）、
\term{后期大多数}（各校同好会的组织化复制，
北航为其显例）、以及\keyw{守门人采纳}
（浙大：官方物料出现该形象）\cite{rogers2003diffusion}。
必须强调两点：其一，这是\keyw{形态上的相容}，
不是对扩散过程的实证检验——案例时间粒度太粗
（精确到月），且五校样本是“出圈案例”的截尾样本
（见 \ref{sec:10:network} 的现存样本偏差讨论）；
其二，序列若成立，它预测 2026 年下半年还会有
更多“守门人采纳”型事件（官方物料化），
这是一个可跟踪的预测。

\section{空间实践：立牌、流动与合影}\label{sec:10:space}

\subsection{立牌：从梗图到实物}

立牌是奶蛙符号在物理空间的锚点。第六章指出，
“摆架势”是衍生动作中“从形象动作迁移为使用者
身体姿态的通道”；立牌把这一通道反过来用——
\keyw{把形象固定为可合影的实物}，使身体程式
（摆架势、合影）有了一个空间坐标。
用 Jenkins 等人的语言说，立牌是可传播性的
\keyw{物化形态}：内容不依赖特定载体传播
（命题 \ref{prop:09:unit}），
而立牌让“载体”本身可以被搬运、被巡游、被拍下再传回线上
\cite{jenkins2013spreadable}。

\subsection{流动奶蛙：制造偶遇}

上交的“流动奶蛙”把立牌从静物变为\keyw{移动事件}：
“把奶蛙带到学校的各个角落，在任意一处地方偶遇惊喜”
（\conf{C}，第一版教材口径）。
这一实践的文化逻辑是\keyw{偶遇的仪式化}——
把概率性事件（碰巧刷到梗）改写为空间性事件
（碰巧遇见蛙）。Turner 的阈限理论提供了一种读法：
毕业典礼、校园绿地本是日常结构的通道，
立牌的闯入制造了一个\keyw{微阈限}时刻——
参与者短暂脱离日常角色（学生、行人），
进入“与一只大笑青蛙合影”的反结构状态，
合影与转发则是阈限体验的\keyw{存档}
\cite{turner1969ritual}。
这一读法是本书的诠释（\conf{M} 级），
它的可检验推论是：流动奶蛙的合影率
（遇见者拍照比例）应显著高于固定装置的合影率——
因为“偶遇”比“陈设”更接近阈限体验的触发条件。

\subsection{FoMO 与纳新}

第二章预言了“准入测试在纳新中反复出现”
（\ref{sec:02:meta}），北航同好会的增长提供了
检验材料：2026-06-24 成立，8 月初破 100 人，
26 级新生入学后 10 月初破 200 人
（\conf{C}，第一版教材口径）。
新生入学与会员数翻倍的时间重合提示
\keyw{错失恐惧}（\en{FoMO}）的动因结构：
FoMO 指对他人正在经历自己缺席的有益经验的
普遍性忧虑，它驱动社交媒介使用与从众行为
\cite{przybylski2013fomo}。
在纳新语境下，刷屏的立牌照片与同好会群聊
构成“大家都在里面”的可见证据，
FoMO 把兴趣转化为入会动作
（\conf{M} 级分析；第十一章问卷含 FoMO 量表
以便检验）。

\section{组织形态：同好会}\label{sec:10:org}

第一版教材记录的北航奶龙同好会成立目标
（\conf{C}）是：“让原飞梦社的抽象老登们
有个神圣的水群”。“水群”（社群用语）指低强度的
日常闲聊；记录同时说明该会因“奶蛙前所未有的强横热度”
而人数激增，“成为飞梦之后北航校园文化新的中坚力量之一”
\conf{C}。

这条记录的学理价值在于它显示了\keyw{组织的双轨性}：
同好会的名义功能（水群）与实际功能（校园迷因的
组织载体：纳新、贺图、表情包生产、立牌活动）
是分开的。Hebdige 笔下的亚文化很少有正式组织，
奶蛙的同好会则是\keyw{半正式}的——
有成立日期、有人数里程碑（100 人贺图、200 人贺图）、
有名称与身份，但无正式社团注册的记载
（\conf{C} 级材料未提及注册状态，本书不推断）。
第二章“准入测试”的功能分析在此落位：
贺图与校徽奶蛙表情包（“数分奶蛙”“捧着航概的奶蛙”
\conf{C}）是\keyw{校本化的迷因资产}，
它们把全国性梗转写为本校身份标记——
持有并使用校本奶蛙表情包，
即是“本校圈内人”的默示声明。

\section{协商：9 小时立牌事件与协商强度谱}\label{sec:10:negotiation}

\subsection{事件经过}

据第一版教材记录：2026 年 9 月 28 日，
一个匿名者在北航学院路校区绿园投放流动奶蛙立牌，
旁边随后出现丰川祥子立牌；学院路学生纷纷与之合影；
当晚立牌被保安发现并带走，存活约 9 小时；
同好会“对此深表遗憾”，成员自发在群内刷屏一段
化用宗教语体的赞诗（“我的肉化作黄桃罐头……
为使世界获得笑容”，\conf{C}，第一版教材记录；
本书只作语体分析，不复制全文）。

\subsection{低强度协商的三维读数}

事实手册指出该事件的关键特征：
官方（保安/校方）与自组织（同好会）之间
\keyw{并未发生冲突性对抗}，而是以
“移除—纪念—再创作”完成一轮低强度协商。
本书把这一观察升级为可分级的教学框架：

\begin{table}[htbp]\centering
  \caption{立牌活动的协商强度谱（本书构造的教学模型，\conf{M} 级；
          案例读数基于 \conf{C} 级材料）}
  \label{tab:10:spectrum}
  \small
  \begin{tabularx}{\textwidth}{@{}lXXX@{}}
    \toprule
    维度 & 低强度（0） & 中强度（1） & 高强度（2） \\ \midrule
    校方反应 & 例行移除、未追责 &
      公开通告、约谈组织者 &
      纪律处分、定性为违规活动 \\
    学生反应 & 自嘲、默默再投放 &
      纪念文本、纪念二创 &
      组织化抗议、对抗性表态 \\
    持久物 & 无 & 线上纪念物 &
      制度化安排（如申请合规展位） \\ \bottomrule
  \end{tabularx}
\end{table}

北航事件的读数：校方反应 0（例行移除，
无公开追责记录）；学生反应 1（纪念文本与刷屏，
无对抗）；持久物 1（线上纪念物；且“流动奶蛙”
实践本身在事件后延续——最早的流动奶蛙仍在使用）。
合计 2，属\keyw{低—中强度区间}。
这个分级的用途不是给事件打分，而是让\keyw{比较}成为可能：
不同学校、不同时间的立牌事件可以放在同一谱上对照，
“协商强度随什么变量变化”（学校管理风格？
投放场合的正式程度？立牌内容？）由此变成可提问的经验问题。

\subsection{赞诗的语体分析}

赞诗值得单独一读，因为它是第六章“抽象文化”
定义中\keyw{对严肃语体的戏仿}的最佳样本：
它化用宗教圣体文的句式（“我的肉……谁若吃了……
必要……直到永远”），把一只橡皮鸭子式的玩偶立牌
写进最庄重的语体框架。语体错配在这里承担三件事：
其一，\keyw{哀悼的合法化}——为一个玩偶立牌“致哀”
本会尴尬，宗教语体的庄重性允许成员
以夸张完成真诚；其二，\keyw{圈内测试}——
识别这段文本的互文来源需要特定文化库存，
读懂即认证；其三，\keyw{软性表态}——
刷屏行为本身向“带走立牌的管理方”传递了
不满，但语体的戏谑性使该表态随时可撤回
（“我们只是在玩梗”）。
这与第七章 M2（群体稀释）的“印象缓冲”结构同源：
\keyw{用可撤回的语体，说不可撤回的态度}。

\begin{nlcritique}[批判性反思：管理方的视角与道德恐慌的风险]
  本章的案例都以学生视角记录（来源即学生教材），
  校方视角完全缺席。公平起见必须列出管理方的
  合理关切：未报备的立体装置涉及公共空间秩序、
  消防与通行安全、以及“下一个投放者放什么”
  的滑坡问题（丰川祥子立牌的出现已经演示了
  跟随投放的即兴性）。同时要警惕反向的过度反应：
  若把校园玩梗上升为“不良文化”进行运动式清理，
  即落入 Cohen 所述\keyw{道德恐慌}的框架——
  把一个低危害的风格现象建构成对秩序的威胁
  \cite{cohen1972folkdevils}。
  协商强度谱的意义正在于提供第三条路：
  双方都以低强度动作表达立场（移除对纪念），
  边界被重申（不经报备的装置会被移走），
  符号实践得以延续（流动奶蛙照常）。
  \keyw{低强度协商不是冲突的失败，而是冲突的替代}。
\end{nlcritique}

\section{传播网络：校园作为高密度混合网络}\label{sec:10:network}

第七章 M4 预测“线下共同在场密度延长素材寿命”
（假设 \ref{hyp:07:m4}），并承诺校园传播网络的结构分析。
本节给出框架（分析框架 \conf{M} 级）。

\subsection{弱连接与强连接的分工}

Granovetter 的弱连接优势命题：新信息更可能
经由弱连接（泛泛之交）到达，因为强连接圈子内部
信息早已同质\cite{granovetter1973weak}。
校园网络是罕见的\keyw{高密度混合}结构：
宿舍与班级提供强连接（复读、互发表情包，
强化趣味、延迟疲劳，对应 M4），
跨院系社团、校园广场与信息流提供弱连接
（新变体、新梗的输入，对应第七章“新变体重新激活”
的 $\omega$ 项）。预测：\keyw{一个校园迷因的寿命
取决于强连接复读维持与弱连接变异供给的乘积}——
只有强连接而无新变体，趣味在复读中缓慢耗尽；
只有弱连接而无复读，梗无法沉淀为本校身份标记。
校本表情包（数分奶蛙、航概奶蛙）正是
弱连接输入被强连接复读固化后的\keyw{沉淀物}。

\subsection{多群 SIR 与现存样本偏差}

第五章批判框承认单群 SIR 模型的第一处局限
（“高校”“平台”混在一个 $N$ 里），并要求第十章
给出多群版本。本书不给完整方程（超出教学需要），
只说明多群化的两个必要改动：
其一，把 $N$ 拆为校园群与公域群，
群间接触率（弱连接强度）与群内接触率
（强连接密度）分开设参；
其二，校园群引入\keyw{学期节律项}——
寒暑假相当于大规模的暂时离场
（$S$、$I$ 的时间调制项），
它解释了北航同好会人数“新生入学后突破 200 人”
的时间点（\conf{C}）：\keyw{每年九月是一次
系统性的易感者注入}。

同时必须重申偏差警告：本章五校案例是
\keyw{被记录、被出圈}的现存样本——
全国高校中奶蛙活动存在但未被记录的院校
不可观测，扩散序列（表 \ref{tab:10:cases}）
因此是截尾样本上的形态归纳，
不是全总体估计。这与第四章“不可复核性”
（\ref{sec:04:irrevocable}）构成一对方法论提醒：
\keyw{网络文化研究的样本，一端在消失
（证据损耗），一端在缺席（未记录）}。

\section{本章小结}\label{sec:10:summary}

本章整理了五校案例（全部 \conf{C} 级），
给出扩散序列假设（采用者类别的形态相容，
\conf{M} 级）、空间实践分析（立牌作为物化符号、
流动奶蛙作为微阈限事件、FoMO 驱动纳新）、
组织形态分析（同好会的名义—实际双轨功能、
校本化表情包作为身份标记）、协商分析
（9 小时立牌事件、协商强度谱、赞诗的语体三功能）
与网络框架（强连接复读与弱连接供给的乘积、
学期节律作为易感者注入）。
全书最重的证据警告也落在本章：
所有案例均为单一来源的出圈样本，分析结论
无一例外是 \conf{M} 级诠释。遗留问题：
本章描述了协商的发生，但未测量协商的心理后果
——参与者、旁观者与管理者各自如何理解这场
“9 小时”，需要第十一章的访谈设计进入现场。

\begin{nlexercises}
  \item 用表 \ref{tab:10:spectrum} 的三维谱，
        为你所在学校设想一次假想的立牌投放事件，
        分别写出“低—中强度收尾”与“高强度升级”两个分支，
        并指出哪个环节的哪个动作是升级的转折点。
  \item “流动奶蛙合影率高于固定装置合影率”
        （\ref{sec:10:space} 的可检验推论）：
        设计一个不干预现场的观察方案
        （观察点选择、观察时段、计数规则、
        遇到的混杂因素及处理）。
  \item 本章称赞诗“用可撤回的语体，说不可撤回的态度”。
        请找一个你所见的其他社群文本（如弹幕、
        群公告、评论区接龙）检验这一概括：
        若成立，指出其可撤回装置；若不成立，
        说明该文本用了什么别的机制来处理表态风险。
\end{nlexercises}

\nlkeywords 立牌、流动奶蛙、奶 OUC、同好会、
创新扩散、采用者类别、阈限、错失恐惧（FoMO）、
弱连接、强连接、协商强度谱、低强度协商、
道德恐慌、校本化、现存样本偏差

\chapter{受体研究设计：问卷、编码与访谈}\label{ch:11}

\begin{nlkeypoints}
  \begin{enumerate}[label=(\arabic*),leftmargin=2.4em]
    \item 本章交付三套\keyw{研究设计}而非研究结果：
          面向 H7.1--H7.4 的问卷（五模块、含假设映射表）、
          面向窄通道与成本不对称命题的内容编码方案
          （含双人独立编码与一致性要求）、
          以及面向立牌活动参与者的半结构化访谈提纲。
          全部设计是供施测的工具，\keyw{本书不含任何经验数据}。
    \item 设计的三个方法论支柱：
          \keyw{学期节律对齐}（压力机制的检验窗口）、
          \keyw{不对称指标}（区分功能机制与消遣/算法解释）、
          \keyw{行为—文本—自陈三角}（防止单一来源偏差）。
    \item 伦理条款先于工具交付：
          匿名化、知情同意、不收集可识别信息、
          与第四章 E1--E3 规则衔接——
          受体研究中“研究对象是人”，
          证据分级制度之外另有全套义务。
  \end{enumerate}
\end{nlkeypoints}

第七章留下了五个待检验假设，第六章与第八章
留下了两个待检验命题，第十章留下了未测量的协商后果。
本章的任务是把这些“欠账”变成可以开工的设计。
必须在一开始就重申全书最重要的声明：
\keyw{本章不含任何测量结果}。所有量表、编码表、
提纲与示例数据都是设计（\conf{M} 级），
它们的价值在于可被批评、修订、施测与再修订
——这正是本书与“梗评”文章的最后一道分界线：
前者交付可检验的工具，后者交付不可检验的断言。
中文网络传播研究的问卷与内容分析的一般方法，
可参照彭兰的网络传播学教材与中文迷因研究的
代表性梳理\cite{penglan2007meme,caoqi2015meme}，
本章的设计与其方法传统对接，不另起炉灶。

\section{研究对象与抽样}\label{sec:11:sampling}

\subsection{总体与抽样框}

\begin{definition}[研究总体（本书界定）]\label{def:11:pop}
  研究总体为\keyw{近 12 个月内在中国大陆高校注册就读、
  且使用任一主流社交平台的在校学生}；
  目标子总体为其中\keyw{接触过奶蛙类内容}者
  （操作化：能在提示下识别奶蛙形象，
  即第六章归并标准 (C3) 的个体版本）。
\end{definition}

抽样方案的诚实版本是：\keyw{不存在可获得的完整抽样框}
（高校学生名单不可得，接触者不可枚举），
因此任何施测都是非概率抽样。设计的应对有三：
其一，\keyw{配额}：按年级（尤其区分毕业年级与非毕业年级）
与院校层次设置配额下限——毕业年级是 H7.1 与 H7.3
的关键对比组；其二，\keyw{多入口}：同好会群、
课程班级、校园公共账号三入口分开发放，
降低单入口的社群同质性；其三，\keyw{报告偏差}：
无论样本多大，结论只表述为“样本内关系”，
不得外推为“大学生总体”。
孙立平对自嘲式表达之社会心态背景的讨论
\cite{sunliping2010} 在这里进入设计动机：
若压力—自嘲的关联是广泛的社会心态，
样本内应能观察到奶蛙使用与一般自嘲表达的
\keyw{区分度}——这是第七章竞争性解释的又一条判别线。

\subsection{样本量的量级}

教学设计的参考量级：若目标是检出
“两组频率差 $\ge 15$ 个百分点”的差异
（显著性水平 0.05、功效 0.80 的双侧检验），
每组约需 130--150 份有效问卷；
考虑配额与废卷，总投放量建议不低于 600 份。
这一估算按常规两比例检验的样本量公式即可复核，
本书不在此展开推导（\conf{M} 级）。

\section{问卷设计}\label{sec:11:questionnaire}

问卷分五模块（A--E），设计原则：
每题服务于一个明确的假设或对照，
不设“你觉得奶蛙为什么火”这类开放式归因题
——归因题测到的是受访者的民间理论，不是机制本身。

\begin{table}[htbp]\centering
  \caption{问卷模块与假设映射（\conf{M} 级设计）}
  \label{tab:11:modules}
  \small
  \begin{tabularx}{\textwidth}{@{}lXX@{}}
    \toprule
    模块 & 内容 & 服务的假设 / 对照 \\ \midrule
    A 使用行为 & 各类内容的使用频率（奶蛙类 / 总娱乐类，
      按学期时段分层回填） & H7.1 的选择性上升；
      算法解释的曝光控制 \\
    B 动机与情绪 & 压力自评（PSS 风格简表）+ FoMO 量表 +
      使用情境勾选（自述 / 群聊复读 / 评论区） &
      H7.1、H7.2、H7.4、FoMO 动因 \\
    C 圈内知识 & 场景语法测验（衍生语—场景配对，
      含干扰项）+ 谱系识别 & 立场维度的客观化；
      圈内 / 圈外分组 \\
    D 自我呈现 & 主页人设自评 + 敏感期发布行为回填 &
      H7.3 的劣化内容占比与人设关联 \\
    E 元信息 & 年级、住宿类型、平台、入口来源 &
      配额核查；H7.4 的在场密度分组 \\
    \bottomrule
  \end{tabularx}
\end{table}

\subsection{示例题目}

以下给出各模块的代表题（完整量表施测前须经预试与
信度检验，\conf{M} 级）：

\begin{enumerate}[label=\textbf{Q11.\arabic*},leftmargin=4.2em,itemsep=0.3em]
  \item \textbf{（A，频率分层）}回想本学期的三个时段
        （开学初 / 期中前后 / 期末与考试周），
        你在其中刷到并\keyw{主动转发或发送}奶蛙类内容的频率
        分别是：0 次 / 每月 1--3 次 / 每周 1--2 次 /
        每周 3--5 次 / 几乎每天。
        （同题式样对“其他搞笑内容”重复一遍，作对照基线。）
  \item \textbf{（B，情境勾选）}你发送奶蛙类内容时，
        通常是为了：（可多选）表达自己的失败 / 焦虑 /
        群里有人先发了跟着发 / 纯粹觉得好笑 /
        回应别人的自嘲 / 其他\_\_。
  \item \textbf{（C，场景语法）}“无了”最适合出现的场景是：
        （单选）奶蛙出场 / 奶蛙吃瘪或惊讶 /
        奶蛙加油 / 奶蛙思考。
        （干扰项从表 \ref{tab:06:langs} 的其他场景取。）
  \item \textbf{（D，人设与发布）}考试周里，
        你是否发布过“故意做坏”的内容
        （低画质梗图、摆烂文案、自我丑化图）？
        频率与平时相比：更多 / 持平 / 更少 / 我从不发布。
\end{enumerate}

\subsection{从题目到假设的检验}

\begin{table}[htbp]\centering
  \caption{假设—题目—判别统计量对照（\conf{M} 级设计）}
  \label{tab:11:hypmap}
  \small
  \begin{tabularx}{\textwidth}{@{}lXX@{}}
    \toprule
    假设 & 所需题目 & 判别统计量 \\ \midrule
    H7.1 & Q11.1 两版频率 + B 模块压力分 &
      奶蛙类与总娱乐类的“考试周/平时”频率比之差；
      该差与压力自评的相关 \\
    H7.2 & Q11.2 情境 + （可选）行为日志 &
      情境勾选中“自述”占比的圈内 / 圈外差异 \\
    H7.3 & Q11.4 + D 模块人设自评 &
      敏感期劣化发布的组间（高/低人设）差异 \\
    H7.4 & E 模块住宿类型 + A 模块衰减回填 &
      集体 / 非集体居住组的持续使用时长比较 \\
    FoMO & B 模块 FoMO 量表 &
      FoMO 得分对入会 / 转发行为的增量解释力 \\
    \bottomrule
  \end{tabularx}
\end{table}

注意 Q11.2 的局限：情境勾选测到的仍是
\keyw{自陈动机}，受社会赞许性影响
（“跟着发”比“表达焦虑”更易勾选）。
H7.2 的强检验需要\keyw{行为数据}
（自愿提供的聊天记录片段，经脱敏），
问卷只提供弱证据——这一主次关系写进设计，
防止施测者把自陈当行为。

\section{内容编码方案}\label{sec:11:coding}

第六章命题 \ref{prop:06:channel}（窄通道）与第八章
命题 \ref{prop:08:asym}（成本不对称）需要内容编码。
方案如下（\conf{M} 级设计）：

\begin{enumerate}[label=\arabic*.]
  \item \textbf{样本框}：选定平台、选定时间窗
        （建议覆盖一次期中考试周），以固定检索式
        抓取 / 人工登记奶蛙相关内容，
        记录登记规则使他人可重复；
  \item \textbf{特征清单}：轮廓特征（黄色主体、大笑姿态、
        眼睛形态）与细节特征（配色变化、道具、
        娘化 / 物化元素、背景元素）分列，
        每条内容按清单二值编码；
  \item \textbf{窄通道检验}：比较轮廓特征保有率与
        细节特征保有率的离散系数
        （第六章预测：前者显著更小）；
  \item \textbf{成本不对称检验}：对每条内容另编
        “形式新颖度”（是否改动音轨、结尾套路、时长结构）
        与“内容新颖度”（题材、场景、跨界对象），
        比较两列的多样性指数随时间的变化
        （第八章预测：内容升、形式不升）；
  \item \textbf{质量控制}：双人独立编码，
        报告每项特征的 Krippendorff $\alpha$
        或 Cohen's $\kappa$；
        $\alpha < 0.67$ 的特征项作废重定义，
        不得带病入表；
  \item \textbf{“故意劣化”编码}（服务 H7.3）：
        只编码可观察标记（刻意低分辨率、错位排版、
        拒绝修正的手绘痕迹），
        \keyw{不}编码“看起来做得很差”——
        后者无法与真的做差区分（第七章已警告）。
\end{enumerate}

\section{访谈设计}\label{sec:11:interview}

访谈服务于问卷够不到的三类问题：
机制的\keyw{过程描述}（一次转发是怎么决定的）、
协商的\keyw{多方视角}（第十章遗留问题）、
以及圈内知识的\keyw{习得史}。
设计为半结构化，提纲如下（\conf{M} 级）：

\begin{enumerate}[label=I11.\arabic*,leftmargin=3.2em,itemsep=0.2em]
  \item 请回忆你第一次见到奶蛙的场合。当时你在做什么？
  \item 讲一次你\keyw{没有}转发但差点转发的经历——
        是什么让你停下的？
  \item 你的群聊里，谁是发奶蛙最多的人？
        你怎么看他 / 她？
  \item （立牌活动参与者）从决定投放到立牌被带走，
        哪个时刻你的想法变了？
  \item 你觉得学校管理方怎么看这些立牌？
        他们说得有道理的部分是什么？
  \item 如果明天起所有平台都屏蔽了奶蛙，
        你会做什么？
\end{enumerate}

设计要点：I11.2 用“没转发的瞬间”绕开动机的
事后美化（行为发生在 30 秒内，自陈发生在事后，
两者天然有落差）；I11.5 是第十章“管理方缺席”的
直接回应——请受访者替对方陈述理由，
可部分逼近协商的全景；I11.6 是 M5 的个体版
试探（压制下的行为意向），其结果只能作假设
生成材料，不作检验。

\section{伦理条款}\label{sec:11:ethics}

受体研究的伦理义务独立于、且优先于方法设计：

\begin{enumerate}[label=(\arabic*),leftmargin=2.4em]
  \item \textbf{知情同意}：施测前告知研究目的、
        数据用途、退出权利；未成年人不在总体内，
        若问卷中出现则终止作答；
  \item \textbf{最小化收集}：不收集姓名、学号、
        精确院系等可识别信息；年级用区间，
        院校由受访者自选披露层级；
  \item \textbf{脱敏与销毁}：行为数据（聊天记录片段）
        须去除对话方身份后提交，
        研究结束后按预先声明的期限销毁；
  \item \textbf{访谈伦理}：访谈录音须单独同意；
        涉及校方管理者的叙述时，遵守第四章 E1 规则
        ——以\keyw{角色}而非\keyw{个人}为宾语记录；
  \item \textbf{结果回馈}：向参与者社群（如同好会）
        公开去标识的结果摘要，回报其协作。
\end{enumerate}

\begin{nlmethod}[示例数据的使用纪律]
  为训练分析流程，本章配套示例数据表
  （虚构！\conf{M} 级）：40 个个案、
  六个变量（年级、住宿、考试周频率比、压力分、
  FoMO 分、圈内测验得分），文件位于
  \fname{figures/sample\_data\_example.csv}。
  纪律：(1) 该表数字为本书构造，
  \keyw{任何场合不得引用为其代表真实分布}；
  (2) 用途限于练习 \ref{tab:11:hypmap} 的统计流程；
  (3) 练习结论与真实检验无关。
  示例表的存在本身是一个教学装置：
  让读者在拿到真实数据之前，
  先见过“数据$\to$判别统计量$\to$假设裁决”的完整链路。
\end{nlmethod}

\section{本章小结}\label{sec:11:summary}

本章交付了三套设计：问卷（五模块、示例题、
假设—题目—统计量对照）、内容编码（特征清单、
窄通道与成本不对称的检验步骤、双人编码一致性要求、
“故意劣化”的标记式编码）、访谈（六题提纲、
过程视角与多方视角的设计要点），以及伦理条款五条。
设计的三个支柱——学期节律对齐、不对称指标、
行为—文本—自陈三角——分别是第七章
竞争性解释框架、第一章方法承诺与第四章考据伦理
在受体研究中的落位。至此，全书前十章欠下的
检验债务都有了对应的工具。遗留问题只有一个，
也是最大的一个：\keyw{工具尚未上过战场}。
第十二章在方法论层面收束全书，
其中“可复现性清单”即为这些工具的通用操作规程。

\begin{nlexercises}
  \item Q11.1 用“回想三个时段”回填频率，
        这种回溯式测量有系统性偏差
        （近因效应、情绪一致性回忆）。
        设计一个前瞻式的替代测量
        （提示：固定样本日记法），
        并比较两者在成本与效度上的取舍。
  \item 编码方案第 6 条拒绝编码“看起来做得很差”。
        有人反驳：“故意劣化的辨识恰恰靠整体观感，
        拆成标记就丢了。”请为双方各写一条最强理由，
        然后给出你的裁决与折中方案。
  \item I11.6（“如果平台全部屏蔽”）在设计审查中
        可能被质疑“引导受访者进入受害想象”。
        修改该题的措辞与位置（相对其他题），
        使诱导最小化；若修改后题目失去分析价值，
        说明为什么。
\end{nlexercises}

\nlkeywords 受体研究、抽样框、配额抽样、学期节律对齐、
不对称指标、自陈动机、行为数据、内容编码、
编码者一致性、半结构化访谈、知情同意、最小化收集、
示例数据

\chapter{方法论收束：证据分级、可复现性与研究伦理}\label{ch:12}

\begin{nlkeypoints}
  \begin{enumerate}[label=(\arabic*),leftmargin=2.4em]
    \item 全书的证据分级制度（\conf{A}/\conf{B}/\conf{C}/\conf{M}）
          在本章给出设计理由、操作细则与三条局限；
          它是本书的\keyw{约定}而非行业标准，
          其价值在执行的一致性，不在分级的权威性。
    \item 本章交付“\keyw{可复现性清单}”——
          把散落各章的纪律（口径并列、E1--E3、
          措辞模板、模型声明、编码公开）收拢为
          修订者可直接执行的十项检查
          （第三、四、六章承诺的兑现）。
    \item 第一章遗留的元问题在此收束：
          当研究对象以“解构严肃性”为乐时，
          严肃教材的自洽出路不是假装对象严肃，
          而是\keyw{让体例承担幽默、让方法承担严肃}——
          语域错配因此不是噱头，而是方法立场
          （\ref{sec:02:legitimacy} 的最终答案）。
  \end{enumerate}
\end{nlkeypoints}

全书以个案始（第一章的现象描述），
以方法终。这样安排不是惯例使然，而是因为
奶蛙案例真正可迁移的部分，恰恰不是“奶蛙是什么”
（它是速朽的），而是\keyw{研究一个正在变异、
正在消失、正在被商业与权威觊觎的网络对象时，
研究者应当遵守什么规程}。本章把散落全书的
方法装置收拢为三组：证据分级（第一节）、
可复现性清单（第二节）、伦理三条线（第四节），
并在第三节处理网络材料的时效问题，
在第五节回答那个“教材自身立场”的元问题。

\section{证据分级制度：设计、执行与局限}\label{sec:12:grades}

\subsection{设计理由}

\conf{A}/\conf{B}/\conf{C}/\conf{M} 四级
（第一章 \ref{sec:01:method}）的区分维度是
\keyw{来源的产生方式}，不是内容的真伪：
A 级是职业编辑流程与署名担保的产物，
B 级是机构信誉担保的公开信息，
C 级是社群生产、无编辑流程的材料，
M 级是本书为教学构造的人工物。
用产生方式而非真伪作区分维度，理由有二：
其一，真伪通常无法当场判定，
产生方式却可以当场核验；
其二，\keyw{真伪判断必须由读者带着分级信息自行完成}，
制度不应替读者下结论——这是本书与
“本人已核实一切”式写作的分野。

\subsection{操作细则的五个要点}

第一，\keyw{级别跟随陈述而非来源}：同一来源的
不同陈述可以标不同级别（第三章“上线日期”取
\conf{B} 级倒推即是实例）。第二，\keyw{多口径并列}：
同一指标的不同来源数值并列呈现，不合并、不取均值
（表 \ref{tab:03:numbers}）。第三，\keyw{C 级不得作为
论证前提}：C 级材料可以引出问题、提供线索、
充当被检验对象，但结论不得悬于其上——
第七章 H7.5 的证伪条件设计即遵守此条。第四，
\keyw{M 级必须显式自首}：模型、示例数据、教学命题
在首次出现处声明构造身份，防止教学装置
在二次引用中洗白为经验发现。第五，
\keyw{措辞随级别走}：“据……报道”（A/B）与
“据社群考据”“有说法认为”（C）的句式差异
是分级制度在语言层的执行（各章实例从略）。

\subsection{三条局限}

\begin{enumerate}[label=\arabic*.]
  \item \textbf{级别是粗粒度的}：同为 \conf{C} 级，
        失传媒体维基的词条（有引用、有修订史）与
        无出处的水贴证据地位不同，四级制无法表达
        这一梯度。使用中的补救是\keyw{在正文里
        保留材料自身的限定语}（第四章“不替考据者
        删掉他们自己写的‘但是’”）；
  \item \textbf{分级判断本身可错}：把一则伪装成
        官方通稿的营销稿标为 \conf{B} 级，
        制度无能为力。分级制度降低错误的传导率，
        不保证零错误；
  \item \textbf{制度有成本}：严格的 C 级纪律
        意味着大量“有趣但不可用”的材料被搁置
        ——本书对奶蛙周边的具体价格、销量
        一字不提，就是这一纪律的代价。
        我们认为这个代价值得付，
        但修订者应知道它是\keyw{选择}的结果。
\end{enumerate}

\section{可复现性清单}\label{sec:12:checklist}

以下是散落各章纪律的收拢版，供修订者
在交付新版本前逐项自查（第三章、第四章、第六章
承诺的“操作化版本”）：

\begin{table}[htbp]\centering
  \caption{可复现性清单（本书构造，\conf{M} 级）}
  \label{tab:12:checklist}
  \small
  \begin{tabularx}{\textwidth}{@{}clXl@{}}
    \toprule
    序号 & 检查项 & 操作要点 & 来源 \\ \midrule
    1 & 引用可追溯 & 每条事实有 \conf{A/B} 级出处或明确标注 \conf{C} &
      第一章 \\
    2 & 口径并列 & 同一指标多来源数值并列，不合并 &
      第三章 \\
    3 & C 级措辞 & “据社群考据 / 有说法认为”句式 &
      第一章 \\
    4 & 自我限定保留 & 不删改考据材料自带的“但是” &
      第四章 \\
    5 & E1 对象限定 & 考据叙述宾语限于账号、帖子等公开对象 &
      第四章 \\
    6 & E2 推测隔离 & 个人境况推测“说明存在、不转述内容” &
      第四章 \\
    7 & E3 权利保留 & 二创形象叙述附署名权归属 &
      第四章 \\
    8 & M 级自首 & 模型 / 示例数据首次出现处声明构造身份 &
      第五章 \\
    9 & 方法可复现 & 判定标准、样本框、计数规则写定 &
      第八章 \\
    10 & 编码与量表公开 & 编码手册、问卷、权重随文交付 &
      第十一章 \\ \bottomrule
  \end{tabularx}
\end{table}

清单的用法说明两点。其一，清单是\keyw{拦截装置}：
第二章警告过，语域错配的严肃性不能替代证据的严肃性，
“若某处用 C 级考据支撑了 A 级口吻的结论，
错配就从立场表达滑向误导”——第 3、4 项正是
为拦截这种滑动而设。其二，清单是\keyw{最低标准}：
通过十项不等于研究合格，违反任何一项则肯定不合格。
修订者可以在其上追加本校条款（如增设校内案例的
核对规程），但不得删减既有条目——删减须在新版
前言中显式声明并说明理由。

\section{网络材料的时效：引文腐烂与现存样本}\label{sec:12:rot}

第四章把“不可复核性”处理为个案
（智能体 \fname{zNCvQDKfXy} 的消失），
本节把它一般化。网络引文面临系统性的
\keyw{引文腐烂}（\en{link rot}）：链接失效、
内容被编辑、账号注销、平台下架。
本书的三项应对及适用边界：

\begin{enumerate}[label=\arabic*.]
  \item \textbf{访问日期制度}：所有网络来源记录
        首次访问日期（附录 D），读者以该日期为
        内容状态的基线。边界：访问日期只证明
        “当时如此”，不担保“现在如此”；
  \item \textbf{现存样本立场}：一切基于可检索内容的
        分析，对象都是\keyw{幸存下来的样本}——
        已删除、已注销、未备份的部分永久缺席
        （第十章“样本一端在消失、一端在缺席”）。
        边界：该立场要求结论措辞留有
        “在现存材料中”的限定，本书执行不一贯处，
        以本条为准修订；
  \item \textbf{不可复核性的显式登记}：凡关键证据
        载体已消失的指向（如奶蛙起源），必须在
        结论处写明“不可复核”，不得以交叉印证
        冒充复核。边界：登记不改变指向的
        证明力，只防止证明力被高估。
\end{enumerate}

宏观的网民与短视频使用口径应引用权威统计的
当期报告并注明报告期数\cite{baidu-index}，
避免以过期口径描述现状——这是 \conf{B} 级材料
在时效维度上的延伸纪律。

\section{研究伦理的三条线}\label{sec:12:ethics}

\subsection{对象伦理：E1--E3}

第四章已给出 E1（对象限定）、E2（推测隔离）、
E3（权利保留）的完整表述，此处只补一点
为什么它们必须升格为全书规则而非第四章的
地方性纪律：奶蛙案例中，考据完成与舆论要求
“当事人出面”之间只有一步，任何转述个人境况的
文字——哪怕以同情的形式——都在给这一步施压。
第十一章把 E1--E3 接入访谈设计（以角色为宾语），
第十章的赞诗分析对其只作语体分析、不复制全文，
都是同一规则在不同材料上的执行。

\subsection{版权伦理：分析与背书的界线}

本书把未经授权的形象再生产作为\keyw{分析对象}
（第三、八、九章），附录 D 申明了使用限于
学术评论与教学示例。方法层面补一条操作界线：
\keyw{分析陈述的宾语是结构，不是行为人}——“
换皮生产降低了变异成本”是分析，“你应该去换皮”
是背书，两者的句式区别与 E1 同构。
凡是发现本书某处从结构滑向了怂恿，
按可复现性清单第 1 项的标准，那是缺陷。

\subsection{社群伦理：研究者与社群的互惠}

研究依赖社群产生的材料（考据词条、第一版教材、
同好会记录），却常以“中立观察者”自居。
本书的立场：\keyw{依赖即负债}。可行的偿还方式包括：
引用时忠实署名（C 级材料的来源标注即对
考据者劳动的承认）、研究结果向社群回馈
（第十一章伦理条款第 5 条）、以及不利用
信息差损伤社群（例如不公开尚未被社群公开的
考据线索）。孙立平式的宏观社会心态分析
\cite{sunliping2010} 为“自嘲式表达”提供的
背景框架，也应向受访者开放讨论——
受访者有权知道研究者用何种理论框定他们的笑。

\section{语域错配的最终答案}\label{sec:12:register}

第一章遗留了全书最深的元问题：
当研究对象本身以“解构严肃性”为乐时，
一本严肃教材如何保持自身立场的一致性？
第二章的答案（\ref{sec:02:legitimacy}）是：
命名与体例的语域错配是自觉的方法表达，
但其严肃性不能替代证据的严肃性。
本节给出最终版本，分三点。

第一，\keyw{幽默的安放位置是体例，不是论证}。
章名、框标题、脚注可以玩笑，因为它们不承担
证明功能；定义、数据、推理链必须中性，因为
证明功能不可戏仿。这一分工在本书的执行是：
所有彩色框承担“人声”，所有编号环境
（定义、命题、假设）承担“制度声”。
第二，\keyw{对象的无厘头性是数据，不是许可}。
奶蛙的“意义不明”是第六章的分析对象，
不是本书放松论证标准的理由——恰恰相反，
对象越混沌，方法越要清晰，否则分析
不过是给混沌再添一层。第三，
\keyw{严肃的最终检验是可批判性}：
本书自始至终把可被读者纠正当作目标状态
（第一章：“读者若发现某条 C 级陈述被当作
事实使用，那正是本书鼓励你指出的错误”）。
一本鼓励读者推翻自己的教材，
它的严肃不需要靠板着脸维持。

\begin{nlcritique}[批判性反思：对本书自身的最后检验]
  以自己的标准检验自己：全书是否处处通过了
  表 \ref{tab:12:checklist}？诚实回答是
  \keyw{没有完全通过}。已知的薄弱处：
  其一，第二章“抽象文化”定义标注了 \conf{M} 级，
  但后续章节调用它时偶有接近当作既定概念的表述；
  其二，第十章的协商强度谱三维读数依赖单一来源，
  读数本身应视为“待复核的示例”；
  其三，各章对 \conf{C} 级材料的措辞纪律
  执行不一——个别句子仍用陈述句式引出了
  本应“据社群考据”的内容。这些薄弱处
  不掩盖、不辩解，留作修订第一优先级。
  一本教人识别证据的方法书，
  最不能做的事就是暗示自己没有病灶。
\end{nlcritique}

\section{面向修订者}\label{sec:12:revisers}

本书以 CC BY-NC-SA 4.0 发布（附录 D），
修订与续编是被预期的使用方式。操作规程：

\begin{enumerate}[label=\arabic*.]
  \item \textbf{事实手册先行}：新版修订从更新
        事实手册开始——核对每条 \conf{A/B} 级口径的
        现时状态、重新访问 \conf{C} 级来源并更新日期、
        把新事件补入年表（附录 A）；
  \item \textbf{分级制度保留}：新版可以修正分级、
        补充级别，但不得废除四级制或删除
        表 \ref{tab:12:checklist} 的条目
        （\ref{sec:12:checklist}）；
  \item \textbf{模型重做自由}：SIR 模型、协商强度谱、
        语音距离指标均可替换参数、权重甚至整体架构，
        保留的约束只有一条——“每个构件都声明
        自己的教学身份”；
  \item \textbf{本校化}：补充本校案例时遵守
        第十章的 C 级纪律与 \ref{sec:12:ethics}
        三条线；校本材料（贺图、群聊记录）
        须获当事人同意后方可收录；
  \item \textbf{版本记录}：每次修订在前言后附
        修订说明（改了什么、为什么、依据哪条清单项），
        使读者能追踪本书自身的演变——
        一本关于“考据与版本”的书，
        应当对自己的版本负责。
\end{enumerate}

\section{本章小结}\label{sec:12:summary}

本章把全书的方法装置收拢为三组并交代其边界：
证据分级制度（产生方式为区分维度、五条操作细则、
三条诚实的局限）、可复现性清单（十项拦截装置，
最低标准）、时效纪律（访问日期、现存样本立场、
不可复核性登记）；伦理部分把 E1--E3、
版权界线、社群互惠并为三条线；
最后以三点收束了第一章的元问题——
幽默安放于体例、无厘头性是数据而非许可、
严肃的检验是可批判性——并以自我检验的方式
列出了本书自身的三处病灶。全书在此结束：
十二章走过了一个现象的描述（一、三章）、
考据（四章）、工具（五章）、解剖（六章）、
机制（七章）、生产（八章）、经济（九章）、
空间（十章）、设计的交付（十一章）与方法的自审
（本章）。它没有解决奶蛙为什么好笑——
那要留给读者自己的笑；它解决的是
\keyw{谈论这种好笑时，怎样谈论才算数}。

\begin{nlexercises}
  \item 从本书正文挑出一个你认为\keyw{违反了}
        可复现性清单某一项的具体句子
        （第十二章批判框列出的三处病灶是提示，不是上限），
        写出：清单条目号、原句、修改后的句子、
        修改依据。
  \item 设想你所在的高校同好会决定续编本书并补充
        本校立牌事件。按 \ref{sec:12:revisers} 的规程
        写出续编的前三步操作（含每步的伦理检查点），
        并指出哪一步最容易省略、省略的后果是什么。
  \item “幽默的安放位置是体例，不是论证”
        ——为这句话找一个\keyw{反例}：
        一部（任意领域的）严肃作品，
        其幽默进入了论证本身且因此更有说服力。
        分析该例成立需要什么前提条件，
        并说明这些条件在奶蛙研究中是否具备。
\end{nlexercises}

\nlkeywords 证据分级、可复现性清单、引文腐烂、
现存样本、对象伦理、版权界线、社群互惠、
语域错配、修订规程、自我检验


%% ===================== 附录 =====================
\appendix
%% 附录模式：章标签切换为“附录 X”，节编号切换为“X.1”样式
%% （LaTeX 的 \appendix 会把章计数器清零，正文用的“第\zhnum{chapter}章”
%%   与 \arabic{chapter}.\arabic{section} 在此会错误地显示为“第一章”“1.1”）
\renewcommand{\NLchapterlabel}{附录\,\thechapter}
\renewcommand{\thesection}{\Alph{chapter}.\arabic{section}}
\chapter{大事年表}\label{app:timeline}

本附录汇总本书涉及的关键事件，供快速检索。可信度标注沿用正文制度：
\conf{A} 一手权威来源，\conf{B} 机构或平台公开信息，\conf{C} 社群考据或二手评述。
同一事件存在多个来源口径时并列列出，不作合并。

\section{正典奶龙：IP 与产业线}\label{app:timeline:canon}

\begin{longtable}{@{}p{3.4cm}p{8.6cm}p{2.2cm}@{}}
  \toprule
  时间 & 事件 & 可信度 \\ \midrule
  2017 & 谢承恩与团队扎根深圳，成立第七印象文化传媒（深圳）有限公司 & \conf{A} \\ \midrule
  2019 & 开始构思奶龙形象，定位为“传递治愈、欢乐的积极情感” & \conf{A} \\ \midrule
  2020-01-22 & 奶龙 IP 在短视频平台正式上线 & \conf{B} \\ \midrule
  2020 & 抖音粉丝破千万，单条视频 3 天点播破亿（第一版教材口径，未获独立权威源印证） & \conf{C} \\ \midrule
  2020--2021 & 短视频、有声读物专辑陆续上线；有声读物 8 个月播放量破亿（第一版教材口径） & \conf{C} \\ \midrule
  2023 & 第一季番剧《奶龙搞怪大作战》播出 & \conf{B} \\ \midrule
  2024 & 第二季番剧《奶龙与小七之大战暴暴龙》播出（第一版教材口径） & \conf{C} \\ \midrule
  2024 & 第七印象参评深圳市新型业态企业，获 25 万元扶持奖励；开始参加深圳文博会 & \conf{A} \\ \midrule
  2025-01 & 第三季番剧《奶龙之奇妙冒险》上线，一个月内全网播放量破 1 亿 & \conf{B} \\ \midrule
  2025 底 & 授权合作品牌超 200 家；衍生 SKU 逾 2000 个；动画译制 15 种语言；微信表情包累计发送量突破 80 亿次 & \conf{A/B} \\ \midrule
  2026-01-15 & 《人民日报》报道“一只‘奶龙’背后的创新链”（全网粉丝超 3800 万、流量破千亿、销往 57 个国家和地区） & \conf{A} \\ \midrule
  2026-01-19 至 25 & “奶龙六周年”生日会：线下观影、五城地标大屏（武汉、青岛、天津、成都、苏州）、三城门店见面会 & \conf{B} \\ \midrule
  2026-01-22 & 奶龙六周年；以“歌手”身份发布首支单曲《奶龙的小确幸》 & \conf{B} \\ \midrule
  2026-01-24 & 深圳海岸影城六周年包场观影，约 200 位粉丝观看《奶龙搞怪大作战 2》精选内容 & \conf{B} \\ \midrule
  2026-01-25 & 亮相中央广播电视总台（CMG）《动漫盛典》，以开场节目之一身份出现 & \conf{B} \\ \midrule
  2026-01-28 & 新番《奶龙搞怪大作战 2》在腾讯、优酷、爱奇艺等平台上线，覆盖风行、聚好看、银河少儿等大屏终端；嘉佳卡通、金鹰卡通等电视台陆续播出 & \conf{B} \\ \midrule
  荣誉（未标注具体届次时间） & 亮相昂西国际动画节；获第 21 届中美电视节“年度最佳动画剧集”金天使奖；入围第 30 届亚洲电视大奖“最佳儿童节目” & \conf{B} \\ \bottomrule
\end{longtable}

\section{民间奶蛙：起源与传播线}\label{app:timeline:naiwa}

\begin{longtable}{@{}p{3.4cm}p{8.6cm}p{2.2cm}@{}}
  \toprule
  时间 & 事件 & 可信度 \\ \midrule
  2024-03-07 & 有网友上传将“奶蛙”形象与抖音博主“季洪行”视频进行 AI 融合的内容；“奶蛙”长期被误认为与此有关 & \conf{C} \\ \midrule
  2024-10-02 & 星野 AI 平台用户“纲手【百豪】”创建“奶龙”智能体（ID \fname{zNCvQDKfXy}），平台 AI 生成的头像出现形体畸变，被社群认定为最初版奶蛙形象；该智能体现已无法访问，评论区关闭 & \conf{C} \\ \midrule
  2024-12 前后 & 抖音用户对“欢迎大家恐龙世界”的声讨截图中发现奶蛙照片，并由此出现“星野 AI 生成”的推测（截图含 AI 生成水印） & \conf{C} \\ \midrule
  2026（上半年，具体日期不详） & 抖音视频生成模型 Seedance 推出，奶蛙二创视频生产门槛显著降低（二手评述口径） & \conf{C} \\ \midrule
  2026-03 & 奶龙官方呼吁停止丑化二创，关闭视频评论区并批量拉黑发布奶蛙的账号 & \conf{C} \\ \midrule
  2026-08-04 & B 站用户“雾蚀心忘症”在星野 App 中找到符合形象的奶蛙智能体并发布评论，起源考据取得关键突破 & \conf{C} \\ \midrule
  2026-08 & 以 B 站 UP 主“看仙碎”为代表的一批爱好者推进起源考据，发布考据视频（视频号 BV1ndMk65EwM，第一版教材口径） & \conf{C} \\ \midrule
  2026 全年 & 奶蛙成为中文互联网高传播度迷因；被二手评述称为“2026 年传播最广、影响最深、二创最活跃的 meme” & \conf{C} \\ \bottomrule
\end{longtable}

\section{高校与线下活动线}\label{app:timeline:campus}

\begin{longtable}{@{}p{3.4cm}p{8.6cm}p{2.2cm}@{}}
  \toprule
  时间 & 事件 & 可信度 \\ \midrule
  2026-06 & 哈尔滨工业大学 2026 年毕业典礼暨学位授予仪式出现第一个广为人知的奶蛙立牌 & \conf{C} \\ \midrule
  2026-06 & 上海交通大学发起“流动奶蛙”活动（最早的流动奶蛙至今仍在用） & \conf{C} \\ \midrule
  2026-06 & 中国海洋大学“奶 OUC”四奶蛙立牌出现（第一个广为人知的复数奶蛙立牌） & \conf{C} \\ \midrule
  2026-06-24 & 北京航空航天大学奶龙同好会成立 & \conf{C} \\ \midrule
  2026-08 初 & 北航同好会人数突破 100 人 & \conf{C} \\ \midrule
  2026-09-28 & 北京航空航天大学学院路校区绿园投放流动奶蛙立牌；旁边随后出现丰川祥子立牌；当晚被保安带走，存活约 9 小时 & \conf{C} \\ \midrule
  2026-10 初 & 北航同好会人数突破 200 人（26 级新生入学后） & \conf{C} \\ \midrule
  2026（具体日期不详） & 浙江大学某地招生宣传立牌出现浙大奶蛙形象 & \conf{C} \\ \bottomrule
\end{longtable}

\noindent\srcnote{年表中“具体日期不详”的条目，
来源为第一版教材或社群考据词条，未给出确切日期；
本书不推测日期。2020 年“抖音粉丝破千万”等
未获独立权威源印证的条目，引用时请保留 \conf{C} 标注。}

\chapter{术语表}\label{app:glossary}

本附录汇总全书使用的关键术语，按拼音字母排序。
释义为检索用的压缩版，不替代正文论证；
“详见”列给出术语首次获得完整定义或分析的章节。
凡属本书自建的教学构件，沿用正文的 \conf{M} 级标注。
英文对照均随词条给出，仅为主要文献用词，不作翻译仲裁。

\begin{longtable}{@{}p{3.0cm}p{9.2cm}p{1.8cm}@{}}
  \toprule
  术语 & 释义 & 详见 \\ \midrule
  \endhead
  \multicolumn{3}{@{}l}{\textbf{A}} \\ \midrule
  AIGC（\en{AI-Generated Content}）&
    人工智能生成内容。第八章把 2026 年奶蛙视频二创的爆发
    归因于以 Seedance 为代表的视频生成模型把生产门槛降到接近零，
    并据此提出“成本不对称”命题。 &
    第八章 \\ \midrule
  \multicolumn{3}{@{}l}{\textbf{B}} \\ \midrule
  半结构化访谈 &
    介于固定问卷与自由谈话之间的访谈形式：按提纲推进、
    不限于既定顺序与措辞。第十一章交付的六题提纲即属此类。 &
    第十一章 \\ \midrule
  编码者一致性 &
    多名编码者对同一材料独立编码时结果的一致程度，
    以 Krippendorff $\alpha$ 或 Cohen's $\kappa$ 报告；
    $\alpha<0.67$ 的特征项应作废重定义，不得带病入表。 &
    第十一章 \\ \midrule
  变异引擎 &
    本书对 AIGC 工具在生产链中所承担环节的称呼：
    以近零成本、高频率产生形象变体，供给第六章所述的
    “内容软”维度。 &
    第八章 \\ \midrule
  变异权 &
    社群对形象进行再创作与变异的实际自由度。
    与法律上的改编权不同：它存在于意义层面，
    不因个体账号被封禁而消失。 &
    第八章 \\ \midrule
  不对称指标 &
    判别设计：功能机制应\keyw{不对称地}改变某指标，
    而消遣 / 算法等竞争解释对其\keyw{对称地}改变；
    两侧表现之差即机制的证据空间。 &
    第七章 \\ \midrule
  不可复核性 &
    关键证据载体已消失或关闭、任何第三方无法再检验的状态。
    奶蛙起源指向的智能体即典型；本书要求在结论处显式登记之。 &
    第四章 \\ \midrule
  \multicolumn{3}{@{}l}{\textbf{C}} \\ \midrule
  采用者类别 &
    Rogers 创新扩散理论中创新者、早期采用者、
    早期 / 后期大多数等类别；第十章用于五校案例的
    形态归位（\conf{M} 级结构假设）。 &
    第十章 \\ \midrule
  参与式文化 &
    受众积极参与生产、分享与改写内容的文化形态；
    本书以之自我定位：教材既是研究也是参与。 &
    第一章 \\ \midrule
  场景语法 &
    衍生语的使用受场景约束的规则系统（“无了”配吃瘪、
    “糖”配反讽）；配错场景即被识别为外行。 &
    第六章 \\ \midrule
  成本不对称 &
    第八章命题：AIGC 使改\keyw{内容}成本近零、
    改\keyw{形式}（音轨、结尾套路）成本仍高，
    由此产生内容多样、形式固化的分化。 &
    第八章 \\ \midrule
  抽象文化 &
    本书工作定义（\conf{M} 级）：以意义悬置、空耳、
    语体戏仿为可观察标记的中文网络风格类型；
    严肃语体承载无厘头内容是其中心装置。 &
    第二章 \\ \midrule
  抽样框 &
    可从中抽取样本的名单或程序。奶蛙受体研究的总体
    不可枚举，故以“配额 + 多入口”替代完整抽样框。 &
    第十一章 \\ \midrule
  错失恐惧（FoMO） &
    对他人正在经历自己缺席的有益经验的普遍性忧虑；
    第十章用于解释新生入学期的同好会纳新增长。 &
    第十章 \\ \midrule
  \multicolumn{3}{@{}l}{\textbf{D}} \\ \midrule
  道德恐慌（\en{moral panic}） &
    Cohen 术语：媒体与权威把低危害的风格现象
    建构为对秩序的威胁的过程；第十章用于警示
    对校园玩梗的运动式清理。 &
    第十章 \\ \midrule
  倒放大法 &
    社群加工手法：把正典音轨倒放，使“我是奶龙，我才是奶龙”
    变为空耳文本；属“捧腹大笑”声音符号的两条加工线之一。 &
    第六章 \\ \midrule
  低强度协商 &
    对立双方以可撤回、不对抗的动作（移除对纪念）
    表达立场并重申边界的互动形态；
    “低强度协商不是冲突的失败，而是冲突的替代”。 &
    第十章 \\ \midrule
  \multicolumn{3}{@{}l}{\textbf{E}} \\ \midrule
  E1/E2/E3 &
    考据伦理三规则：E1 对象限定（只叙述账号、帖子等
    公开对象）；E2 推测隔离（个人境况只说明存在、
    不转述内容）；E3 权利保留（二创形象叙述附署名权归属）。 &
    第四章 \\ \midrule
  二次爆发 &
    SIR 教学模型中新变体输入项触发的第二条感染峰；
    第五章以双峰曲线示意 2024 与 2026 两波热度（\conf{M} 级）。 &
    第五章 \\ \midrule
  \multicolumn{3}{@{}l}{\textbf{F}} \\ \midrule
  反精致 &
    第七章 M3：以故意做差、画质劣化对抗精致内容的姿态；
    其判别问题是“反主流的规范”还是“另一种人设消费”。 &
    第七章 \\ \midrule
  分层装配模型 &
    本书教学综合（\conf{M} 级）：压力层（M1、M2）、
    风格层（M3）与场景层（M4、M5）叠加装配出单一笑点；
    层间可独立证伪。 &
    第七章 \\ \midrule
  复制子（\en{replicator}） &
    Dawkins 术语：能复制自身的单位。第五章把“梗”处理为
    文化复制子，其存续取决于保真度、多产性与存续性。 &
    第五章 \\ \midrule
  \multicolumn{3}{@{}l}{\textbf{G}} \\ \midrule
  关评拉黑 &
    官方奶龙 2026-03 关闭视频评论区并批量拉黑奶蛙账号的
    治理动作（\conf{C} 级）；第九章解读为
    “只能作用于法律层面可执行对象”的不彻底抑制。 &
    第三、九章 \\ \midrule
  \multicolumn{3}{@{}l}{\textbf{H}} \\ \midrule
  换皮奶蛙 / 异形奶蛙 &
    形象谱系两大类：保留奶蛙形体轮廓、仅换配饰与场景者
    为换皮；形体本身偏离（简化、娘化、物化、跨物种）者为异形。 &
    第六章 \\ \midrule
  \multicolumn{3}{@{}l}{\textbf{J}} \\ \midrule
  机械性（\en{mechanical encrustment}） &
    Bergson 笑论：可笑源于“生命的机械化”——
    活人做出僵硬机器般的程式动作；奶蛙的重复仪式形态
    为其当代样本。 &
    第七章 \\ \midrule
  机制矩阵 &
    第七章交付的四列模板（压力源—补偿行为—
    可观察指标—竞争性解释），供书写任何机制假设使用
    （\conf{M} 级）。 &
    第七章 \\ \midrule
  降格（\en{tendentious wit}） &
    Freud 诙谐理论：笑话经由压缩与凝缩绕过心理抑制，
    释放被压抑的攻击或性冲动能量。 &
    第七章 \\ \midrule
  竞争性解释 &
    与目标机制竞争的替代解释（纯消遣、算法推荐、单纯从众）。
    本书纪律：每个机制假设必须并列写出竞争解释与判别设计。 &
    第七章 \\ \midrule
  \multicolumn{3}{@{}l}{\textbf{K}} \\ \midrule
  可传播性（\en{spreadability}） &
    Jenkins 等人术语：内容被受众主动携带、复制、改写的潜力；
    内容的价值边界由受众的可传播性行为划定。 &
    第九章 \\ \midrule
  可传播性单位 &
    本书教学命题：搬运中不依赖原始语境（内容可换、
    署名可弃、载体可换）的文化构件；
    “结尾集体大笑”为奶蛙语料之最（\conf{M} 级）。 &
    第九章 \\ \midrule
  可复现性清单 &
    第十二章交付的十项自查表（引用可追溯、口径并列、
    C 级措辞、E1--E3、M 级自首等）；
    修订者交付新版本前的最低标准（\conf{M} 级）。 &
    第十二章 \\ \midrule
  空耳（\en{soramimi}） &
    把一种语言的语音听成另一种词句的再诠释实践。
    第六章作为社群知识分析，第八章以“声韵母编辑距离”
    为其操作化（\conf{M} 级）。 &
    第六、八章 \\ \midrule
  口径并列 &
    本书证据纪律：同一指标的不同来源数值并列呈现、
    注明来源与日期，不合并、不取均值
    （如 3800 万与 3900 万并存）。 &
    第三章 \\ \midrule
  \multicolumn{3}{@{}l}{\textbf{L}} \\ \midrule
  立牌 &
    奶蛙符号在物理空间的锚点物：把线上迷因固定为
    可合影、可搬运、可巡游的实物，使符号实践进入
    物理空间并接受校方管理。 &
    第十章 \\ \midrule
  流动奶蛙 &
    上海交通大学发起的空间实践：把立牌带到校园各处
    制造偶遇事件。其合影率高于固定装置为可检验推论。 &
    第十章 \\ \midrule
  \multicolumn{3}{@{}l}{\textbf{M}} \\ \midrule
  meme（迷因） &
    Dawkins：通过模仿传播的文化复制子；
    Shifman 将其操作化为内容 / 形式 / 立场三维。 &
    第二、五章 \\ \midrule
  免疫衰减 &
    SIR 教学模型中恢复者随时间回到易感者状态的通道，
    解释同一梗的复发；与二次爆发配套（\conf{M} 级）。 &
    第五章 \\ \midrule
  迷因簇 &
    第五章定义：共享一组可辨识特征的一群变体；
    归并标准 (C1) 特征集、(C2) 谱系、(C3) 识别。 &
    第五章 \\ \midrule
  迷因生态缝隙 &
    控制工具（作用于法律层面）与控制对象
    （部分转移到意义层面）的错位空间；
    “失败的控制不是控制不努力，而是层面错配”。 &
    第九章 \\ \midrule
  模板固化 &
    第八章观察性命题（\conf{M} 级）：AIGC 二创
    在题材上丰富化的同时，形式收敛于少数模板
    （音轨、结尾套路、时长结构）。 &
    第八章 \\ \midrule
  \multicolumn{3}{@{}l}{\textbf{N}} \\ \midrule
  奶蛋 &
    异形奶蛙中简化到极致的产物：形体特征只剩
    “眼睛 + 大笑”，标定谱系的简化极限；
    因“奶蛋轮椅”叙事被赋予强运人设。 &
    第六章 \\ \midrule
  奶龙 &
    第七印象文化传媒 2019 年起开发的动画 IP（正典），
    2020-01-22 于短视频平台上线；
    产业口径见大事年表。 &
    第三章 \\ \midrule
  奶蛙 &
    奶龙二创中诞生的黄色大笑蛙形迷因形象；
    起源指向星野平台某智能体（\conf{C} 级，不可复核）。 &
    第四章 \\ \midrule
  能量节省 &
    Spencer 以来功能主义笑论：笑是心理能量
    经意外通路释放的“节省装置”；第七章古典理论之一。 &
    第七章 \\ \midrule
  逆反认证 &
    第七章 M5：权利人打压反而为迷因提供“官方认证”；
    其判别设计与治理两难分析见第七章。 &
    第七章 \\ \midrule
  \multicolumn{3}{@{}l}{\textbf{P}} \\ \midrule
  判定标准 &
    第八章纪律：任何“是 / 否奶蛙内容”的编码必须有
    事先写定的可执行规则（样本框、检索式、计数规则），
    使他人可重复。 &
    第八章 \\ \midrule
  配额抽样 &
    按关键亚组（年级、院校层次）设下限的非概率抽样；
    无完整抽样框时的替代方案，毕业年级为关键对比组。 &
    第十一章 \\ \midrule
  \multicolumn{3}{@{}l}{\textbf{Q}} \\ \midrule
  情绪代偿 &
    第七章 M1：以“失败—大笑”结构为压力情绪提供
    即时替代性出口；奶蛙语料的中心机制假设（H7.1）。 &
    第七章 \\ \midrule
  群体稀释 &
    第七章 M2：转发把个体压力“稀释”进群体，
    换取“不止我一个人这样”的归属信号；
    与社交笑研究同源。 &
    第七章 \\ \midrule
  \multicolumn{3}{@{}l}{\textbf{R}} \\ \midrule
  $R_0$（基本再生数） &
    SIR 模型中单个感染者在易感者群体中造成的平均新感染数；
    $R_0>1$ 则梗扩散。本书取值为教学示意，非经验估计。 &
    第五章 \\ \midrule
  \multicolumn{3}{@{}l}{\textbf{S}} \\ \midrule
  SIR 型教学模型 &
    第五章把易感—感染—移出流行病模型改造为
    迷因传播的教学模型（\conf{M} 级，非经验模型）。 &
    第五章 \\ \midrule
  三维框架 &
    Shifman 对迷因的操作化：内容（说什么）、
    形式（怎么说）、立场（对谁说、以何姿态说）；
    本书符号分析的工作框架。 &
    第五、六章 \\ \midrule
  社交笑（\en{social laugh}） &
    Provine 笑研究：笑首先是社会性信号
    （群居情境频率约为独处 30 倍），而非仅对幽默的反应。 &
    第七章 \\ \midrule
  社群互惠 &
    第十二章伦理第三线：研究依赖社群材料即负债，
    应以忠实署名、结果回馈、不利用信息差偿还。 &
    第十二章 \\ \midrule
  生产链 &
    第八章整理的奶蛙内容生产环节序列
    （梗图模板 $\to$ 视频剪辑 $\to$ AIGC 生成）；
    各环节成本结构不同，解释变异供给的来源。 &
    第八章 \\ \midrule
  声韵母编辑距离 &
    第八章定义（\conf{M} 级）：按声母 / 韵母 / 声调
    计算两个短语的编辑代价并归一化得 $d$，
    作为空耳成立度的可计算代理指标。 &
    第八章 \\ \midrule
  示例数据 &
    第十一章配套的虚构 40 例 CSV 文件（\conf{M} 级），
    供练习统计流程；任何场合不得引用为代表真实分布。 &
    第十一章 \\ \midrule
  收编（\en{incorporation}） &
    Hebdige：主流与商业把亚文化风格转化为商品、
    消解其抵抗性的过程；第九章指出奶蛙案例中
    收编是双向运动，并提醒先确认抵抗性证据。 &
    第九章 \\ \midrule
  受众外溢 &
    第三章：奶龙受众从亲子向学生群体扩展的三层解释
    （用户年龄结构、二创作者构成、观看场景迁移）。 &
    第三章 \\ \midrule
  双轨分发 &
    正典内容由制作方控制渠道分发，变种内容由
    二创者网络搬运流通；两条轨道的媒介与经济结构均不同。 &
    第二、三章 \\ \midrule
  \multicolumn{3}{@{}l}{\textbf{T}} \\ \midrule
  糖 &
    社群评级词（反讽用法）：指正典内容中幼稚、
    尬感强的特质；社群以讽刺性彩虹屁完成反向评级。 &
    第六章 \\ \midrule
  同好会 &
    校园迷因社群的半正式组织（如北航奶龙同好会）：
    名义功能是水群，实际功能是组织纳新、贺图、
    表情包生产与立牌活动。 &
    第十章 \\ \midrule
  \multicolumn{3}{@{}l}{\textbf{W}} \\ \midrule
  外部性 &
    第八章：二创的生产成本由个体承担、
    拥挤成本（低质内容刷屏）由公共信息流承担；
    平台侧治理机制设计的分析基础。 &
    第八章 \\ \midrule
  未结案证据链 &
    本书对奶蛙起源的定性：指向存在、时间线连贯，
    但关键载体不可复核；只能作教学示范，不作已证事实。 &
    第四章 \\ \midrule
  文本盗猎（\en{textual poaching}） &
    Jenkins（承 de Certeau）：粉丝对文本的
    选择性挪用与再创作；第二章用于定位二创的意义。 &
    第二章 \\ \midrule
  文化资本 &
    Bourdieu 术语：圈内知识（场景语法、谱系识别）
    的持有量构成圈内 / 圈外的边界，是准入测试的给分依据。 &
    第二、五章 \\ \midrule
  误指 &
    考据阶段术语：把错误来源广泛指认为起源
    （如“奶蛙来自某博主”）；误指是考据的起点而非障碍。 &
    第四章 \\ \midrule
  \multicolumn{3}{@{}l}{\textbf{X}} \\ \midrule
  现存样本偏差 &
    以可检索内容为对象的研究只能看到幸存样本
    （未删除、未注销、被记录者）；结论措辞应带
    “在现存材料中”的限定。 &
    第十、十二章 \\ \midrule
  校本化 &
    全国性梗被转写为本校身份标记（“数分奶蛙”
    “航概奶蛙”）；持有并使用即“本校圈内人”的默示声明。 &
    第十章 \\ \midrule
  协商强度谱 &
    第十章教学模型（\conf{M} 级）：按校方反应 /
    学生反应 / 持久物三维给立牌事件分级（0/1/2）；
    北航 9 小时事件合计 2，属低—中强度区间。 &
    第十章 \\ \midrule
  学期节律对齐 &
    第十一章设计支柱：把测量窗口对齐期中与期末考试周，
    使压力机制的效应在数据中可见。 &
    第十一章 \\ \midrule
  行为数据 &
    自愿提供的脱敏聊天记录片段等实际行为痕迹；
    对机制假设的检验优先于自陈问卷。 &
    第十一章 \\ \midrule
  \multicolumn{3}{@{}l}{\textbf{Y}} \\ \midrule
  衍生语 &
    社群从语料中沉淀的固定用语（“无了”“安迪”等），
    受场景语法约束，兼作圈内识别与身份标记。 &
    第六章 \\ \midrule
  引文腐烂（\en{link rot}） &
    网络引文的系统性失效（链接失效、内容被编辑、
    账号注销）；本书三项应对：访问日期、
    现存样本立场、不可复核性登记。 &
    第十二章 \\ \midrule
  一软两硬 &
    第六章对奶蛙符号系统的结构性判断（\conf{M} 级）：
    内容高替换（软）、形式高稳定（硬）、
    立场高约定化（硬）——用最硬的形式承载最软的内容。 &
    第六章 \\ \midrule
  阈限（\en{liminality}） &
    Turner：脱离日常结构、处于角色之间的过渡状态；
    “流动奶蛙合影”被读作微阈限事件（\conf{M} 级诠释）。 &
    第十章 \\ \midrule
  语域错配 &
    本书的命名与体例策略：严肃语域承载无厘头对象；
    其严肃性不能替代证据的严肃性，最终答案见第十二章。 &
    第一、二章 \\ \midrule
  \multicolumn{3}{@{}l}{\textbf{Z}} \\ \midrule
  窄通道规律 &
    第六章命题：变异集中于少数“通道”（道具、配饰、题材），
    形体轮廓等硬核极少被动摇；第八章以成本不对称
    解释其生产端成因。 &
    第六、八章 \\ \midrule
  正典（\en{canon}） &
    权利主体认定的官方设定与内容（正典奶龙）；
    与变种的区分在法律层面与意义层面上并不重合。 &
    第二、三章 \\ \midrule
  证据分级（A/B/C/M） &
    本书制度：A 一手权威、B 机构公开、C 社群考据
    （不作论证前提）、M 本书自建教学构件（须显式声明）；
    区分维度是来源的产生方式，不是内容的真伪。 &
    第一、十二章 \\ \midrule
  证伪条件 &
    第七章纪律：每个机制假设写出“什么观察会使我放弃它”；
    H7.5（逆反认证）的证伪条件设计为范例。 &
    第七章 \\ \midrule
  知情同意 &
    受体研究伦理：施测 / 访谈前告知目的、用途、
    退出权并取得同意；未成年人不在总体内。 &
    第十一章 \\ \midrule
  自陈动机 &
    问卷测得的受访者自述原因；受社会赞许性影响，
    只作弱证据，强检验需行为数据。 &
    第十一章 \\ \midrule
  自我检验 &
    第十二章：本书以自身标准检验自身并列出三处病灶，
    不掩盖、不辩解；方法书最不能做的是
    暗示自己没有病灶。 &
    第十二章 \\ \midrule
  最小化收集 &
    伦理原则：不收集姓名、学号等可识别信息；
    年级用区间，院校由受访者自选披露层级。 &
    第十一章 \\ \bottomrule
\end{longtable}

\chapter{习题参考答案}\label{app:answers}

本书习题多为开放题：答案不止一个，评分对象是
\keyw{方法与论证}，不是结论本身。本附录给出各题的
\keyw{答题要点}与\keyw{评分点}，供自答后对照、
或供施教者批改参考。凡要求读者自选实例的题目，
答案只给判定标准；凡有计算成分的题目，给出完整过程。
建议次序：先自答，再对照；若你的答案与本附录冲突，
先检查方法，再检查结论——冲突本身可能就是
一道新的习题。

\section{第一章}\label{app:ans:01}

\begin{enumerate}[label=\textbf{1.},leftmargin=2.6em]
  \item \textbf{要点}：三维各写一句——内容（它说了什么、
        主角是谁、发生了什么事）；形式（怎么呈现：模板、
        音轨、排版、结尾套路）；立场（对谁说、以什么姿态说：
        自嘲 / 圈内认证 / 反讽评级）。“硬约定”维的判别标志：
        \keyw{用错会被社群识别为外行}（如衍生语配错场景、
        缺失结尾集体大笑）。多数梗的硬维在形式或立场，
        内容维可自由发挥——但归属取决于所选梗。
        \textbf{评分点}：能否举出“用错会被识别”的具体证据，
        而非硬维的归属结论本身。
\end{enumerate}

\begin{enumerate}[label=\textbf{2.},leftmargin=2.6em]
  \item \textbf{要点}：压力源假设指标示例——同一批账号在
        考试周内奶蛙类转发占其全部转发的比例，相对平时周的
        上升幅度（单一比例，可计数）。算法假设指标示例——
        同期话题曝光总量（平台播放量、话题页浏览）的增幅。
        若两者方向相反（曝光大增而个体转发占比不动）：
        总热度由供给侧与推荐侧驱动，使用侧的压力选择机制
        未获得支持——\keyw{竞争性解释获得相对支持}。
        \textbf{评分点}：指标是单一可计数数字；能说出
        “方向相反不是证伪，而是相对支持”这一逻辑关系。
\end{enumerate}

\begin{enumerate}[label=\textbf{3.},leftmargin=2.6em]
  \item \textbf{要点}：如全网粉丝 3800 万（《人民日报》口径）
        与 3900 万（另一来源口径）。并列而非取均值的理由：
        两数字的统计时点、口径边界（是否含海外、是否含子账号）
        可能不同；均值会制造出\keyw{任何一方都不曾主张的假口径}。
        对引用者的影响：引用时必须携带“来源 + 日期”，
        不得只引数字；也不得声称“约 3850 万”——
        并列制度把不确定性显式交还给引用者。
        \textbf{评分点}：能说出“均值 = 新造口径”这一核心。
\end{enumerate}

\section{第二章}\label{app:ans:02}

\begin{enumerate}[label=\textbf{1.},leftmargin=2.6em]
  \item \textbf{要点}：法律层面实例——官方关闭评论区、
        批量拉黑（\conf{C} 级）、法务团队与快速维权机制
        （\conf{A} 级）。意义层面实例——二创者反讽征用正典
        桥段（“糖”评级）、“潘多拉奶盒”叙事宣布脱离正典。
        两者不总是一致：法律层面的支配关系持续存在，
        意义层面的流通已不服从该边界（第九章命题
        \ref{prop:09:dual}）。
        \textbf{评分点}：两实例分属两层面；能指出
        “不一致正是迷因生态缝隙”。
\end{enumerate}

\begin{enumerate}[label=\textbf{2.},leftmargin=2.6em]
  \item \textbf{要点}：最易贴标签的通常是“文本盗猎”
        （答“抽象文化”且自圆亦可）。以文本盗猎为例，
        使用前必答的三个结构问题：(1) 盗猎\keyw{什么}——
        被挪用的文本单元（形象 / 桥段 / 音轨）；(2) \keyw{怎么}盗猎
        ——再创作的技术路径与门槛；(3) 盗猎后\keyw{功能}的变化
        ——意义、受众、流通结构各变成什么。
        \textbf{评分点}：三问缺一即视为贴标签。
\end{enumerate}

\begin{enumerate}[label=\textbf{3.},leftmargin=2.6em]
  \item \textbf{要点}：按第二章定义的三项标记逐条判定。
        (a) 讣告体作业封面：命中语体戏仿（庄重语体承载
        校园小事）+ 意义悬置（框架与对象严重错配），
        判“属于”。(b) 画质劣化但语义完整：意义未悬置；
        能否命中取决于劣化是否构成对“精致规范”的
        可辨认姿态——判“边缘”，关键看能否指认戏仿对象。
        (c) 仅画质劣化、语义清晰：三项皆不命中，
        判“不属于”。
        \textbf{评分点}：逐条对标记作答；禁止以
        “感觉很抽象”代替判定。
\end{enumerate}

\section{第三章}\label{app:ans:03}

\begin{enumerate}[label=\textbf{1.},leftmargin=2.6em]
  \item \textbf{要点}：第三季 2025-01 上线与“一个月破亿”
        有机构报道支撑（\conf{B}）；第二季播出信息仅见于
        第一版教材口径（\conf{C}）。同系列不同年份可信度
        不同，体现“\keyw{级别跟随陈述}”：等级属于
        具体陈述（某季某时点的报道是否存在），
        不属于品牌整体。
        \textbf{评分点}：能复述“级别跟随陈述而非来源”。
\end{enumerate}

\begin{enumerate}[label=\textbf{2.},leftmargin=2.6em]
  \item \textbf{要点}：层 1（用户年龄）——平台年报中
        奶龙内容 20 岁以上观众占比的时间序列；层 2
        （二创作者构成）——抽样二创作者并检索其公开简介中
        高校关键词的计数占比；层 3（场景迁移）——门店
        非亲子时段客群观察、校园立牌出现频率等自建观察。
        最难测的是层 3：场景是弥散的线下行为，
        无现成平台数据，任何代理指标都离行为更远。
        \textbf{评分点}：承认所有指标都是代理；
        层 3 指标不得依赖自陈。
\end{enumerate}

\begin{enumerate}[label=\textbf{3.},leftmargin=2.6em]
  \item \textbf{要点}：不矛盾。法务 + 维权（\conf{A}）
        是授权体系的\keyw{前置条件}（第九章）；
        关评拉黑（\conf{C}）针对“丑化二创”这一失控流通。
        两者共同服务于“形象可租赁”。依赖 C 级的部分：
        关评拉黑的时间、规模、动机细节；只依赖 A 级的部分：
        维权机制的存在及其产业逻辑。
        \textbf{评分点}：能把结论拆开，分别交代证据级别。
\end{enumerate}

\section{第四章}\label{app:ans:04}

\begin{enumerate}[label=\textbf{1.},leftmargin=2.6em]
  \item \textbf{要点}：第一句是断言形态（要求起源已证：
        载体可复核、证据链闭合）；第二句是考据形态
        （指向 + 限定：指向存在，载体不可复核须写明）。
        第一句当前不成立：关键智能体已无法访问，
        任何第三方无法复核，材料只支持“指向”，
        不支持“来自”。
        \textbf{评分点}：能说出“不可复核”在此的作用位置。
\end{enumerate}

\begin{enumerate}[label=\textbf{2.},leftmargin=2.6em]
  \item \textbf{要点}：学生自选实例。三阶段划分须正确：
        误指（被广泛指认的错误来源）、线索（推动定位的
        可观察痕迹）、发现（可复核的载体）。若该梗起源
        已公认，须指出跳过的阶段及原因（如起源系官方公布，
        则“误指”阶段无生长空间）。
        \textbf{评分点}：阶段划分与该梗实际考据史相符。
\end{enumerate}

\begin{enumerate}[label=\textbf{3.},leftmargin=2.6em]
  \item \textbf{要点}：检查程序示例。第一步：圈出修订稿中
        一切涉及具体个人的句子（关键词：本人、作者、
        当事人、其账号等）。第二步：逐句套用句式模板——
        允许：“社群中存在关于某人的说法”“该说法无法证实”；
        禁止：转述具体内容的从句（境况、去向、私生活细节）。
        第三步：核对所有“说明存在”句是否只含出处与
        存在声明。违例句数为零则通过。
        \textbf{评分点}：程序能在五分钟内执行
        （关键词清单 + 句式模板两件套齐全）。
\end{enumerate}

\section{第五章}\label{app:ans:05}

\begin{enumerate}[label=\textbf{1.},leftmargin=2.6em]
  \item \textbf{要点}：(a) 高数课本表情包：特征集保留（C1）、
        谱系同源（C2）、可识别（C3），同簇；决定性标准 C1。
        (b) 黄色大笑正常体型龙：C1 部分命中（黄、大笑），
        但若无畸变谱系出身，C2 失败，不同簇；决定性标准 C2。
        (c) 异形象同音轨：C1 失败、C3 依赖场景；若社群在
        奶蛙语境中稳定使用，可划入“形式衍生”边缘。
        \textbf{评分点}：本题意图在让读者发现三条标准
        可能给出不一致判定——\keyw{承认模糊比强行裁决
        更符合第五章立场}，能写出这一层为优。
\end{enumerate}

\begin{enumerate}[label=\textbf{2.},leftmargin=2.6em]
  \item \textbf{要点}：按第五章记号，感染变化率
        $\mathrm{d}I/\mathrm{d}t = I \times
        (\text{接触项} - \text{恢复项})$。$\omega = 0$ 时
        $S(t)$ 单调不增，故接触项随时间不增；恢复项不增减
        （或单调）。令 $\mathrm{d}I/\mathrm{d}t = 0$：
        等式至多在一点成立（或成立为一个区间，对应平台期）。
        因此 $I(t)$ 先增、后（至多）减，至多一个极大值。
        若 $S(t)$ 始终高于阈值，则 $I$ 单调增——峰在
        观测窗之外，这正是命题中“\keyw{至多}”一词的含义。
        \textbf{评分点}：用到“$S$ 单调”与“驻点条件”
        两个要素；“至多”的边界情形。
\end{enumerate}

\begin{enumerate}[label=\textbf{3.},leftmargin=2.6em]
  \item \textbf{要点}：判据示例——先作 7 日平滑；
        峰 = 局部极大且两侧回落超过峰值的三成；
        双峰 = 两峰间隔不小于两周且谷值低于两峰均值六成。
        稳健化：对平滑窗宽与谷值比例做扰动，形态不变者
        才报告。不能识别的：无法区分“同一梗的二次爆发”
        与“两个独立梗的叠加”；无法识别免疫衰减参数本身；
        无法排除平台限流造成的假谷。
        \textbf{评分点}：判据可执行 + 至少一条
        “本方案不能识别什么”。
\end{enumerate}

\section{第六章}\label{app:ans:06}

\begin{enumerate}[label=\textbf{1.},leftmargin=2.6em]
  \item \textbf{要点}：逐条对照 (C1)(C2)(C3)，指出
        起决定作用的一条。主张移出某变种者，必须指认
        其丢失了哪个“不可再丢”的特征（如只剩黄色、
        失去大笑与谱系）。无标准答案。
        \textbf{评分点}：判定过程可复核（三条标准逐一给出
        满足 / 不满足及证据）。
\end{enumerate}

\begin{enumerate}[label=\textbf{2.},leftmargin=2.6em]
  \item \textbf{要点}：以“典”为例——判定操作：
        复读已成套路的论证模板（可观察特征：句式雷同、
        引用链封闭）。被官方或广告征用后的变化：
        评级权从社群转移到被评对象，词义“通货膨胀”
        （什么都典 = 无物可典），社群或弃用、或转义为
        更深的反讽。
        \textbf{评分点}：把“效力变化”操作化
        （使用场景的收敛或迁移），不停留在“变味了”。
\end{enumerate}

\begin{enumerate}[label=\textbf{3.},leftmargin=2.6em]
  \item \textbf{要点}：行为指标示例——统计转发该文本的用户
        在其后一段时间内是否仍参与官方相关负面二创；
        文本指标示例——该文本变体的演化方向（强化立场、
        还是滑向纯玩梗）。不一致时\keyw{接受行为指标}：
        文本只能证明“社群愿意如此自我叙述”（自我理解），
        不能证明个体按此行动（动机证据）——这正是第六章
        对“奶蛙不叛”的处理方式。
        \textbf{评分点}：说出两指标的分界理由，
        且裁决与第六章定位一致。
\end{enumerate}

\section{第七章}\label{app:ans:07}

\begin{enumerate}[label=\textbf{1.},leftmargin=2.6em]
  \item \textbf{要点}：四列齐全（压力源—补偿行为—
        可观察指标—竞争性解释）；指标可在不接触隐私的
        前提下计数（公开帖计数、话题标签频次）；
        证伪条件具体到观察层面。
        \textbf{评分点}：自查一步（指标是否隐私安全）
        必须真的执行——若原始指标涉隐私且未改写，本题不及格。
\end{enumerate}

\begin{enumerate}[label=\textbf{2.},leftmargin=2.6em]
  \item \textbf{要点}：判别实验示例——把“圈内身份”与
        “压力水平”正交：高压力圈外组（临近大考的圈外学生）
        与低压力圈内组（假期中的同好会核心成员）。转发频率
        随压力变 → 支持群体稀释；随身份变 → 支持
        身份认证。自然实验版本：考试周前后同一人群的
        频率变化（压力变、身份不变）。
        \textbf{评分点}：两组至少在一个关键变量上正交；
        说清什么结果对应哪个解释。
\end{enumerate}

\begin{enumerate}[label=\textbf{3.},leftmargin=2.6em]
  \item \textbf{要点}：若 M1 证伪而 M3 独立成立：现象
        可持续性\keyw{上升}——情绪压力是周期性驱动
        （考试周），风格规范是长期性的（社群身份），
        现象从“周期性情绪出口”转为“稳定风格规范”，
        寿命预期延长、季节波动减弱。反之（仅 M1 成立）：
        现象随压力节律波动、随毕业潮流失受体，
        寿命与校园节律绑定。
        \textbf{评分点}：两个方向都写；把“分层可独立证伪”
        与“寿命预测”连接起来。
\end{enumerate}

\begin{enumerate}[label=\textbf{4.},leftmargin=2.6em]
  \item \textbf{要点}：策略选项示例。(1) 转移视线：
        正面内容供给挤压丑化空间——改变 M4 的场景供给结构；
        (2) 官方下场玩梗：把 M5 的逆反能量转化为授权流量，
        风险是触发 M3 的反精致反弹（收编指控）；
        (3) 选择性维权：对商用侵权维权、对玩梗沉默——
        把 M5 的“认证信号”与“法律信号”分离；
        (4) 开放共创授权：重设 M5 的权利结构，
        使打压—认证回路失去输入。
        \textbf{评分点}：每个选项对应到具体机制的变动，
        并给出副作用。
\end{enumerate}

\section{第八章}\label{app:ans:08}

\begin{enumerate}[label=\textbf{1.},leftmargin=2.6em]
  \item \textbf{要点}：逐位对齐（两词均 5 音节）：
        嘎 / 刚——声母 \texttt{g} 同（0）、韵母
        \texttt{a}/\texttt{ang}（1）；嘎 / 刚（1）；
        滴 / 点——声母 \texttt{d} 同（0）、韵母
        \texttt{i}/\texttt{ian}（1）；辣 / 了——声母
        \texttt{l} 同（0）、韵母 \texttt{a}/\texttt{e}（1）；
        虾 / 谁——声母 \texttt{x}/\texttt{sh}（1）+
        韵母 \texttt{ia}/\texttt{ui}（1）= 2。
        总代价 6 分，$d = 6/5 = 1.20$。逐位对齐即最优：
        任何增删（每次 1.5）与重排的代价都不低于它。
        仍可成立的理由：等价关系由三条标准\keyw{联合}判定
        ——(2) 场景稳定（该空耳在特定桥段被社群稳定使用）、
        (3) 实例可检索（存在公开二创实例）；$d$ 只服务
        标准 (1) 的\keyw{排序}，而标准 (1) 的文字表述是
        “大体对应”，无数值阈值——音节数一致、叠音结构
        对应（嘎嘎$\leftrightarrow$刚刚）已满足其质性要求。
        \textbf{评分点}：对齐方案 + 计算正确；
        能引用 (2)(3) 与“$d$ 是排序指标”两条警告。
\end{enumerate}

\begin{enumerate}[label=\textbf{2.},leftmargin=2.6em]
  \item \textbf{要点}：最小对照示例。样本：同一平台、
        同一时间窗内“改结尾”与“保留结尾”的奶蛙二创视频
        各 $N$ 条（题材配对）。变量：是否改结尾（二值）、
        传播表现（转发 / 点赞，取对数）。统计量：两组中位数
        比较或配对差值检验。混杂排除：只取作者
        \keyw{首次发布}的作品（人人都不熟练），
        或把作者历史发布量作为配对 / 协变量。
        \textbf{评分点}：样本—变量—统计量三件套齐全；
        混杂处理具体到操作。
\end{enumerate}

\begin{enumerate}[label=\textbf{3.},leftmargin=2.6em]
  \item \textbf{要点}：机制设计示例——对单账号单位时间的
        批量低原创发布设“梯度曝光衰减”（类累进排放税：
        发布量越大，单条初始曝光越低），费率公开。
        反弹分析：社群会把限流本身再梗化
        （“被吞了”成为新素材），触发 M5 逆反认证——
        \keyw{治理动作成为新的认证信号}；同时限流损伤 M2
        的群体稀释通道，压力表达将寻找替代出口。
        \textbf{评分点}：有“费率”结构；反弹对应到
        具体机制并说明传播后果。
\end{enumerate}

\section{第九章}\label{app:ans:09}

\begin{enumerate}[label=\textbf{1.},leftmargin=2.6em]
  \item \textbf{要点}：排序示例：结尾集体大笑 $>$ 奶蛋 $>$
        品牌联名奶茶杯。依据三维：内容可换——大笑不依赖
        具体桥段，奶蛋仍绑定“奶伴”对战叙事，联名杯绑定
        品牌与产品；载体可换——大笑跨平台流通，奶蛋限于
        B 站语境，联名杯限于线下商品；署名可弃——三者
        均已匿名化，此项不构成区分。
        \textbf{评分点}：三维分别论证，不笼统排序。
\end{enumerate}

\begin{enumerate}[label=\textbf{2.},leftmargin=2.6em]
  \item \textbf{要点}：判据示例。官方视角：\keyw{可执行性}
        ——对象超出可执法清单（匿名、跨平台、线下化）
        即失控；社群视角：\keyw{伤害}——出现实质损害
        （诈骗、人身攻击）才失控。差异根源：官方以权利
        结构为界，社群以内容性质为界。
        \textbf{评分点}：判据可操作（能裁决具体案例）；
        能说明两视角的界线性质不同。
\end{enumerate}

\begin{enumerate}[label=\textbf{3.},leftmargin=2.6em]
  \item \textbf{要点}：观点一隐含前提：匿名 = 放弃权利主张
        （把“不可识别”当作“权利放弃”）。观点二隐含前提：
        社群劳动创造的价值应归属劳动者（劳动价值预设 +
        可识别的贡献分配）。E3 支持的\keyw{分析姿态}：
        叙述二创形象时附署名权归属，使创作者作为权利主体
        在文本中可见；E3 不支持的\keyw{结论}：两种观点都是
        规范性结论，E3 作为考据伦理规则不裁决分配正义
        ——它只防止权利主体在叙述中消失。
        \textbf{评分点}：分清“姿态”与“结论”两个层次。
\end{enumerate}

\section{第十章}\label{app:ans:10}

\begin{enumerate}[label=\textbf{1.},leftmargin=2.6em]
  \item \textbf{要点}：低—中强度分支：例行移除（校方 0）
        + 纪念文本（学生 1）+ 线上存续（持久物 1）。
        高强度分支：公开通报或追责（校方 2）$\to$
        组织化表态（学生 2）$\to$ 制度化对抗或申请合规化
        （持久物 2）。转折点候选：校方是否\keyw{公开追责}
        （把例行处置升格为立场表态）、学生是否出现
        \keyw{组织化}动作（从个体玩梗变为集体对抗）。
        \textbf{评分点}：两个分支的三维读数齐全；
        转折点指认到具体动作。
\end{enumerate}

\begin{enumerate}[label=\textbf{2.},leftmargin=2.6em]
  \item \textbf{要点}：观察点——流动投放点与固定立牌
        各一（可比）；时段——投放后固定窗口
        （如午间 60 分钟 $\times$ 5 天）；计数规则——
        合影率 = 驻足者中拍照人数 / 驻足总人数，
        “驻足”定义为停留不短于 5 秒；混杂——天气、
        考试周、人流构成（逐场记录作控制变量）；
        伦理——公共场所匿名计数，不摄录人脸，
        不跟踪个体。
        \textbf{评分点}：计数规则事先写定；
        至少两类混杂及处理。
\end{enumerate}

\begin{enumerate}[label=\textbf{3.},leftmargin=2.6em]
  \item \textbf{要点}：找到的文本须有可撤回装置
        （戏谑语体、夸张到不可能当真、匿名性）；
        若不成立，须指出其替代机制（实名表态、
        组织化申诉、直接协商）。加分项：把该文本与
        赞诗作结构对比（同一功能、不同装置）。
        \textbf{评分点}：“可撤回装置”被具体指认，
        而非笼统说“也是玩梗”。
\end{enumerate}

\section{第十一章}\label{app:ans:11}

\begin{enumerate}[label=\textbf{1.},leftmargin=2.6em]
  \item \textbf{要点}：前瞻替代——固定样本日记法：
        招募小样本（如 60 人），考试周前后连续 21 天
        每日推送一题（“今日是否转发奶蛙类内容 0/1 +
        情境勾选”）。取舍：日记法单样本成本高、
        流失率高，但消除近因效应与情绪一致性回忆；
        事件—时间对齐精度远高于回填。
        \textbf{评分点}：说出回填法的两类系统偏差
        及其方向；成本与效度的取舍双向呈现。
\end{enumerate}

\begin{enumerate}[label=\textbf{2.},leftmargin=2.6em]
  \item \textbf{要点}：反方最强理由：辨识“故意劣化”
        依赖整体观感（格式塔），拆成标记会系统性漏掉
        明显故意做差但无标记的内容——\keyw{效度}受损。
        正方最强理由：无标记的整体观感编码无法区分
        故意与真的做差，编码者一致性必然崩溃——
        \keyw{信度}崩塌。裁决示例：信度优先
        （宁可漏、不可滥），折中方案——标记编码为主、
        入假设检验；另设“整体判断”辅助列，不入检验，
        只用于生成后续假设。
        \textbf{评分点}：说清信度—效度冲突的结构；
        裁决有明确取舍而非和稀泥。
\end{enumerate}

\begin{enumerate}[label=\textbf{3.},leftmargin=2.6em]
  \item \textbf{要点}：修改示例——措辞去“屏蔽”改中性
        （“如果你刷奶蛙的频率变低了，你的使用会怎么变？”）；
        位置后移至访谈末尾；附正常化选项（“也许不会变
        ——很多人过段时间就换新梗了”）以降诱导。
        若修改后失去分析价值的理由：M5 检验需要
        “压制—反应”的强刺激设定，中性化的“频率变低”
        测到的只是兴趣衰退（M4 范围），无法生成
        压制反弹假设——此时应承认该题不能温和化，
        改以第十章已发生的移除事件作自然实验替代。
        \textbf{评分点}：修改真正降低诱导（三处之一即可）；
        “失去价值”的论证触及 M5 与 M4 的区分。
\end{enumerate}

\section{第十二章}\label{app:ans:12}

\begin{enumerate}[label=\textbf{1.},leftmargin=2.6em]
  \item \textbf{要点}：条目号正确；原句真实存在
        （注明章节位置）；修改句符合该条目的操作要点；
        依据引用条目原文。批判框所列三处病灶之外的自选
        例子同样有效——\keyw{发现本书未登记的新违例，
        应评为高质量答案}。
        \textbf{评分点}：四要素齐全；修改不引入新违例
        （如把 C 级句改为 A 级口吻）。
\end{enumerate}

\begin{enumerate}[label=\textbf{2.},leftmargin=2.6em]
  \item \textbf{要点}：前三步示例。(1) 更新事实手册
        （伦理检查点：重新访问 C 级来源时仍守 E1--E3）；
        (2) 获取校本材料授权（伦理检查点：当事人同意 +
        匿名化选项）；(3) 按第十章 C 级纪律写案例
        （伦理检查点：单一来源不作论证前提）。
        最易省略的是第 (2) 步——“发在群里好像就是
        公开的”。后果：校本材料成为无授权的个人信息，
        违反最小化收集；续编版本同时破坏 CC BY-NC-SA
        的署名链条。
        \textbf{评分点}：每步有检查点；对“最易省略”
        的论证触及“公开可见 $\ne$ 可自由使用”。
\end{enumerate}

\begin{enumerate}[label=\textbf{3.},leftmargin=2.6em]
  \item \textbf{要点}：反例候选——鲁迅杂文（幽默进入
        论证且增强说服力）、讽喻传统（《动物农场》以
        幽默承担政治论证）、以讽刺组织论据的学术写作。
        成立前提：(1) 幽默与论证指向同一对象；
        (2) 受众能解码幽默；(3) \keyw{幽默组织证据，
        而不替代证据}。奶蛙研究是否具备：部分具备——
        社群文本（赞诗）本身就是幽默论证，可作研究对象；
        但研究者的幽默（教材体例）与论证对象（社群行为）
        分属不同层次，若研究者的幽默直接充当论证
        （以段子代替机制分析），前提 (3) 即被破坏——
        这正是第十二章把幽默限制在体例层的理由。
        \textbf{评分点}：说出前提 (3)；能把反例前提
        与本书的分层处理对接。
\end{enumerate}

\chapter{资料来源与版权说明}\label{app:sources}

本章（附录）汇总本书的资料来源构成，并说明本书在版权上的位置。
正文引用以 \ref{app:timeline} 页大事年表与参考文献列表为准，
此处只作\keyw{来源类型学}与\keyw{权利边界}的说明，不复述具体事实。

\section{来源类型与可信度}\label{app:sources:types}

本书的全部材料按产生方式分为四类，与正文的 A/B/C/M 四级对应：

\begin{table}[htbp]\centering
  \caption{材料来源类型与可信度对应}\label{tab:04:source}
  \small
  \begin{tabularx}{\textwidth}{@{}llX@{}}
    \toprule
    类型 & 可信度 & 说明 \\ \midrule
    一手权威报道 & \conf{A} & 如《人民日报》关于奶龙 IP 创新链的报道；
      可作为事实陈述直接引用 \\
    机构/平台公开信息 & \conf{B} & 如中国日报网转载的六周年活动通稿、
      百科与平台条目；可作事实陈述，注明来源与日期，
      多来源冲突时并列 \\
    社群考据与二手评述 & \conf{C} & 如失传媒体中文维基词条、
      自媒体考据文章；只能作为“线索”与“说法”，
      不得作为论证前提，正文须以“据……”“有说法认为”引出 \\
    本书自建模型 & \conf{M} & 本书为教学目的构造的框架、方程、
      示例数据与虚构问卷；必须显式声明非经验事实，
      参数取值为示意值 \\ \bottomrule
  \end{tabularx}
\end{table}

\begin{nlmethod}[使用纪律]
  \textbf{C 级材料的三条纪律}：
  (1) 不作为唯一依据支撑任何 A/B 级口吻的结论；
  (2) 涉及具体个人时，不转述其健康、生死、私德等个人境况的推测；
  (3) 当 C 级考据自身声明“来源缺失/仅为推测”时，
  本书照实保留这一自我限定，不替考据者“补全证据”。
\end{nlmethod}

\section{奶蛙起源考据的来源特别说明}\label{app:sources:naiwa}

奶蛙起源线索（第三章、第四章）主要依赖社群考据（\conf{C}）。
本书在引用时保留考据材料自身的三处自我限定：

\begin{enumerate}[label=\arabic*.]
  \item 考据指向的智能体（ID \fname{zNCvQDKfXy}）现已无法搜索、
        分享，评论区关闭，\keyw{不可复核}；
  \item 失传媒体维基词条自陈“待改进”“来源缺失”，
        且“星野 App 出现 bug”的结论“仅为推测、缺乏直接证据”；
  \item 多个生成路径推测均被考据者证伪，复现亦未成功。
\end{enumerate}

因此本书将奶蛙起源定性为一条\keyw{“未结案的证据链”}，
作为教学示范材料，而非已证事实。这一处理是全书
证据分级制度最严格的一次执行，读者可将其视为
“C 级材料应如何被使用”的标准样例。

\section{版权与合理使用}\label{app:sources:copyright}

\subsection{本书对形象的使用}

本书\keyw{不嵌入}受版权保护的原始动画截图、表情包原图、
周边商品图片。全部插图均为本书以 TikZ 自制的
抽象化示意图，不复制任何受保护的具体画面。
涉及“奶龙”动画形象、剧集名称及其相关商标时，
本书将其知识产权归属第七印象文化传媒（深圳）有限公司；
涉及“奶蛙”二次创作形象时，其署名权归属原始创作者。

\begin{nlcritique}[权利边界]
  本书对上述形象的使用限于\textbf{学术评论、教学示例与文化研究}
  目的，不主张任何权利，亦不构成对原作品著作权的替代或商业利用。
  正文将奶蛙二创的“未经授权形象再生产”作为\keyw{分析对象}
  加以讨论，不构成对任何具体侵权行为的背书或鼓励。
  如权利人认为本书某一处使用不当，请联系编写组予以替换或删除。
\end{nlcritique}

\subsection{引用社区用语与剧集名称}

本书引用的社群用语（如“糖”“倒放大法”“鲱鱼罐头”等）
与剧集名称（如《奶龙搞怪大作战》《奶龙之奇妙冒险》），
属于对公开传播内容的\keyw{指称性使用}，
用于识别与讨论特定文化对象，不构成对表达本身的复制。

\subsection{本书的再传播许可}

本书遵循\keyw{署名—非商业性使用—相同方式共享
（CC BY-NC-SA 4.0）}协议发布。各高校同好会以本校
活动材料为基础进行修订与续编时，须满足：
(1) 署名本书编写组与原出处；(2) 不用于商业目的；
(3) 衍生版本采用相同协议。修订时请同步更新
“事实手册”中相应条目的可信度标注与访问日期。

\section{访问日期与时效}

网络来源均记录首次访问日期（见参考文献列表）。
鉴于网络材料的时效性（“引文腐烂”问题，第十二章讨论），
读者在引用 \conf{B} / \conf{C} 级来源时，
应以本书记录的访问日期为基线，并在再引用时
重新核验链接的当前可达性与内容的现时状态。

\noindent\srcnote{本附录为版权与来源的“元说明”，
不替代正文的事实引用；具体引文格式以参考文献列表
（GB/T 7714—2015）为准。}


%% ===================== 参考文献 =====================
\cleardoublepage
\addcontentsline{toc}{chapter}{参考文献}
\bibliography{refs/references}

%% ===================== 索引 =====================
\cleardoublepage
\printindex

%% ===================== 后记 =====================
\chapter*{第一版后记}
\addcontentsline{toc}{chapter}{第一版后记}
\markboth{第一版后记}{第一版后记}

本书到第十二章结束。

从一份二十页的社团讲义到一部十二章的教材，中间隔着的不是篇幅，
而是三件事：把“章节”变成“章—节—目”，把“说法”从“事实”里
分出来，把“感想”变成可检验的假设。这三件事里没有一件
是奶蛙本身的“故事”，它们都是\keyw{关于如何研究奶蛙}的
元工作。如果本书对读者最有用的部分是这些，那它就没白写。

有几句话想在最后再说一遍。

第一，\textbf{本书没有任何一项来自作者亲自采集的经验数据}。
第十一章的问卷、示例数据表，都是供修订者与后续研究者使用的
\keyw{设计}与\keyw{教学样例}，不是本书的测量结果。
谁若把本书当作“已有结论”来引用，请回到第七章的假设
与第十二章的可复现性清单，先看它的证伪条件。

第二，\textbf{本书对具体个人保持沉默}。
奶蛙起源考据指向的创作者，以及一切出现在社群材料里的
真实个体，本书只作指称，不作推断——尤其不转述任何关于
其健康、生死、私德的内容。这是证据分级制度里
最难执行、也最不能妥协的一条。

第三，\textbf{本书把“乐子”当作研究对象，但不消费它}。
我们承认奶蛙二创中存在未经授权的形象再生产，
这在版权法上有争议；本书讨论它，是为了理解
“正典—变种”结构如何运作，而不是为了鼓励或背书
任何一种具体的再生产行为。

欢迎各高校同好会以本校的奶蛙活动为基础，
对本书进行修订与续编——更新事实手册中的条目与访问日期，
补充本校的案例与数据，重做第十一章的问卷。
修订时请遵守 CC BY-NC-SA 4.0 协议的三项义务，
并在新的“事实手册”里保留本书的证据分级制度。

如果读完这本书，你能同时做到两件事——
\keyw{看懂一只青蛙为什么让人捧腹大笑}，
并且\keyw{说不清哪一句是事实、哪一句是说法}——
那么它可能没有完全失败。

\vspace{2em}
\begin{flushright}
\BookOrg\\
\BookDate
\end{flushright}


%% ===================== 封底 =====================
%% 封底：与封面共享边距、配色、字体与分隔线。
\clearpage
\thispagestyle{empty}
\pdfbookmark[0]{封底}{backcover}
\begin{tikzpicture}[remember picture,overlay]
  \fill[NLCoverGreen] (current page.north west) rectangle (current page.south east);
  \begin{scope}[shift={(current page.north west)},x=1cm,y=-1cm]
    \fill[NLCoverGold] (0,0) rectangle (.38,29.7);
    \draw[NLCoverPaper!35!NLCoverGreen,line width=.6pt] (2.1,2.66) -- (18.85,2.66);
    \node[nlcovertext,text=NLCoverPaper,font=\NLCoverSans\fontsize{11.5}{15}\selectfont]
      at (2.1,1.72) {迷因文化研究系列};
    \node[nlcoverright,text=NLCoverPaper!75!NLCoverGreen,font=\NLCoverLatin\fontsize{8.4}{12}\selectfont]
      at (18.85,1.85) {{\addfontfeatures{LetterSpace=9}NAILONG STUDIES}};

    \node[nlcovertext,text=NLCoverPaper,font=\NLCoverBold\bfseries\fontsize{26}{40}\selectfont]
      at (2.1,4.06) {对现象保持幽默，\\对方法保持严肃。};
    \draw[NLCoverGold,line width=2.2pt] (2.1,7.52) -- (5.35,7.52);

    \node[nlcovertext,text=NLCoverGold,font=\NLCoverSans\fontsize{12.5}{18}\selectfont]
      at (2.1,8.36) {一部关于数字迷因与校园文化的通识读本};
    \node[nlcovertext,text=NLCoverPaper!88!NLCoverGreen,text width=16.45cm,
      font=\NLCoverSans\fontsize{11.5}{21}\selectfont]
      at (2.1,9.65) {从奶龙到奶蛙，从一段笑声到一场校园相遇，本书以十二章的篇幅，
      连接形象演变、迷因理论、传播机制与社群实践。通过概念界定、史料辨析、
      案例讨论和研究方法训练，引导读者把熟悉的网络现象转化为可以分析的文化问题。};

    \foreach \y in {13.35,16.55,19.75,22.95} {
      \draw[NLCoverPaper!22!NLCoverGreen,line width=.45pt] (2.1,\y) -- (18.85,\y);
    }
    \node[nlcovertext,text=NLCoverGold,font=\NLCoverLatin\fontsize{21}{27}\selectfont]
      at (2.1,14.12) {01};
    \node[nlcovertext,text=NLCoverPaper,font=\NLCoverBold\bfseries\fontsize{14.5}{19}\selectfont]
      at (4.15,14.10) {形象与历史};
    \node[nlcovertext,text=NLCoverPaper!76!NLCoverGreen,text width=14.4cm,
      font=\NLCoverSans\fontsize{10.2}{17}\selectfont]
      at (4.15,15.07) {厘清奶龙与奶蛙的概念边界，回到传播史与材料本身。};

    \node[nlcovertext,text=NLCoverGold,font=\NLCoverLatin\fontsize{21}{27}\selectfont]
      at (2.1,17.32) {02};
    \node[nlcovertext,text=NLCoverPaper,font=\NLCoverBold\bfseries\fontsize{14.5}{19}\selectfont]
      at (4.15,17.30) {迷因与社群};
    \node[nlcovertext,text=NLCoverPaper!76!NLCoverGreen,text width=14.4cm,
      font=\NLCoverSans\fontsize{10.2}{17}\selectfont]
      at (4.15,18.27) {结合符号学、情绪传播与校园案例，理解网络文化的生成。};

    \node[nlcovertext,text=NLCoverGold,font=\NLCoverLatin\fontsize{21}{27}\selectfont]
      at (2.1,20.52) {03};
    \node[nlcovertext,text=NLCoverPaper,font=\NLCoverBold\bfseries\fontsize{14.5}{19}\selectfont]
      at (4.15,20.50) {方法与实践};
    \node[nlcovertext,text=NLCoverPaper!76!NLCoverGreen,text width=14.4cm,
      font=\NLCoverSans\fontsize{10.2}{17}\selectfont]
      at (4.15,21.47) {从内容分析到访谈与考据，训练证据意识与研究伦理。};

    \node[nlcovertext,text=NLCoverPaper!72!NLCoverGreen,font=\NLCoverSans\fontsize{9.4}{14}\selectfont]
      at (2.1,23.69) {概念\enspace ·\enspace 史料\enspace ·\enspace 案例\enspace ·\enspace 习题\enspace ·\enspace 方法};

    \fill[NLCoverPaper] (.38,25.38) rectangle (21,29.7);
    \node[nlcovertext,text=NLCoverMuted,font=\NLCoverSans\fontsize{9.2}{13}\selectfont]
      at (2.1,26.12) {适用对象};
    \node[nlcovertext,text=NLCoverGreen,font=\NLCoverSans\fontsize{10.5}{16}\selectfont]
      at (2.1,26.89) {高校通识课程\enspace ·\enspace 学生社团\enspace ·\enspace 迷因文化爱好者};
    \node[nlcovertext,text=NLCoverMuted,font=\NLCoverSans\fontsize{8.5}{12}\selectfont]
      at (2.1,28.24) {\BookEdition\enspace ·\enspace \BookDate};
    \node[nlcoverright,text=NLCoverMuted,font=\NLCoverSans\fontsize{8.5}{12}\selectfont]
      at (18.85,28.24) {戏拟教材\enspace /\enspace 非营利性讲义};
  \end{scope}
\end{tikzpicture}
\null\clearpage


\end{document}
